\chapter{Calabi--Yau 流形的度量问题}\label{chap:7}

\SourceParagraph{ch07-intro-001}{519}
在很大程度上，一直忽略了由度量引起的问题。\BookChapterRef{2}讨论的是拓扑场论，而非共形场论。\BookChapterRef{4}与\BookChapterRef{5}主要讨论 D 膜的问题，这使得转而考虑簇上的凝聚层导出范畴；在开始考虑稳定性条件以前，这与度量无关。在镜像的一侧，考虑的是 Fukaya 范畴，对它来说，重要的是辛形式，而非度量。讨论 Strominger--Yau--Zaslow 猜想时，度量开始起作用，因为没有度量，就无法定义特殊 Lagrange 子流形。不过，在\BookChapterRef{6}中，为了给出 SYZ 的一些拓扑版本，很快就放弃了度量；涉及特殊 Lagrange 环面的完整 SYZ 版本，过于困难。

\SourceParagraph{ch07-intro-002}{519}
本章希望更认真地对待度量，以便更深入地理解 SYZ 猜想。将讨论的许多结果都相当困难，因此，本章具有综述的性质，主要解释结果，而不给出证明。

\SourceParagraph{ch07-intro-003}{519}
从 Ricci 平坦度量的例子开始。大部分例子，是复流形上的度量；此时，特别容易陈述 Ricci 平坦的意义。若 $X$ 为 Kähler 流形，Kähler 形式为 $\omega$，$\Omega$ 为 $X$ 上处处非零的全纯 $n$ 形式，则对于某个函数 $f\colon X\to\mathbb R$，有
\[
\omega^n=e^f\,(-1)^{n(n-1)/2}\frac{\mathrm i^n n!}{2^n}\,
\Omega\wedge\overline\Omega,
\]
Ricci 曲率形式 $\rho$ 为
\[
\rho=-\mathrm i\,\partial\bar\partial f.
\]
当 $\rho=0$ 时，度量 Ricci 平坦。在 $X$ 紧的情形中，这等价于 $f$ 在每个连通分支上为常数。即使不存在整体的全纯 $n$ 形式，也可以用这种方式局部地定义 Ricci 曲率形式。

\SourceParagraph{ch07-intro-004}{519}
在上述紧 Kähler、第一 Chern 类为零的情形中，由 Yau\BookCite{ref475}的结果，知道 Ricci 平坦度量存在，不过，对这些度量仍了解很少。然而，在具有大量对称性的非紧例子中，可以明确写出 Ricci 平坦度量。将综述文献中出现的一些例子。还将较简略地讨论一些非 Kähler 例子，特别是某些完整群为 $G_2$ 的例子，它们在物理文献中起着越来越重要的作用。

\SourceParagraph{ch07-intro-005}{520}
完成这些例子的讨论以后，将回顾 Dominic Joyce 关于构造特殊 Lagrange 子流形例子的工作。它给出一些令人惊讶的特殊 Lagrange 纤维化例子，其行为与预期根本不同。特别地，这些纤维化的判别集，余维为一，而\BookChapterRef{6}所研究的拓扑环面纤维化，则具有良好的余维二判别集。这项工作尤其令人惊讶的一个结果是，对偶化时，判别集很可能发生改变。因此，\BookChapterRef{6}所设想的强形式 SYZ 对偶，即对非奇异环面纤维作对偶，并按该章解释的方式改变奇异纤维，不能成立。这要求重新评估并重新表述 SYZ 猜想。本质上，Strominger、Yau 与 Zaslow 的物理论证，应当只在大复结构极限处成立；离开这一极限时，可以预期，猜想会出现一些不精确之处。通过研究大复结构极限点附近 Ricci 平坦度量的行为，可以在数学上将这一点说得更精确。因此，最后将综述 Gross 与 Wilson\BookCite{ref210}关于 K3 曲面的大复结构极限的工作。这最终导向对 SYZ 猜想的一种重新表述，或许已在现有技术能够企及的范围内。

\section{Ricci 平坦度量的例子与各种拟设}\label{sec:7.1}

\subsection{一些基本例子}\label{sec:7.1.1}

\SourceParagraph{ch07-7.1.1-001}{520}
\begin{theorem*}[Yau 定理]
若 $X$ 为紧 Kähler 流形，第一 Chern 类在实上同调中为零，则 Yau 定理\BookCite{ref475}保证：对于任意选择的 Kähler 类，存在唯一的 Ricci 平坦 Kähler 度量，其关联的 $(1,1)$ 形式 $\omega$ 属于给定的类。
\end{theorem*}
常常通过把度量称为 $\omega$，混用度量与其关联的 $(1,1)$ 形式。问题化为在 $X$ 上求解复 Monge--Ampère 方程。给定 $X$ 上的 Kähler 度量 $\omega$，其 Ricci 形式为 $\mathrm i\,\partial\bar\partial f$；其中，函数 $f$ 归一化为满足 $\int_Xe^f\omega^n=\int_X\omega^n$。需要找到函数 $\phi$，可以用 $\int_X\phi\omega^n=0$ 作归一化，使得 $\omega+\mathrm i\,\partial\bar\partial\phi$ 为某个度量的 Kähler 形式，并满足
\[
(\omega+\mathrm i\,\partial\bar\partial\phi)^n=e^f\omega^n.
\]
那么，度量 $\omega+\mathrm i\,\partial\bar\partial\phi$ Ricci 平坦。

\begin{proof}[证明思路]
\SourceParagraph{ch07-7.1.1-002}{520}
这个度量的唯一性，源自 Calabi。例如，可以通过一个简单的分部积分论证证明。Yau 用连续性方法证明存在性：对于 $t\in[0,1]$，求函数 $\phi_t$，归一化为满足 $\int_X\phi_t\omega^n=0$，使得
\begin{enumerate}[label=(\arabic*)]
\item $\omega+\mathrm i\,\partial\bar\partial\phi_t$ 是某个度量的 Kähler 形式；
\item
\[
(\omega+\mathrm i\,\partial\bar\partial\phi_t)^n
=e^{tf}\frac{\int_X\omega^n}{\int_Xe^{tf}\omega^n}\,\omega^n.
\]
\end{enumerate}

\SourceParagraph{ch07-7.1.1-003}{521}
因此，$t=0$ 时有解 $\phi_0=0$，而 $t=1$ 时，就是希望求解的方程。令
\[
A=\{t\in[0,1]\mid\text{存在 }C^{k,\alpha}\text{ 解 }\phi_t\},
\]
其中，$k\geq2$、$0<\alpha<1$。这里，$C^{k,\alpha}$ 表示第 $k$ 阶导数具有指数 $\alpha$ 的 Hölder 连续性的函数。实值函数 $f$ 具有指数 $\alpha$ 的 Hölder 连续性，是指存在常数 $C$，使得对于定义域内所有 $x,y$，有 $|f(x)-f(y)|\leq C\|x-y\|^\alpha$。$\alpha=1$ 时，就是 Lipschitz 连续性。这是求解偏微分方程时，典型的正则性假设。Yau 的证明，是说明 $A$ 既开又闭，因此 $A=[0,1]$。开性来自标准分析技术，例如 Hodge 理论，以及 Banach 空间上的逆函数定理；闭性则来自对上述方程族的解所作的高度非平凡、独立于 $t$ 的先验估计。
\end{proof}

\SourceParagraph{ch07-7.1.1-004}{521}
虽然由此知道这些度量存在，不过，在紧情形中，仍没有已知的非平坦显式例子。由于所关心的通常例子缺乏连续的等距对称，实际找出这样的度量将非常困难。本章后面将看到，在适当情形中，仍可以希望通过粘合技术，找到这些度量的良好近似。这些近似，依赖于知道某些非紧流形上 Ricci 平坦 Kähler 度量的显式形式。对于非紧流形，度量不再唯一，不过，可能施加一定的对称性，从而把上述高度复杂的非线性偏微分方程，化为可以相当明确地求解的常微分方程。

\SourceParagraph{ch07-7.1.1-005}{521}
虽然下面不会用到这一点，还指出，数值近似 Ricci 平坦度量的研究已有进展。利用 Donaldson\BookCite{ref130}的思想，Douglas、Karp、Lukic 与 Reinbacher\BookCite{ref140}，已经在超曲面上计算出 Ricci 平坦度量的良好数值近似。

\SourceParagraph{ch07-7.1.1-006}{521}
\BookStatementNumber{1}
\begin{example}[Eguchi--Hanson 度量]\label{ex:7.1}
令 $X=\mathbb C^2/\{\pm1\}$，即 $\mathbb C^3$ 中的二次锥，消解为 $T^*\mathbb{CP}^1$，与 $S^2$ 的实切丛微分同胚。在 $\mathbb C^2$ 上，令 $u=|z_1|^2+|z_2|^2$，寻找势函数形如 $f(u)$ 的 Ricci 平坦 Kähler 度量。求解后，得到显式解，见文献\BookCite{ref307}第293页：
\[
f_a(u)=u\left(1+\frac{a^2}{u^2}\right)^{1/2}
+a\log\frac{u}{\sqrt{u^2+a^2}+a},\qquad a>0.
\]
可以证明，Kähler 形式 $\omega_a=\mathrm i\,\partial\bar\partial f_a$ 延伸过 $T^*\mathbb{CP}^1$ 的零截面 $E$，给出 Eguchi--Hanson 度量。这是 $\operatorname{SO}(3)$ 不变的 Ricci 平坦 ALE Kähler 度量；ALE 指渐近局部 Euclid 型。以 $[E]$ 表示 $E$ 的整 Poincaré 对偶类，其 Kähler 类为 $-2a[E]$。二维球面 $E$ 的曲率为 $a^{-1}$，面积为 $4\pi a$。在 $E$ 的补集中的任意紧集上，当 $a\to0$ 时，Eguchi--Hanson 度量迅速收敛到轨形平坦度量\BookCite{ref307}。

\SourceParagraph{ch07-7.1.1-007}{522}
这个相对简单的例子，为超 Kähler 的，并且十分重要，因为它在复二维的简单节点的光滑化与消解上，都定义 Ricci 平坦 Kähler 度量。例如，若紧 K3 曲面 $X$ 带有具有节点的轨形度量，则取 $a$ 很小，并粘入 Eguchi--Hanson 度量的适当副本，就可以在消解或光滑化上，得到 Ricci 平坦度量的良好近似。更多细节，见文献\BookCite{ref307}。

\SourceParagraph{ch07-7.1.1-008}{522}
这也是余齐性为一的度量的一个例子。这意味着，有群 $G$ 通过等距变换作用于空间，而群的一般轨道的余维为一。在上面的例子中，$G=\operatorname{SO}(3)$。将会看到，通常还有一个维数更低的轨道，形成\BookRef{fig:7.1}{图1}那样的图景。对于 Eguchi--Hanson 度量，一般轨道为 $S^3/\{\pm1\}$，特殊轨道为 $S^2$，即余切丛的零截面。
\begin{figure}[H]
\centering
\BookImage{01_Numbered_Figures/Chapter_07/Ch07_Fig01_PDFp522_Bookp512.pdf}
\BookNumberOnlyCaption\label{fig:7.1}
\end{figure}

\SourceParagraph{ch07-7.1.1-009}{522}
可以用多种方式推广 Eguchi--Hanson 度量。在复二维中，利用超 Kähler 商构造\BookCite{ref238}，Kronheimer\BookCite{ref316}构造出 ALE 引力瞬子，即任意二维有理双点的消解上的超 Kähler 度量。对于 $A_k$ 奇点，这些度量已经通过 Gibbons--Hawking 构造，以多瞬子度量的形式得到。还有向更高维的自然推广，即广义 Eguchi--Hanson 度量\BookCite{ref079,ref177,ref274}。它们是 $\mathbb{CP}^{k-1}$ 上 Hopf 丛的第 $k$ 次幂，即 $\mathcal O(-k)$，上的完备度量，完整群为 $\operatorname{SU}(k)$；等价地，它们位于商奇点 $\mathbb C^k/C_k$ 的保典消解上，其中，$C_k=\langle\zeta\rangle$ 为按对角方式作用的 $k$ 次单位根群。度量可以通过只依赖于 $\sum_i|z_i|^2$ 的 Kähler 势表达\BookCite{ref274}。对于 $k=3$，文献\BookCite{ref340}已经在另一背景下写出这个度量。这些度量为 ALE 度量；在上述更高维情形中，它们是已知的完整群为 $\operatorname{SU}(k)$ 的 ALE 度量中仅有的显式例子。对于以自然方式作用于丛的 $\operatorname{U}(k)$，它们的余齐性为一；余维一的轨道为
\[
\operatorname{U}(k)/(\operatorname{U}(k-1)\times C_k)=S^{2k-1}/C_k,
\]
还有一个奇异轨道，即余维更高的轨道，$\mathbb{CP}^{k-1}$。下面将进一步讨论余齐性为一的度量。实际上，这些度量是以下度量的特殊情形：在任意正曲率 Kähler--Einstein 流形的典范丛上，定义完整群包含于 $\operatorname{SU}(k)$ 的度量，见文献\BookCite{ref079}与文献\BookCite{ref408}的定理8.1。
\end{example}

\SourceParagraph{ch07-7.1.1-010}{523}
与上述 $\operatorname{SU}(k)$ 完整群的度量类似，还有 Calabi 给出的著名超 Kähler 度量。

\SourceParagraph{ch07-7.1.1-011}{523}
\BookStatementNumber{2}
\begin{example}[Calabi 超 Kähler 度量]\label{ex:7.2}
这些是复射影空间的全纯余切丛 $T^*\mathbb{CP}^k$ 上的完备超 Kähler 度量，余齐性也为一，等距群 $\operatorname{U}(k+1)$ 作用于形如
\[
\operatorname{U}(k+1)/(\operatorname{U}(k-1)\times\operatorname{U}(1))
\]
的余维一轨道，并有一个奇异轨道 $\mathbb{CP}^k$，见文献\BookCite{ref079}。在 $\mathbb{CP}^k$ 上赋予 Fubini--Study 度量的一个倍数，其常全纯截面曲率为 $2\kappa$。它在全纯余切丛 $T^*\mathbb{CP}^k$ 上诱导 Hermitian 度量。局部地，在仿射开集 $U\subset\mathbb{CP}^k$ 上，Calabi 写出 $T^*\mathbb P^k|_U$ 上的显式 Kähler 势，并证明，它在全纯余切空间上给出良定义的整体超 Kähler 度量。这个度量限制到 $\mathbb P^k$ 上，就是给定的度量；超 Kähler 的2形式三元组，由度量的 Kähler 形式，与全纯余切丛上自然复辛形式的实部、虚部组成。事实上，若加上条件，要求它在 $T^*\mathbb{CP}^k$ 的纤维上的自然 $\operatorname{U}(1)$ 作用下不变，则这些性质唯一确定这个超 Kähler 度量，见文献\BookCite{ref146,ref455}。$k=1$ 时，采用\BookRef{ex:7.1}{例题7.1}的归一化，Calabi 度量的 Kähler 势，就是参数为 $a$ 的 Eguchi--Hanson 势乘以 $a/2$，其中，$a>0$ 由 $\kappa=a^{-2}$ 给出。因此，度量在 $\mathbb P^1$ 上的限制，具有所需的曲率 $2a^{-2}=2\kappa$。对于任意 $k$，这些度量的余齐性为一，等距群 $\operatorname{U}(k+1)$ 作用于形如 $\operatorname{U}(k+1)/(\operatorname{U}(k-1)\times\operatorname{U}(1))$ 的余维一轨道，并有一个奇异轨道 $\mathbb{CP}^k$。
\end{example}

\SourceParagraph{ch07-7.1.1-012}{523}
不过，还可以考虑更高维中的节点，由此导向 Stenzel 度量。

\SourceParagraph{ch07-7.1.1-013}{523}
\BookStatementNumber{3}
\begin{example}[Stenzel 度量及相关度量]\label{ex:7.3}
Eguchi--Hanson 度量向更高维的另一项相关推广，如下给出。对于 $\mathbb C^{n+1}$ 中由方程 $z_0^2+z_1^2+\cdots+z_n^2=0$ 定义的二次锥，可以令 $\rho=|z_0|^2+|z_1|^2+\cdots+|z_n|^2$，寻找形如 $f(\rho)$ 的 Ricci 平坦 Kähler 势。要求 Ricci 形式为零，得到 $f(\rho)=\rho^{(n-1)/n}$。$n=2$ 时，就是轨形平坦度量。$n=3$ 时，所得 Ricci 平坦 Kähler 度量，可以重新理解为 Einstein 流形 $T^{1,1}$ 上的锥形度量；后一个流形在拓扑上为 $S^3\times S^2$，不过，度量不是乘积 Einstein 度量\BookCite{ref084}。同样，可以在光滑化，即方程
\[
z_0^2+z_1^2+\cdots+z_n^2=\varepsilon
\]
所定义的非奇异二次曲面 $Q$ 上，求解形如 $f(\rho)$ 的 Ricci 平坦 Kähler 势。写成 $z_j=x_j+\mathrm i y_j$，并取 $Q$ 的方程的实部与虚部，就看到，$Q$ 与 $n$ 维球面 $S^n$ 的切丛微分同胚；等价地，也与余切丛微分同胚，利用度量将余切丛与切丛识别。方程化为常微分方程
\[
\rho(f')^n+f''(f')^{n-1}(\rho^2-|\varepsilon|^2)=c>0,
\]
可以显式求解。这些就是 $Q\cong T(S^n)$ 上的 Stenzel 度量。当 $\varepsilon\to0$ 时，恢复锥形奇点度量。它们是节点的光滑化上的 Ricci 平坦 Kähler 度量；$n=2$ 时，与 Eguchi--Hanson 度量一致。$n=3$ 时，Candelas 与 de la Ossa\BookCite{ref084}首先写出了这个度量，它渐近锥形，锥底为 Einstein 流形 $T^{1,1}$。对于 $n=3$，节点的小消解上相应的度量，也已显式给出\BookCite{ref387}，同样渐近于以 Einstein 流形 $T^{1,1}$ 为底的锥形度量。

\SourceParagraph{ch07-7.1.1-014}{524}
在任意秩一紧对称空间的切丛上，都存在 Stenzel 度量的自然推广，文献\BookCite{ref388,ref328}描述了这些度量。文献\BookCite{ref213}也包含对这些度量的优美推导。它们都完备。文献\BookCite{ref028}证明了，在更高秩对称空间的复化上，完备 Ricci 平坦 Kähler 度量的存在性。
\end{example}

\SourceParagraph{ch07-7.1.1-015}{524}
除了上述具体例子，还有各种形式的拟设：假定某类例子具有一定对称性，求出对应的微分方程，寻找满足例如 Ricci 平坦、Einstein，或者具有指定完整群条件的度量。通常无法用解析方式写出一般解，甚至可能找不到个别特殊解。物理文献中有许多论文，致力于以数值方法研究这些拟设的解。在实四维超 Kähler 度量的情形中，的确有一个由 Gibbons 与 Hawking 找到的拟设，可以解析地得到一般解，并且与后面各节密切相关。

\subsection{实四维超 Kähler 度量的 Gibbons--Hawking 拟设}\label{sec:7.1.2}

\SourceParagraph{ch07-7.1.2-001}{525}
在这个拟设中，将描述如何在 $\mathbb R^3$ 的开子集上的 $S^1$ 丛上，构造等变超 Kähler 度量。下面对 Gibbons--Hawking 拟设的描述，基于 Gibbons、Hawking、Hitchin 等人的工作。

\SourceParagraph{ch07-7.1.2-002}{525}
令 $U\subseteq\mathbb R^3$ 为带有 Euclid 度量的开集，坐标为 $u_1,u_2,u_3$。令 $\pi\colon X\to U$ 为 $S^1$ 主丛，$S^1$ 作用 $S^1\times X\to X$ 写为 $(e^{\mathrm it},x)\mapsto e^{\mathrm it}\cdot x$。令 $\theta$ 为 $X$ 上的联络1形式，即取值于 $\mathfrak u(1)=\mathrm i\mathbb R$ 的1形式，在 $S^1$ 作用下不变，并满足 $\theta(\partial/\partial t)=\mathrm i$。联络的曲率为 $\mathrm d\theta=\pi^*F$，其中，$F$ 为 $U$ 上的2形式，$\mathrm iF/(2\pi)$ 代表丛的第一 Chern 类，见文献\BookCite{ref094}的附录。设 $V$ 为 $U$ 上的正实函数，满足 $*\mathrm dV=F/(2\pi\mathrm i)$。令
\[
\begin{aligned}
\omega_1&=\mathrm du_1\wedge\frac{\theta}{2\pi\mathrm i}
+V\,\mathrm du_2\wedge\mathrm du_3,\\
\omega_2&=\mathrm du_2\wedge\frac{\theta}{2\pi\mathrm i}
+V\,\mathrm du_3\wedge\mathrm du_1,\\
\omega_3&=\mathrm du_3\wedge\frac{\theta}{2\pi\mathrm i}
+V\,\mathrm du_1\wedge\mathrm du_2.
\end{aligned}
\]
这里省略从 $U$ 到 $X$ 的拉回记号。那么，$\omega_1^2=\omega_2^2=\omega_3^2$ 处处非零，$i\ne j$ 时，$\omega_i\wedge\omega_j=0$。此外，$*\mathrm dV=F/(2\pi\mathrm i)$ 蕴含，对所有 $i$，有 $\mathrm d\omega_i=0$；例如，
\[
\begin{aligned}
\mathrm d\omega_1
&=-\mathrm du_1\wedge\frac{\mathrm d\theta}{2\pi\mathrm i}
+\mathrm dV\wedge\mathrm du_2\wedge\mathrm du_3\\
&=-\mathrm du_1\wedge*\mathrm dV
+\mathrm dV\wedge\mathrm du_2\wedge\mathrm du_3=0.
\end{aligned}
\]
因此，$\omega_1,\omega_2,\omega_3$ 在 $X$ 上定义超 Kähler 度量。注意，$V$ 调和，因为 $\mathrm dF=0$ 蕴含 $*\mathrm d*\mathrm dV=0$。

\SourceParagraph{ch07-7.1.2-003}{525}
用 $\theta_0$ 表示实1形式 $\theta/(2\pi\mathrm i)$，并注意
\[
-\omega_1-\mathrm i\omega_2
=(\theta_0-\mathrm iV\,\mathrm du_3)\wedge(\mathrm du_1+\mathrm i\,\mathrm du_2).
\]
将其取为 $X$ 上的全纯2形式 $\Omega$，则由\BookRef{prop:6.24}{命题6.24}，确定 $X$ 上的可积概复结构；可积性来自 $\Omega$ 为闭形式。$\mathrm du_1+\mathrm i\,\mathrm du_2$ 与 $\theta_0-\mathrm iV\,\mathrm du_3$，在复化余切空间内张成全纯余切空间。因此，可积概复结构在余切空间上的作用为
\[
J^*(\mathrm du_1)=-\mathrm du_2,\qquad
J^*(\mathrm du_3)=-V^{-1}\theta_0.
\]
将 Kähler 形式 $\omega=\omega_3$ 看作交替张量，并使用 Riemann 度量 $g$ 满足的关系 $g(\zeta,\xi)=\omega(\zeta,J\xi)$，得到度量表达式
\[
\mathrm ds^2=V\,\mathrm du\cdot\mathrm du+V^{-1}\theta_0^2.
\]

\SourceParagraph{ch07-7.1.2-004}{525}
通常，实际上从 $U$ 上的正调和函数 $V$ 出发，要求 $-*\mathrm dV$ 代表丛的 Chern 类。于是，始终可以找到联络1形式 $\theta$，满足 $\mathrm d\theta/(2\pi\mathrm i)=*\mathrm dV$；这样的 $\theta$，在相差 $U$ 上闭1形式的拉回的意义下唯一确定。由此，得到上述超 Kähler 度量。

\SourceParagraph{ch07-7.1.2-005}{526}
\BookStatementNumber{4}
\begin{note}[numbered]\label{note:7.4}
后面需要计算这个度量的一些曲率信息。为此，可以作局部计算，采用正交归一的移动余标架
\[
V^{1/2}\mathrm du_1,\qquad V^{1/2}\mathrm du_2,\qquad
V^{1/2}\mathrm du_3,\qquad V^{-1/2}\theta_0.
\]
此外，局部地，可以将联络形式写为
\[
\theta_0=\frac{\mathrm dt}{2\pi}+A_1\,\mathrm du_1+A_2\,\mathrm du_2+A_3\,\mathrm du_3,
\]
其中，$\nabla V=\nabla\times A$。然后，可以应用 Cartan 的方法计算曲率。经过相当多的计算\BookCite{ref210}，得到曲率的范数平方
\[
\|R\|^2=12V^{-6}|\nabla V|^4+V^{-4}\Delta(|\nabla V|^2)
-6V^{-5}(\nabla V)\cdot\nabla(|\nabla V|^2).
\]
利用 $V$ 调和这一事实，经过进一步计算\BookCite{ref210}，恢复文献\BookCite{ref375}式(32)中的紧凑公式：
\[
\|R\|^2=\tfrac12V^{-1}\Delta\Delta(V^{-1}).
\]
\end{note}

\SourceParagraph{ch07-7.1.2-006}{526}
\BookStatementNumber{5}
\begin{example}\label{ex:7.5}
在\BookSectionRef{sec:6.4.2}{6.4.2}中，考虑过映射 $p\colon X=\mathbb C^2\setminus\{(0,0)\}\to\mathbb R^3\setminus\{(0,0,0)\}$，由
\[
p(z_1,z_2)=(2\operatorname{Re}(z_1z_2),\,2\operatorname{Im}(z_1z_2),\,|z_1|^2-|z_2|^2)
\]
给出。这个映射，将 $X$ 表示为 $\mathbb R^3\setminus\{(0,0,0)\}$ 上的 $S^1$ 丛，Chern 类为 $\pm1$。$S^1$ 在 $X$ 上的作用为 $e^{\mathrm it}\cdot(z_1,z_2)=(e^{\mathrm it}z_1,e^{-\mathrm it}z_2)$。还注意，若将 $p$ 与到前两个因子的投影复合，就得到 $(z_1,z_2)\mapsto2z_1z_2$，它相对于标准复结构为全纯映射。

\SourceParagraph{ch07-7.1.2-007}{526}
现在，在 $\mathbb R^3\setminus\{(0,0,0)\}$ 上选择正调和函数 $V$，使得对于 $\mathbb R^3$ 中的单位球面 $S^2$，有
\[
-\int_{S^2}*\mathrm dV=\int_{S^2}\frac{\mathrm iF}{2\pi}=\pm1,
\]
即 Chern 数正确。这里考虑的具体例子为
\[
V=e+\frac1{4\pi|u|}
=e+\frac1{4\pi\sqrt{u_1^2+u_2^2+u_3^2}},\qquad e\geq0.
\]
容易看到，积分
\[
\int_{S^2}*\mathrm d\!\left(\frac1{4\pi\sqrt{u_1^2+u_2^2+u_3^2}}\right)
\]
为 $\pm1$，取决于球面的定向。取联络形式
\[
\theta=\mathrm i\frac{\operatorname{Im}(\bar z_1\,\mathrm dz_1-\bar z_2\,\mathrm dz_2)}
{|z_1|^2+|z_2|^2}.
\]
则
\[
\frac{\mathrm d\theta}{2\pi\mathrm i}
=\frac{-(u_1\,\mathrm du_2\wedge\mathrm du_3
+u_2\,\mathrm du_3\wedge\mathrm du_1
+u_3\,\mathrm du_1\wedge\mathrm du_2)}
{4\pi(u_1^2+u_2^2+u_3^2)^{3/2}}=*\mathrm dV,
\]
正如所需。因此，得到 $X$ 上的超 Kähler 度量。对于所有 $e\geq0$，可以证明，它们延伸为 $\mathbb C^2$ 上的度量；用\BookSectionRef{sec:6.4.3}{6.4.3}的说法，就是延伸过 NUT 模型。当 $e>0$ 时，这些度量为 ALF，即渐近局部平坦型；当 $|u|\to\infty$ 时，它们接近平坦度量，同时在 $t$ 方向具有周期。$e=1$ 时，得到 Taub--NUT 度量；$e=0$ 时，则为 $\mathbb C^2$ 上的平坦度量。

\SourceParagraph{ch07-7.1.2-008}{527}
要证明 $e=0$ 时的陈述，直接计算可知，写成 $z_j=x_j+\mathrm i y_j$ 后，有
\[
\begin{aligned}
\omega_1&=\frac1\pi(\mathrm dx_2\wedge\mathrm dy_1-\mathrm dx_1\wedge\mathrm dy_2),\\
\omega_2&=\frac1\pi(\mathrm dx_1\wedge\mathrm dx_2-\mathrm dy_1\wedge\mathrm dy_2),\\
\omega_3&=\frac1\pi(\mathrm dx_1\wedge\mathrm dy_1+\mathrm dx_2\wedge\mathrm dy_2).
\end{aligned}
\]
因此，$\omega_1,\omega_2,\omega_3$ 延伸到 $\mathbb C^2$，并给出平坦度量，正如所述。

\SourceParagraph{ch07-7.1.2-009}{527}
这些计算都是局部的，因此，可以将例子推广到多瞬子度量：在 $\mathbb R^3$ 的任意有限点集上方插入 NUT 模型。也就是说，取
\[
V=e+\frac1{4\pi}\sum_j\frac1{r_j},
\]
其中，$r_j$ 表示到该点集的第 $j$ 个点的 Euclid 距离。取 $e=0$，对于两个点，再次得到 Eguchi--Hanson 度量；对于 $k$ 个不同点，得到 Gibbons--Hawking 引力多瞬子。若这些点共线，一个简单论证说明，它给出 $A_{k-1}$ 型有理双点奇点的消解上的完备 Ricci 平坦 Kähler 度量\BookCite{ref231}。如前所述，后来，这些度量被推广到任意有理双点的消解上\BookCite{ref316}。若在上述表达式中取 $e=1$，则得到多 Taub--NUT 度量；它们不是 ALE 度量。
\end{example}

\subsection{\texorpdfstring{$S^1$}{S¹} 不变的椭圆纤维化 Ricci 平坦度量}\label{sec:7.1.3}

\SourceParagraph{ch07-7.1.3-001}{527}
Gibbons--Hawking 拟设的一个应用，是构造局部椭圆纤维化上的 Ricci 平坦度量。具体做法如下。

\SourceParagraph{ch07-7.1.3-002}{527}
在 Gibbons--Hawking 拟设中，可以考虑 $U=B\times\mathbb R$ 的情形，其中 $B$ 是 $\mathbb R^2$ 的一个可缩开子集；特别地，$U$ 上的 $S^1$ 丛 $X$ 在拓扑上平凡。令 $y_1=u_1$、$y_2=u_2$，则 $y=y_1+\mathrm i y_2$ 为 $B$ 上的复坐标。$X$ 的超 Kähler 结构给出一个复结构，在此复结构下，函数 $y$ 是全纯的；也就是说，映射 $X\to B$ 全纯\BookCite{ref210}。固定纤维化的一个全纯截面后，沿纤维积分相对全纯形式 $\theta_0-\mathrm i V\,\mathrm du$，可以在 $X$ 的万有覆盖 $\widetilde X$ 上得到全纯坐标 $x=x_1+\mathrm i x_2$，使得 $X$ 的全纯二形式恰为 $\mathrm dx\wedge\mathrm dy$。这里 $u=u_3$，而 $\mathrm dx$ 与 $\theta_0-\mathrm i V\,\mathrm du$ 在模去 $\mathrm dy$ 后相等。由此，可以把 $\widetilde X$ 识别为 $T^*B$ 的一个开子域；若每条纤维上的坐标 $x$ 取遍 $\mathbb C$，则识别为整个 $T^*B$。$x,y$ 是余切丛的全纯正则坐标：点 $(x,y)$ 对应于一形式 $x\,\mathrm dy\in T_y^*B$，其中 $y$ 也表示 $B$ 中坐标为 $y$ 的点。

\SourceParagraph{ch07-7.1.3-003}{528}
此外，在这些坐标下，$X$ 上的 $S^1$ 作用提升为余切丛上沿 $x_1$ 坐标平移的 $\mathbb R$ 作用。因此，在纤维坐标取遍 $\mathbb C$ 的情形，$X$ 同构于 $(T^*B)/\mathbb Z$；一般情形则为上述开子域的相应商。细节见\BookCite{ref210}中的构造 2.6。反过来，设在 $(T^*B)/\mathbb Z$ 上给定一个上述类型的 $S^1$ 不变 Ricci 平坦 Kähler 度量，即相对于 $T^*B$ 上的全纯正则坐标 $x,y$，度量与 $x_1$ 无关。容易验证\BookCite{ref210}，它确实由 Gibbons--Hawking 拟设得到。

\SourceParagraph{ch07-7.1.3-004}{528}
我们特别关心上述度量中 $V$ 和 $\theta$ 关于 $u=u_3$ 都具有周期的情形。超 Kähler 度量能够下降到 $Y=B\times S^1$ 上相应的 $S^1$ 纤维化，当且仅当三个二形式 $\omega_1,\omega_2,\omega_3$ 在 $u$ 增加一个周期时不变；这也就是说，$u$ 的周期与 $y=u_1+\mathrm i u_2$ 无关。现在改变记号，用 $X$ 表示 $B$ 上这一 $S^1\times S^1$ 纤维化，其万有覆盖 $\widetilde X$ 与先前相同。在此周期情形，纤维的全纯坐标取遍 $\mathbb C$。由于 Kähler 形式 $\omega_3$ 在纤维 $X_y$ 上的限制恰为
\[
\mathrm du\wedge\theta_0=\mathrm du\wedge\frac{\mathrm dt}{2\pi},
\]
每条纤维的体积就是 $u$ 的周期。以 $y$ 为 $B$ 的复坐标，映射 $f\colon X\to B$ 是到 $\mathbb C$ 的开子集 $B$ 的全纯映射，纤维为椭圆曲线。此外，在这一情形，$V$ 和 $\theta$ 也关于 $u$ 周期，而其周期为一个常数，恰好等于 $f\colon X\to B$ 的椭圆纤维的体积。

\SourceParagraph{ch07-7.1.3-005}{528}
现在，将 Gibbons--Hawking 拟设应用于\BookRef{constr:6.86}{构造6.86}中构造的流形 $X$。在那里，$\overline Y=D\times S^1$，$D$ 是 $\mathbb C$ 中以原点为中心的圆盘，而对于某个 $\epsilon>0$，$S^1=\mathbb R/\epsilon\mathbb Z$。若定义在
\[
(D\times\mathbb R)\setminus(\{0\}\times\epsilon\mathbb Z)
\]
上的 $V$ 和 $\theta$ 关于 $u=u_3$ 以 $\epsilon$ 为周期，就得到 $X$ 上的超 Kähler 度量。若进一步选择 $V$，使其在原点的局部形式为 $1/(4\pi r)$ 加上一个调和函数，则在所有点 $(0,0,n\epsilon)$ 处也具有这一局部形式，度量因而延伸为相对紧化 $\overline X\to\overline Y$ 上的超 Kähler 度量。族 $\overline X\to D$ 是全纯的，可以识别为 $I_1$ 型 Kodaira 退化：奇异纤维是一条具有普通双点的椭圆曲线，在拓扑上，就是将环面上的一个圆夹缩成一个点，如\BookRef{constr:6.86}{构造6.86}所述。任意纤维的体积均为常数 $\epsilon$。这就是 Ooguri--Vafa 度量构造的基本思想，我们将在\BookSectionRef{sec:7.3.4}{7.3.4}介绍这一度量。事实上，我们希望从一个固定的 $I_1$ 型 Kodaira 退化出发，也就是预先给定 $\overline X$ 的复结构。因此，必须仔细选择 $V$ 中的常数，以保证所定义的超 Kähler 度量对应于给定的复结构。更多细节见\BookSectionRef{sec:7.3.4}{7.3.4}。

\subsection{其他完整群}\label{sec:7.1.4}

\SourceParagraph{ch07-7.1.4-001}{528}
下面简要介绍 Ricci 平坦度量的一些其他构造。

\SourceParagraph{ch07-7.1.4-002}{528}
文献\BookCite{ref389}利用扭量空间方法推广 Gibbons--Hawking 拟设，研究容许 $n$ 维环面 $T^n$ 作用的 $4n$ 维超 Kähler 度量。文献\BookCite{ref171}进一步研究了这一超 Kähler 情形，考察容许 $n$ 维环面 $T^n$ 三重全纯自由作用的 $4n$ 维超 Kähler 流形 $X$。结果表明，可以在适应环面作用的坐标中，将度量写成类似于四维 Gibbons--Hawking 拟设的形式，使非线性 Einstein 方程化为一组线性方程，其内容基本上是要求 Euclid 三维空间上的某些函数调和。对于八维流形，即 $n=2$，可以借助 Euclid 六维空间中三维线性子空间的配置，对解作几何描述\BookCite{ref171}。构造超 Kähler 度量的另一个重要方法，是文献\BookCite{ref238}提出的超 Kähler 商构造；由此得到的超 Kähler 度量的讨论见\BookCite{ref046}。从平坦四元数向量空间 $\mathbb H^n\cong\mathbb R^{4n}$ 及 $\operatorname{Sp}(n)$ 的某个子群 $G$ 的作用出发，超 Kähler 商构造给出许多例子。特别地，当 $G$ 为环面时，上述论文判定了哪些超 Kähler 度量能够由这一方法得到。

\SourceParagraph{ch07-7.1.4-003}{529}
事实上，非紧流形上完整群为 $\operatorname{SU}(n)$ 或 $\operatorname{Sp}(2n)$ 的度量，已有许多显式例子，其中不少仅见于物理文献。其他自动给出 Ricci 平坦度量的完整群，是七维的特殊完整群 $G_2$ 和八维的 $\operatorname{Spin}(7)$。对于紧致例子，同样只有度量存在性的结论，见\BookCite{ref274}\BookCite{ref314}第 11--15 章；要给出这类度量的显式例子，必须考虑非紧情形。直到不久前，已知的显式完备度量只有三个七维 $G_2$ 完整群的例子，以及一个八维 $\operatorname{Spin}(7)$ 完整群的例子\BookCite{ref076}\BookCite{ref178}；在此之前，Bryant 已构造出不完备的例子。这三个 $G_2$ 完整群的例子都是渐近锥形的，分别定义在 $S^4$ 的自对偶二形式丛、$\mathbb{CP}^2$ 的自对偶二形式丛，以及 $S^3$ 的自旋丛上；最后一个丛的拓扑为 $\mathbb R^4\times S^3$。这些度量分别相对于 Lie 群 $\operatorname{SO}(5)$、$\operatorname{SU}(3)$ 和 $\operatorname{SU}(2)\times\operatorname{SU}(2)\times\operatorname{SU}(2)$ 具有余齐性一。回顾，余齐性一的度量容许 Lie 群通过等距变换作用，其一般轨道，即主轨道，具有实余维一。文献\BookCite{ref107}证明了一个分类定理，分类对象是维数 $4n>4$、相对于某个紧单 Lie 群具有余齐性一的不可约超 Kähler 度量。特别地，若度量完备，则 $X$ 是射影 $n$ 维空间的全纯余切丛 $T^*\mathbb{CP}^n$，度量为前述 Calabi 超 Kähler 度量。

\SourceParagraph{ch07-7.1.4-004}{529}
上述三个 $G_2$ 完整群的例子，主轨道分别为 $\mathbb{CP}^3$、$\operatorname{SU}(3)/T^2$ 和 $S^3\times S^3$。它们均完备，轨道空间为 $\mathbb R_+$，并且只有一个奇异轨道，即余维更高的轨道。物理学家将这种轨道称为螺栓轨道（bolt）；NUT 则是退化为一个点的螺栓轨道。在这三个例子中，螺栓轨道就是相应丛的底空间，分别为 $S^4$、$\mathbb{CP}^2$ 和 $S^3$。实际上，可以将这些例子看作主轨道上的锥形度量的变形，用螺栓轨道代替锥顶。最后一个例子尤其值得关注：它的主轨道为 $S^3\times S^3$，由于这一情形与锥形奇点转变有关，对物理学家颇为重要。

\SourceParagraph{ch07-7.1.4-005}{530}
这些 $G_2$ 完整群的例子，都容许 Lie 群 $G$ 以低余维轨道作用。这是显式 Einstein 度量的一个普遍特征。最简单的情形是群作用只有一个轨道，此时度量流形为齐性流形。对于齐性流形上的度量，Einstein 条件可以完全写成代数条件。此外，所有 Ricci 平坦齐性流形都是平坦的\BookCite{ref044}，因此不会出现有趣的度量。其次可以考虑相对于 $G$ 具有余齐性一的情形，即一般轨道具有余维一。这时，Einstein 条件化为关于一个变量的非线性常微分方程组，此变量就是轨道空间上的参数。在所考虑的完备、不可约且非平坦的 Ricci 平坦情形，Cheeger--Gromoll 定理\BookCite{ref092}蕴含流形至多只有一个端：若有两个端，则度量沿一条直线方向分裂。由于我们考虑非紧情形，轨道空间为 $\mathbb R_+$，且只有一个奇异轨道。几何上，若主轨道为 $G/K$，则奇异轨道，即螺栓轨道，为 $G/H$，其中 $H\supset K$。在相应等变法丛模型中，若 $G$ 紧，则空间成为光滑流形的充要条件，是 $H/K$ 为法空间中的球面轨道，因而微分同胚于球面\BookCite{ref369}。在许多例子中，群 $G$ 的形式使这一条件不可能满足，于是任何由此构造的度量都不会完备。因此，若应用某种拟设，例如下述 Hitchin 拟设，寻找完整群为 $G_2$ 或 $\operatorname{Spin}(7)$ 的余齐性一度量族，并且希望度量完备，首先要解决螺栓轨道处的正则性，其次要判断解能否一直积分到无穷远。后一条件通常会减少允许的参数个数。

\SourceParagraph{ch07-7.1.4-006}{530}
近来，为寻找余齐性一的 $G_2$ 或 $\operatorname{Spin}(7)$ 完整群度量，人们提出了各种拟设，物理文献中也出现了更多此类度量的例子。不过，其中大多数都可以用 Hitchin 最近提出的一个拟设\BookCite{ref235}来描述；它在理论和实际计算上都很重要。特别地，它已促成一些特殊完整群为 $G_2$ 和 $\operatorname{Spin}(7)$ 的新余齐性一度量的构造。$G_2$ 完整群余齐性一度量的部分分类见\BookCite{ref099}。

\subsection{特殊完整群度量的 Hitchin 拟设}\label{sec:7.1.5}

\SourceParagraph{ch07-7.1.5-001}{530}
这一拟设的介绍见 2002 年 Clay 暑期学校的论文集\BookCite{ref237}，以及原始论文\BookCite{ref235}；因此，这里只解释其主要特点。出发点是稳定形式的概念。实 $n$ 维流形 $X$ 上的 $p$ 形式称为稳定的，是指在每个切空间 $V=T_PX$ 上，它都位于 $\operatorname{GL}(n)$ 在 $\bigwedge^pV^*$ 上作用的一个开轨道中。若 $p=1$，稳定性等价于处处非零；若 $p=2$ 且 $n=2m$ 或 $2m+1$，则稳定性等价于处处秩为 $2m$。更有意思的是稳定三形式：除可由低次形式对偶得到的低维情形外，新的情形只出现在六、七、八维。注意，对偶性使 $\bigwedge^pV^*$ 中的开轨道对应于 $\bigwedge^{n-p}V^*$ 中的开轨道，从而也得到稳定的 $(n-p)$ 形式。至此，所有可能的稳定形式均已包括在内，不再有其他情形。

\SourceParagraph{ch07-7.1.5-002}{531}
对于 $n=6$，$\bigwedge^3V^*$ 中存在一个开轨道，其稳定子单位分支为 $\operatorname{SL}(3,\mathbb C)$；另一个开轨道的稳定子单位分支为 $\operatorname{SL}(3,\mathbb R)\times\operatorname{SL}(3,\mathbb R)$。我们关心每点稳定子为前一种类型的稳定形式。同样，对于 $n=7$，要求稳定子为 $G_2$；对于 $n=8$，要求稳定子单位分支为 $\operatorname{PSU}(3)=\operatorname{SU}(3)/\mathbb Z_3$。后二者也对应于 $\bigwedge^3V^*$ 中的开轨道。从现在起，我们将这些稳定子类型也纳入稳定三形式的定义。

\SourceParagraph{ch07-7.1.5-003}{531}
现在限于最有意思的情形，即 $p=3$。任何给定的稳定 $p$ 形式 $\rho$ 都对应一个体积形式 $\phi(\rho)$，它关于 $\rho$ 为 $n/p$ 次齐次，可以用 $\rho$ 明确写出；这些体积形式的显式表达式见\BookCite{ref235}的附录。此外，体积形式的一阶变分正则地确定一个 $(n-p)$ 形式 $\eta_\rho$：以
\[
D\phi_\rho(\dot\rho)=\dot\rho\wedge\eta_\rho
\]
为归一化约定，齐次性给出
\[
\rho\wedge\eta_\rho=\frac np\phi(\rho).
\]
在 $n=6$、$p=3$ 时，将伴随三形式记为 $\widehat\rho$，它由 $\Omega=\rho+\mathrm i\widehat\rho$ 是 $\operatorname{SL}(3,\mathbb C)$ 保持的复 $(3,0)$ 形式这一性质确定\BookCite{ref233}；在下文所用的体积归一化下，$\phi(\rho)=\tfrac12\rho\wedge\widehat\rho$。

\SourceParagraph{ch07-7.1.5-004}{531}
当 $n=7$ 时，稳定三形式 $\rho$ 确定一个度量。具体地，可以如下定义：在切空间上取取值于最高次外幂的双线性形式
\[
b_\rho(u,v)=-\frac16\iota(u)\rho\wedge\iota(v)\rho\wedge\rho.
\]
它定义线性映射
\[
K_\rho\colon V\longrightarrow V^*\otimes\bigwedge^7V^*.
\]
体积形式为 $\phi=(\det K_\rho)^{1/9}$，该点的度量则为
\[
g(u,v)=b_\rho(u,v)(\det K_\rho)^{-1/9}.
\]
度量确定 Hodge $*$ 运算；以度量归一化记伴随形式时，有 $\widehat\rho=*\rho$\BookCite{ref233}，相应的一阶变分形式为 $\eta_\rho=\tfrac13*\rho$。$n=8$ 时也有类似情形，伴随形式按常用约定写成 $\widehat\rho=-*\rho$。

\SourceParagraph{ch07-7.1.5-005}{531}
随后，Hitchin 在固定上同调类中考虑体积泛函
\[
\Phi(\rho)=\int_X\phi(\rho).
\]
它在 $\operatorname{Diff}(X)$ 的作用下不变。Hitchin 证明，此泛函的临界点 $\rho$ 对应于方程 $\mathrm d\rho=0=\mathrm d\widehat\rho$ 的解。在 $n=6$ 时，这意味着复 $(3,0)$ 形式 $\Omega$ 闭合，并赋予 $X$ 一个典范丛平凡的复三维流形结构。$n=7$ 时，前面已指出，三形式 $\rho$ 决定度量。若 $\rho$ 是代表给定上同调类的闭稳定形式，则度量的完整群包含于 $G_2$ 的条件是 $\mathrm d*\rho=0$，见\BookCite{ref147}或\BookCite{ref274}命题 10.1.3；这正是说 $\rho$ 为 $\Phi$ 的临界点。$n=8$ 的情形仍有一些谜团：在当时，还没有已知例子。

\SourceParagraph{ch07-7.1.5-006}{531}
紧单连通流形 $M$ 上一组相关的度量，自然对应于 $M$ 上的锥 $X$ 的特殊完整群度量；这里锥是翘曲乘积 $M\times\mathbb R_+$。最值得我们关注的维数有两个。$n=6$ 时，这类度量传统上称为近 Kähler 度量，对应于锥上的 $G_2$ 度量；$n=7$ 时，则称流形具有弱完整群 $G_2$，对应于锥上的 $\operatorname{Spin}(7)$ 度量。要使锥的完整群恰为所说的类型，还须排除少数特殊情形\BookCite{ref029}。文献\BookCite{ref235}的基本思想，是将这些度量实现为某个泛函在一定约束下的临界点。

\SourceParagraph{ch07-7.1.5-007}{532}
现在集中讨论六维近 Kähler 度量，以及由此得到的七维 $G_2$ 度量；细节见\BookCite{ref235}\BookCite{ref237}。要指定 $M$ 上的辛形式，可以给出一个闭稳定二形式 $\omega$，也可以给出一个稳定四形式 $\omega^2/2$，要求它为体积泛函的临界点。后一方法更有用。现在，将先前记为 $\Phi$ 的体积泛函，在稳定三形式上记为 $V$，在稳定四形式上记为 $W$。选择上同调类 $\mathcal A\in H^3(M,\mathbb R)$ 和 $\mathcal B\in H^4(M,\mathbb R)$，并将它们看作由所有代表相应上同调类的闭形式组成的无限维仿射空间。于是，$\mathcal A\times\mathcal B$ 是仿射空间的乘积，其任意一点的切空间，是恰当三形式空间与恰当四形式空间的乘积。给定恰当三形式 $\rho=\mathrm dF$ 和恰当四形式 $\sigma=\mathrm d\beta$，定义
\[
\langle\rho,\sigma\rangle=\int_MF\wedge\sigma
=-\int_M\rho\wedge\beta.
\]
由公式 $\langle\rho_1,\sigma_2\rangle-\langle\rho_2,\sigma_1\rangle$，得到 $\mathcal A\times\mathcal B$ 上的形式辛结构。

\SourceParagraph{ch07-7.1.5-008}{532}
设在 $M$ 上给定上述稳定形式 $\rho$ 和 $\sigma=\omega^2/2$，这里不要求它们恰当。于是，$\rho$ 确定 $M$ 上的概复结构，以及一个处处非零的 $(3,0)$ 形式 $\Omega=\rho+\mathrm i\widehat\rho$；此结构可积且 $\Omega$ 全纯，当且仅当 $\mathrm d\rho=0=\mathrm d\widehat\rho$。当概复结构与二形式 $\omega$ 相容时，它们确定 $M$ 上的 Hermitian 形式；若此形式为正定度量，则称 $(\rho,\sigma)$ 为正型。若 $(\rho,\sigma)$ 是一对正型恰当稳定形式，则在约束 $\langle\rho,\sigma\rangle$ 为常数的条件下，它们成为
\[
3V(\rho)+8W(\sigma)
\]
的临界点，等价于相应度量为近 Kähler 度量。这里的具体常数并不关键，因为总可以缩放度量。直接验证可知，存在一个非零常数 $\lambda$，使得
\[
\mathrm d\widehat\rho=-4\lambda\sigma,
\qquad \mathrm d\omega=3\lambda\rho;
\]
而 $M\times\mathbb R_+$ 上的三形式
\[
\frac{r^2}{\lambda}\,\mathrm dr\wedge\omega+r^3\rho
\]
稳定，且同时闭合与余闭合，因而定义完整群包含于 $G_2$ 的度量。$M$ 上的度量为近 Kähler 度量这一事实，来自\BookCite{ref029}。

\SourceParagraph{ch07-7.1.5-009}{532}
回到一般情形，我们希望对稳定形式 $\rho$ 与 $\sigma$ 加上两个相容条件。第一个条件为 $\rho\wedge\omega=0$，其含义是：相对于 $\rho$ 诱导的概复结构，$\omega$ 为 $(1,1)$ 型；也可以说，$\rho$ 相对于 $\omega$ 为本原形式。第二个条件是体积形式成比例，即存在常数 $c$，使 $\phi(\rho)=c\phi(\sigma)$。假定这对形式为正型，这意味着 $\Omega=\rho+\mathrm i\widehat\rho$ 相对于度量具有常长度。在上述恰当稳定形式的例子中，这些条件自动成立，无须另外施加。如果 $\rho$ 和 $\sigma$ 都为各自体积泛函的临界点，就会得到 $M$ 上完整群包含于 $\operatorname{SU}(3)$ 的度量。

\SourceParagraph{ch07-7.1.5-010}{533}
文献\BookCite{ref235}\BookCite{ref237}关于 $G_2$ 完整群度量存在性的定理如下。
\begin{theorem*}\label{thm:7.hitchin-flow}
设 $(\rho(t),\sigma(t))\in\mathcal A\times\mathcal B$ 是如上所述、依赖于参数 $t$ 的一对正型稳定形式，其中 $\sigma(t)=\omega(t)^2/2$。假设它们相对于上述 $\mathcal A\times\mathcal B$ 的辛形式，沿泛函
\[
H(\rho,\sigma)=V(\rho)-2W(\sigma)
\]
的 Hamilton 流演化。这一条件化为两个方程
\[
\frac{\partial\rho}{\partial t}=\mathrm d\omega,
\qquad
\frac{\partial\sigma}{\partial t}
=\omega\wedge\frac{\partial\omega}{\partial t}
=-\mathrm d\widehat\rho.
\]
如果在 $t=t_0$ 时，形式满足相容条件
\[
\rho\wedge\omega=0,\qquad\phi(\rho)=2\phi(\sigma),
\]
则这些条件对所有 $t$ 都继续成立\BookCite{ref235}；特别地，解满足 Hamilton 约束 $H=0$。此外，三形式
\[
\varphi=\omega\wedge\mathrm dt+\rho
\]
在七维流形 $M\times(a,b)$ 上定义一个完整群包含于 $G_2$ 的度量，其中 $t_0\in(a,b)\subset\mathbb R$ 为某个开区间。上述第一个方程对应于 $\mathrm d\varphi=0$，第二个对应于 $\mathrm d*\varphi=0$。
\end{theorem*}
在 $\mathcal A$ 和 $\mathcal B$ 均为零类的情形，前述受约束泛函的正型临界点给出近 Kähler 纤维，并可以作为这一流的初始数据。

\SourceParagraph{ch07-7.1.5-011}{533}
为了实际计算，必须化为有限维的形式空间；施加对称性就能达到这一目的。第一个有意义且能够计算的情形，是余齐性一的度量。需要指出，文献\BookCite{ref099}表明，相对于紧连通 Lie 群具有余齐性一、且完整群为 $G_2$ 的条件，严格限制了可能的主轨道；除有限商以外，它们只能是 $S^6$、$\mathbb{CP}^3$、$\operatorname{SU}(3)/T^2$、$S^3\times S^3$、$S^3\times T^3$ 或 $T^6$。当时所有已知的七维、完整群恰为 $G_2$ 的显式完备度量，都具有余齐性一，主轨道为 $S^3\times S^3$、$\operatorname{SU}(3)/(\operatorname U(1)\times\operatorname U(1))$ 或 $\mathbb{CP}^3$。

\SourceParagraph{ch07-7.1.5-012}{533}
例如，取 $M=S^3\times S^3$，令 $G=\operatorname{SU}(2)\times\operatorname{SU}(2)$ 以自然方式从左作用。可以考虑 $M$ 上在 $G$ 下不变、且在交换两个因子的 $C_2$ 作用下具有规定奇偶性的形式：三形式变号，四形式不变。仍使用上述上同调类 $\mathcal A$ 和 $\mathcal B$，但现在只考虑代表这些类、满足 $G$ 不变性和上述 $C_2$ 奇偶性的形式空间。由于 $H^4(M,\mathbb Z)=0$，$\mathcal B$ 必为零类；$\mathcal A$ 则取为从两个因子拉回 $H^3(S^3,\mathbb Z)=\mathbb Z$ 的生成元之后，两者之差所代表的类。$\mathcal A$ 和 $\mathcal B$ 都是三维仿射空间。由 Hamilton 流条件导出的六个一阶方程，明确写在\BookCite{ref235}中，等价于\BookCite{ref057}导出的方程。施加更多对称性后，可以解出一个显式单参数 $G_2$ 度量族\BookCite{ref057}；在对称性最大的情形，它对应于前述\BookCite{ref076}\BookCite{ref178}的例子，主轨道为 $S^3\times S^3$。一般而言，这些度量不再渐近锥形，但参数充分为正或充分为负时，度量完备\BookCite{ref057}。通过对方程组取极限，作者还重新得到\BookCite{ref084}的 $\operatorname{SU}(3)$ 完整群变形锥形奇点度量，即 $n=3$ 时的 Stenzel 度量。文献\BookCite{ref105}也得到了相关结果。文献\BookCite{ref056}\BookCite{ref105}写出了一个相当一般的方程组，描述主轨道为 $S^3\times S^3$、完整群为 $G_2$ 的余齐性一度量；其中\BookCite{ref105}识别出三类解，在 $G_2$ 的框架中统一了上述变形锥形奇点与消解锥形奇点的 $\operatorname{SU}(3)$ 完整群度量。适当选择参数，这些解族给出正则且完备的度量。还可以指出，文献\BookCite{ref057}直接从 $G_2$ 完整群的条件出发，要求三形式闭合且余闭合；不过，如\BookCite{ref235}那样，将计算置于 Hitchin 拟设的框架中，可以使方法更清楚。

\SourceParagraph{ch07-7.1.5-013}{534}
文献\BookCite{ref215}利用 Hitchin 拟设，构造出主轨道为 $S^3\times T^3$ 的余齐性一 $G_2$ 完整群度量的显式例子；另见\BookCite{ref237}。Hamilton 流条件化为两个一阶微分方程，可以显式求解。如前所述，当时还没有已知的、主轨道为 $S^3\times T^3$ 的余齐性一 $G_2$ 完整群显式完备度量；这些例子也确实不完备。事实上，无法在保持光滑性的同时插入一个螺栓轨道。

\SourceParagraph{ch07-7.1.5-014}{534}
文献\BookCite{ref097}再次应用同一拟设，构造了更多主轨道为 $S^3\times S^3$ 的余齐性一 $G_2$ 完整群度量；另见\BookCite{ref237}。这里选取在 $\operatorname{SU}(2)\times\operatorname{SU}(2)$ 的自然作用下不变的形式，但不再要求 $\mathcal A$ 对应于从两个因子拉回 $H^3(S^3,\mathbb Z)$ 的生成元所得之差，即整数对 $(1,-1)$；而是取对应于任意整数对 $(m,n)$ 的更一般类 $\mathcal A$。形式不再要求在 $C_2$ 作用下反不变。不过，在 $\mathcal A$ 中，两个 $S^3$ 因子上不变三形式的拉回，可以按适当系数组合成一个自然闭三形式；作者只考虑与此形式之差是在 $C_2$ 作用下反不变的恰当形式。因此，所考虑的形式仿射空间仍为三维，Hitchin 拟设显然作一个小推广后仍有效。Hamilton 流条件化为两个一阶微分方程\BookCite{ref097}\BookCite{ref237}，随后得到显式解\BookCite{ref097}。同样，也可以直接定义三形式，再直接检验闭合与余闭合条件；但通过 Hitchin 拟设来论证，思路更透明。

\SourceParagraph{ch07-7.1.5-015}{534}
文献\BookCite{ref097}导出的方程，实际推广了先前找到的、相对于 $\operatorname{SU}(2)\times\operatorname{SU}(2)$ 具有余齐性一、主轨道为 $S^3\times S^3$ 的 $G_2$ 完整群度量例子。此外，\BookCite{ref215}的一阶方程组可以作为这些方程的极限重新得到，其度量也可以从这里的解得到。随后，作者论证，以这一极限方式得到的、主轨道为 $S^3\times T^3$ 的完备非奇异解只有两种：平坦的 $\mathbb R^4\times T^3$，以及 Eguchi--Hanson 空间与 $T^3$ 的直积；显然，二者的完整群都退化。特别地，作者由此推出\BookCite{ref215}中的例子不完备。这些例子以及\BookCite{ref097}中的许多其他例子，都不能在螺栓轨道处保持正则。

\SourceParagraph{ch07-7.1.5-016}{535}
对于七维弱完整群 $G_2$ 度量，从而对于八维非紧 $\operatorname{Spin}(7)$ 完整群流形，也存在类似理论\BookCite{ref235}\BookCite{ref237}。在紧七维流形 $M$ 上，寻找一个适当的四形式 $\rho$。下面简述其基本思想。$H^4(M,\mathbb R)$ 中任一给定上同调类 $\mathcal A$ 的闭形式代表组成一个仿射空间；其上有自然定义的不定度量，对应的二次型在恰当四形式上将 $\mathrm d\gamma$ 映为
\[
\int_M\gamma\wedge\mathrm d\gamma.
\]
然后，选择闭稳定四形式族 $\rho(t)$，使其沿体积泛函 $\Phi(\rho)$ 的梯度流演化。可以证明，四形式
\[
\varphi=\mathrm dt\wedge*\rho+\rho
\]
满足在 $X=M\times(a,b)$ 上定义 $\operatorname{Spin}(7)$ 完整群度量的条件，其中 $(a,b)\subset\mathbb R$ 是某个开区间。限于余齐性一的度量，这一理论也为寻找 $\operatorname{Spin}(7)$ 完整群度量提供了实用拟设。通过此方法得到的例子包括一些新的完备例子\BookCite{ref105}。文献\BookCite{ref097}\BookCite{ref105}的参考文献，选列了这一主题近期文献中的不少佳作。

\SourceParagraph{ch07-7.1.5-017}{535}
此外，文献\BookCite{ref120}将六维与七维的 Hitchin 泛函，同拓扑弦论和拓扑 M 理论联系起来。

\section{特殊 Lagrange 子流形的例子}\label{sec:7.2}
\subsection{特殊 Lagrange 子流形数学中的困难}\label{sec:7.2.1}

\SourceParagraph{ch07-7.2.1-001}{535}
现在回到\BookChapterRef{6}讨论的 SYZ 猜想的原始表述。读者或许已经注意到，第 6 章讨论特殊 Lagrange 子流形的模空间之后，随即将特殊 Lagrange 条件置于一旁。原因在于，进展最大的是 SYZ 猜想的一些较简单形式；在那里，模空间作为仿射流形，起着最重要的作用。事实上，本章将论证，这确实是一个合理的观点。不过，首先想解释证明 SYZ 猜想时遇到的一些困难。

\SourceParagraph{ch07-7.2.1-002}{535}
第一个困难，是寻找特殊 Lagrange 子流形的例子。如\BookSectionRef{sec:7.1}{7.1}所述，还没有人写出紧 Calabi--Yau 流形上的显式 Ricci 平坦度量。因此，不能指望直接解出特殊 Lagrange 方程
\[
\omega|_M=0,\qquad\operatorname{Im}\Omega|_M=0.
\]
也许可以期待，无须直接求解这些方程，而能借助其他技巧寻找这样的子流形。

\SourceParagraph{ch07-7.2.1-003}{536}
回顾\BookSectionRef{sec:1.1.4}{1.1.4}中定义的更一般概念：校准形式与校准子流形。
\begin{definition*}
Riemann 流形 $X$ 上的闭 $p$ 形式 $\phi$ 称为校准形式，是指对于 $X$ 的任意切空间中的任意有向 $p$ 维子空间 $V$，都有 $\phi|_V\leq\operatorname{Vol}(V)$。若 $X$ 的有向 $p$ 维子流形 $M$ 满足 $\phi|_M=\operatorname{Vol}(M)$，则称 $M$ 被 $\phi$ 校准。
\end{definition*}

\SourceParagraph{ch07-7.2.1-004}{536}
对于其他类型的校准子流形，存在一些构造技巧。例如，设 $X$ 是具有 Kähler 形式 $\omega$ 的复 $n$ 维 Kähler 流形，则对任意 $p\leq n$，$\omega^p/p!$ 都是校准形式。Wirtinger 定理告诉我们，实 $2p$ 维子流形 $M\subseteq X$ 满足
\[
\left.\frac{\omega^p}{p!}\right|_M=\operatorname{Vol}(M)
\]
当且仅当 $M$ 为复子流形。另一方面，众所周知，在射影代数情形，闭复子流形是代数的，由多项式方程定义。因此，很容易写出被 $\omega^p/p!$ 校准的子流形。

\SourceParagraph{ch07-7.2.1-005}{536}
遗憾的是，特殊 Lagrange 情形还没有类似的一般技巧。只有一种情形，容易写出例子：若 Calabi--Yau 流形 $X$ 定义在 $\mathbb R$ 上，即有一个反全纯对合 $\iota\colon X\to X$，则选择与这一实结构相容的 Ricci 平坦 Kähler 度量后，$\iota$ 的不动点集，也就是 $X$ 的实点集，是特殊 Lagrange 子流形。事实上，设 $\omega$ 和 $\Omega$ 分别为 Kähler 形式和全纯 $n$ 形式。选择满足 $-\iota^*[\omega]=[\omega]$ 的 Kähler 类；这样的类可以由任意 Kähler 类与其负拉回取平均得到。由于 $\Omega$ 为 $(n,0)$ 型，$\iota^*\Omega$ 为 $(0,n)$ 型。调整 $\Omega$ 的常数相位后，可以取 $\iota^*\Omega=\pm\overline\Omega$。而 $\iota^*\omega$ 仍为 $(1,1)$ 型，但变为负定。因此，$-\iota^*\omega$ 正定，仍定义一个 Ricci 平坦度量，且属于同一 Kähler 类；由 Ricci 平坦度量的唯一性，$-\iota^*\omega=\omega$。由此，对于
\[
M=\{x\in X\mid\iota(x)=x\},
\]
有 $\omega|_M=0$，并且 $\operatorname{Re}\Omega|_M=0$ 或 $\operatorname{Im}\Omega|_M=0$。无论哪种情形，适当选择相位后，$M$ 都是特殊 Lagrange 子流形。

\SourceParagraph{ch07-7.2.1-006}{536}
文献\BookCite{ref306}\BookCite{ref075}应用这一方法，在紧 Calabi--Yau 三维流形，包括五次超曲面中，得到了特殊 Lagrange 环面及其他流形的例子。遗憾的是，这一情形相当特殊，不太可能据此得到寻找环面的一般方法。

\SourceParagraph{ch07-7.2.1-007}{536}
另一种取得一定成功的构造方法，适用于能够相当精确地近似 Ricci 平坦度量的情形。Peng Lu\BookCite{ref340}以及 A. Kovalev 的未发表工作，都应用过这一方法。以\BookCite{ref340}中最简单的例子说明：设 $E$ 是周期为 $1$ 和 $\xi=\exp(2\pi\mathrm i/3)$ 的椭圆曲线，因而具有三阶自同构。$\mathbb Z_3$ 对角作用于 $E\times E\times E$，这一作用有 $3^3=27$ 个稳定子阶数为三的不动点。于是，
\[
X=(E\times E\times E)/\mathbb Z_3
\]
有 27 个奇点，局部形式为 $\mathbb C^3/\mathbb Z_3$，其中作用为
\[
(z_1,z_2,z_3)\longmapsto(\xi z_1,\xi z_2,\xi z_3).
\]
将这 27 个点爆破，可以得到奇点消解 $\pi\colon\widetilde X\to X$。

\SourceParagraph{ch07-7.2.1-008}{537}
$E\times E\times E$ 具有平坦度量，因此 $X$ 具有平坦轨形度量。可以将\BookSectionRef{sec:7.1.1}{7.1.1}提到的复三维广义 Eguchi--Hanson 度量粘入；它是 $\mathcal O_{\mathbb P^2}(-3)$ 总空间上的度量。后一空间是 $\mathbb C^3/\mathbb Z_3$ 的消解，例外集为 $\mathbb P^2$。可以将例外集的体积作为参数，把这些度量粘到 $\pi$ 的例外除子的邻域中，得到一个近似 Ricci 平坦的度量。在例外集的一个紧邻域之外，度量 Ricci 平坦，事实上就是平坦的；在例外集的一个小邻域内，它也 Ricci 平坦，与广义 Eguchi--Hanson 度量一致。只有粘合区域的度量不再 Ricci 平坦。不过，利用 R. Kobayashi\BookCite{ref307}从 Yau 关于 Ricci 平坦度量存在性的定理\BookCite{ref475}的证明中改造出的技术，可以证明：只要例外除子的体积接近零，这一度量就接近真正的 Ricci 平坦度量。随后可以证明，$X$ 上避开奇点的特殊 Lagrange 子流形，会变形为 $\widetilde X$ 上的特殊 Lagrange 子流形。

\SourceParagraph{ch07-7.2.1-009}{537}
事实上，这是在更一般情形中构造例子的最有希望的方法。特别地，\BookSectionRef{sec:7.3}{7.3}将解释，这类方法如何有可能构造出“正确的”特殊 Lagrange 环面。这里暂不多谈。

\SourceParagraph{ch07-7.2.1-010}{537}
因此，第一个困难是构造任何一个特殊 Lagrange 子流形。此外还有其他困难。假设已在复 $n$ 维 Calabi--Yau 流形 $X$ 中构造出特殊 Lagrange 环面 $M$。McLean 的变形理论结果，见\BookSectionRef{sec:6.1.1}{6.1.1}，告诉我们，这个环面在一个 $n$ 维族中变形。但接下来的问题是，甚至不知道这一族是否在 $M$ 的某个邻域中给出 $X$ 的叶状结构。事实上，回顾\BookSectionRef{sec:6.1.1}{6.1.1}，$M$ 在 $X$ 中的无穷小变形，对应于 $M$ 上的调和一形式：若法向量场 $\nu$ 决定一个无穷小变形，则对应的一形式为 $\iota(\nu)\omega|_M$。注意，若一个单参数特殊 Lagrange 变形族持续与 $M$ 相交，相应的法向量场就可能出现零点；而 $\nu$ 有零点，当且仅当 $\iota(\nu)\omega$ 有零点。在非紧 Calabi--Yau 流形中，确有特殊 Lagrange 环面的例子，其上的非零调和一形式具有零点，见\BookCite{ref353}。因此，$M$ 在变形时，有可能与原先的环面或同族其他成员相交。无法排除这一可能性，也就不能保证 $M$ 的变形在 $X$ 中局部给出叶状结构。注意，复维数二时没有这一问题，因为二维环面上的非零调和一形式从不出现零点。

\SourceParagraph{ch07-7.2.1-011}{537}
下一个问题是，McLean 的结果是局部的。也许无法让 $M$ 变形，使其变形填满整个 $X$。下面的设想说明一种可能的原因：假设子集 $Z\subseteq X$ 在一个 $n$ 维闭球 $B$ 上纤维化，写作 $Z\to B$，纤维均为特殊 Lagrange 的，但位于 $B$ 边界上方的所有纤维都是奇异的。由于 McLean 的结果完全不涉及奇异特殊 Lagrange 子流形，就无法保证还能继续变形。

\SourceParagraph{ch07-7.2.1-012}{538}
为克服这一问题，可以研究特殊 Lagrange 子流形在变形时，奇点如何产生。\BookSectionRef{sec:7.2.5}{7.2.5}将介绍这一方向的一些工作。不过，研究此问题最一般的框架，是几何测度论\BookCite{ref145}。这一主题技术性太强，无法在这里展开；一个很好的入门介绍见\BookCite{ref364}。几何测度论用可求长流的概念代替子流形。$p$ 维流是紧支撑光滑 $p$ 形式空间上的连续线性泛函；粗略地说，可求长流通过在一个性质适当的有向 $p$ 维集合上，带重数地积分 $p$ 形式而得到。这一理论旨在研究体积最小化。假设希望在给定同调类中，寻找体积最小的 $p$ 维整流，就取一列代表该类的整流，使其体积趋向所有代表的体积下确界。几何测度论的一个重要结果，Federer 紧致性定理，告诉我们：在紧环境中、质量及边界质量有统一上界时，该序列在适当的流空间拓扑中有收敛子列；体积的下半连续性与最小化性质进一步说明，这一最小化序列的极限体积，就是原序列体积的极限。因此，容易得到体积最小的可求长整流。之所以必须考虑流而非子集，是因为很难为子集空间描述一种使上述收敛结论成立的拓扑。

\SourceParagraph{ch07-7.2.1-013}{538}
特殊 Lagrange 子流形本来就使体积最小，因此，在相应紧致性条件下，这一理论的一个推论是：特殊 Lagrange 子流形的序列存在极限，该极限是体积最小、且仍为特殊 Lagrange 的可求长流。于是，问题转变为：特殊 Lagrange 可求长流究竟是什么样的？

\SourceParagraph{ch07-7.2.1-014}{538}
遗憾的是，当时对这一问题还没有进展。唯一已知的结果，是 Almgren 著名的正则性定理：除去一个 Hausdorff 余维至少为二的集合，体积最小的整流实际上就是带重数的流形。证明在他去世后才发表，印刷本长达 955 页\BookCite{ref008}。仅仅判断对于特殊 Lagrange 流，Almgren 结果是否有更短的证明，就已是一个有挑战性的问题。但要理解特殊 Lagrange 流的变形理论，还需要强得多的信息。巨大的技术挑战仍然存在，看起来，在不远的将来发现强形式 SYZ 猜想证明的可能性很小。另一方面，对特殊 Lagrange 奇点的理解已有许多进展，本节将作介绍。此外，\BookSectionRef{sec:7.3.6}{7.3.6}将给出一个稍弱形式的 SYZ 猜想；它似乎更有希望得到证明，而且大概足以满足多数用途。

\subsection{在环面作用下不变的特殊 Lagrange 子流形}\label{sec:7.2.2}

\SourceParagraph{ch07-7.2.2-001}{538}
对特殊 Lagrange 子流形的认识，大多来自某些类型的例子。在存在大量对称性时，这些例子最容易描述。文献考虑过各种情形；这里集中讨论环面作用下不变的例子。

\SourceParagraph{ch07-7.2.2-002}{539}
固定三元组 $(X,\omega,\Omega)$，其中 $X$ 为复 $n$ 维流形，$\omega$ 为 Kähler 形式，$\Omega$ 为处处非零的全纯 $n$ 形式。这里不坚持要求 Kähler 度量 Ricci 平坦，原因马上就会看到。Dominic Joyce 将这组数据称为概 Calabi--Yau 流形。

\SourceParagraph{ch07-7.2.2-003}{539}
假设环面 $T:=T^m$ 在 $X$ 上 Hamilton 作用，同时保持 $\omega$ 和 $\Omega$。这一群作用有矩映射 $\mu\colon X\to\mathfrak t^*$，其中 $\mathfrak t^*$ 为 $T$ 的 Lie 代数的对偶。向量 $\xi\in\mathfrak t$ 在 $X$ 上通过 Hamilton 向量场 $\rho(\xi)$ 无穷小作用；此向量场由
\[
\iota(\rho(\xi))\omega=\mathrm d\langle\mu,\xi\rangle
\]
定义。这里，$\langle\mu,\xi\rangle$ 是将 $x$ 映为 $\langle\mu(x),\xi\rangle$ 的函数 $X\to\mathbb R$。

\SourceParagraph{ch07-7.2.2-004}{539}
一个值得记住的模型，是 $T^2$ 在 $\mathbb C^3$ 上的作用
\[
(z_1,z_2,z_3)\longmapsto
\bigl(e^{\mathrm i(\theta_1+\theta_2)}z_1,
e^{-\mathrm i\theta_1}z_2,e^{-\mathrm i\theta_2}z_3\bigr),
\qquad(e^{\mathrm i\theta_1},e^{\mathrm i\theta_2})\in T^2.
\]
它恰为 $\operatorname{SU}(3)$ 的对角子群在 $\mathbb C^3$ 上的作用。为保持全纯三形式 $\mathrm dz_1\wedge\mathrm dz_2\wedge\mathrm dz_3$，需要 $\operatorname{SU}(3)$ 这一条件。取标准辛形式 $\omega=\tfrac{\mathrm i}{2}\sum_j\mathrm dz_j\wedge\mathrm d\bar z_j$，按照上述 Hamilton 向量场的符号约定，矩映射为
\[
\mu(z_1,z_2,z_3)
=-\frac12\bigl(|z_1|^2-|z_2|^2,\ |z_1|^2-|z_3|^2\bigr).
\]

\SourceParagraph{ch07-7.2.2-005}{539}
给定这种作用，选择 $p\in\mathfrak t^*$，就得到辛商
\[
X':=X/\!/T=\mu^{-1}(p)/T,
\]
以及诱导的辛形式 $\omega'$。此外，设 $X_1,\ldots,X_m$ 是生成 $T=T^m$ 作用的向量场，即对应于 $\mathfrak t$ 的一组基。令
\[
\Omega'=\iota(X_1,\ldots,X_m)\Omega,
\]
也就是用这些向量场收缩 $\Omega$。可以验证，它下降为 $X'$ 上的形式 $\Omega'$。因此，只要 $p$ 为 $\mu$ 的正则值，且该水平集上的 $T$ 作用自由，就得到概 Calabi--Yau 流形 $(X',\omega',\Omega')$；若只有有限稳定子，则相应得到轨形。注意，即使 $(X,\omega,\Omega)$ 为 Calabi--Yau，而不仅仅是概 Calabi--Yau，也没有理由保证 $(X',\omega',\Omega')$ 为 Calabi--Yau。

\SourceParagraph{ch07-7.2.2-006}{539}
下面这一简单结果，曾被多人以不同形式观察到，见\BookCite{ref227}\BookCite{ref185}\BookCite{ref199}\BookCite{ref266}；可以作为练习证明。
\BookStatementNumber{6}
\begin{proposition}\label{prop:7.6}
设 $M'\subseteq X'$ 为特殊 Lagrange 子流形，即 $\omega'|_{M'}=0$、$\operatorname{Im}\Omega'|_{M'}=0$。则 $M'$ 在 $\mu^{-1}(p)\subseteq X$ 中的逆像 $M$，是 $X$ 的特殊 Lagrange 子流形。
\end{proposition}

\SourceParagraph{ch07-7.2.2-007}{539}
\BookStatementNumber{7}
\begin{example}\label{ex:7.7}
对于上述 $T^2$ 在 $\mathbb C^3$ 上的作用，容易验证，$\mathbb C^3/\!/T^2$ 是一条复曲线，具有复坐标 $w=z_1z_2z_3$；将诱导全纯形式的常数实因子归一化后，$\Omega'=\mathrm dw$。$X'$ 中任一实曲线都是 Lagrange 的，因此，$X'$ 中的特殊 Lagrange 子流形恰由 $\operatorname{Im}w=\text{常数}$ 给出。于是，得到 $\mathbb C^3$ 中由方程
\[
|z_1|^2-|z_2|^2=c_1,\qquad
|z_1|^2-|z_3|^2=c_2,\qquad
\operatorname{Im}(z_1z_2z_3)=c_3
\]
定义的特殊 Lagrange 子流形。

\SourceParagraph{ch07-7.2.2-008}{540}
因此，映射 $f\colon\mathbb C^3\to\mathbb R^3$
\[
f(z_1,z_2,z_3)
=\bigl(|z_1|^2-|z_2|^2,\ |z_1|^2-|z_3|^2,
\operatorname{Im}(z_1z_2z_3)\bigr)
\]
的光滑纤维均为特殊 Lagrange 子流形。这是特殊 Lagrange 纤维化的第一个例子，由 Harvey 和 Lawson 在最初的论文\BookCite{ref226}中给出；可与\BookRef{constr:6.92}{构造6.92}比较。
\end{example}

\subsection{\texorpdfstring{$T^{n-1}$}{Tⁿ⁻¹} 不变的例子：余齐性一}\label{sec:7.2.3}

\SourceParagraph{ch07-7.2.3-001}{540}
环面几何给出许多非紧概 Calabi--Yau 流形的例子，它们是环面簇的消解或光滑化。本节所述例子的更多细节，见\BookCite{ref199}。这里解释 Gorenstein 环面奇点消解的一般框架，给出其中几个例子，以及一个光滑化的例子。假定读者熟悉环面几何的基础，参见\BookBibText{ref242}{MS1}第 7 章。

\SourceParagraph{ch07-7.2.3-002}{540}
像通常一样，令 $M\cong\mathbb Z^n$、$N=\operatorname{Hom}_{\mathbb Z}(M,\mathbb Z)$、$M_{\mathbb R}:=M\otimes_{\mathbb Z}\mathbb R$、$N_{\mathbb R}:=N\otimes_{\mathbb Z}\mathbb R$，并设 $\sigma\subseteq M_{\mathbb R}$ 是强凸有理多面锥。假定存在 $n_0\in N$ 以及 $\sigma$ 的生成元 $m_1,\ldots,m_s\in M$，使每个 $i$ 都满足 $\langle n_0,m_i\rangle=1$。此外，若 $m\in\sigma\cap(M\setminus\{0\})$，则 $\langle n_0,m\rangle\geq1$。满足这些条件时，称该锥为 Gorenstein 锥，因为由 $\sigma$ 定义的仿射环面簇 $Y_\sigma$ 是 Gorenstein 簇。在仿射复环面情形，这意味着典范类实际上平凡，因而是 Calabi--Yau 流形的一个局部版本。

\SourceParagraph{ch07-7.2.3-003}{540}
设 $P$ 为 $m_1,\ldots,m_s$ 的凸包，即
\[
P=\{m\in\sigma\mid\langle n_0,m\rangle=1\}.
\]
进一步假定，给定一个只使用整格顶点的 $P$ 的细分；令 $\Sigma$ 为它诱导的 $\sigma$ 的细分。有理多面扇 $\Sigma$ 定义环面簇 $Y_\Sigma$，以及一个保典部分消解态射 $Y_\Sigma\to Y_\sigma$，即 $K_{Y_\Sigma}=0$。涉及光滑度量时，考虑 $Y_\Sigma$ 光滑的情形；对于仍有奇点的部分消解，相应构造先在光滑部分进行。

\SourceParagraph{ch07-7.2.3-004}{540}
\BookStatementNumber{8}
\begin{example}\label{ex:7.8}
\begin{enumerate}
\item 取 $\sigma$ 由 $(1,0,1)$、$(0,1,1)$ 和 $(-1,-1,1)$ 生成。则 $n_0=(0,0,1)$，$P$ 是这三个生成元的凸包。加入点 $(0,0,1)$，就可以如\BookRef{fig:7.2}{图2}那样细分。在此情形，$Y_\sigma$ 可以识别为 $\mathbb C^3/\mathbb Z_3$，其中 $\mathbb Z_3$ 通过对角乘以三次单位根作用；$Y_\Sigma$ 是 $Y_\sigma$ 在奇点处的爆破，同构于 $\mathcal O_{\mathbb{CP}^2}(-3)$ 的总空间。
\item 取 $\sigma$ 由 $(0,0,1)$、$(1,0,1)$、$(0,1,1)$ 和 $(1,1,1)$ 生成。有两种可能的细分，如\BookRef{fig:7.3}{图3}所示。$Y_\sigma$ 为普通双点 $xy-zw=0$，这两个细分对应于普通双点的两个小消解。
\end{enumerate}
\end{example}
\begin{figure}[H]
\centering
\BookImage{01_Numbered_Figures/Chapter_07/Ch07_Fig02_PDFp541_Bookp531.pdf}
\BookNumberOnlyCaption\label{fig:7.2}
\end{figure}
\begin{figure}[H]
\centering
\BookImage{01_Numbered_Figures/Chapter_07/Ch07_Fig03_PDFp541_Bookp531.pdf}
\BookNumberOnlyCaption\label{fig:7.3}
\end{figure}

\SourceParagraph{ch07-7.2.3-005}{540}
可以利用\BookRef{prop:7.6}{命题7.6}，在 $Y_\Sigma$ 上构造特殊 Lagrange 纤维化。首先需要 $Y_\Sigma$ 上的全纯 $n$ 形式和 Kähler 形式。选择 $M$ 的基 $e_1,\ldots,e_n$，其对偶基 $e_1^*,\ldots,e_n^*$ 满足 $n_0=e_1^*+\cdots+e_n^*$。于是，这些对偶基向量对应于 $M\otimes\mathbb C^*$ 上的坐标 $z_1,\ldots,z_n$。若某个 $e_i^*\in\sigma^\vee$，则 $z_i$ 延伸为 $Y_\sigma$，从而也是 $Y_\Sigma$ 上的正则函数。一般并不能使所有 $e_i^*$ 同时位于 $\sigma^\vee$ 中，但乘积 $z_1\cdots z_n$ 总是正则的。可以取
\[
\Omega=\mathrm dz_1\wedge\cdots\wedge\mathrm dz_n;
\]
这一形式也延伸到 $Y_\Sigma$。

\SourceParagraph{ch07-7.2.3-006}{541}
\BookStatementNumber{9}
\begin{exercise}\label{exer:7.9}
验证 $\Omega$ 除一个符号外与基的选择无关。进一步证明，它在 $Y_\Sigma$ 上实际上处处非零。提示：验证沿 $\mathbb C^*\otimes M$ 在 $Y_\Sigma$ 上作用的任意余维一轨道，$\Omega$ 在一般点没有零点。
\end{exercise}

\SourceParagraph{ch07-7.2.3-007}{541}
注意，$\Omega$ 被下述 $T^{n-1}$ 作用保持：$(\theta_1,\ldots,\theta_{n-1})$ 将 $(z_1,\ldots,z_n)$ 映为
\[
\bigl(e^{\mathrm i\theta_1}z_1,\ldots,
e^{\mathrm i\theta_{n-1}}z_{n-1},
e^{-\mathrm i(\theta_1+\cdots+\theta_{n-1})}z_n\bigr).
\]

\SourceParagraph{ch07-7.2.3-008}{541}
对于 Kähler 形式 $\omega$，选择在同一作用下不变的度量。在已知 $Y_\Sigma$ 上存在 Ricci 平坦度量的情形，例如适用 Joyce 结果\BookCite{ref259}\BookCite{ref260}的商奇点，或\BookRef{ex:7.8}{例题7.8}第（2）项的锥形奇点情形，度量的不变性来自相应的唯一性结果。例如，例题 7.8 第（1）项中，固定例外 $\mathbb{CP}^2$ 的体积及渐近模型后，有唯一的 ALE Ricci 平坦度量，因此它必为 $T^2$ 不变的。同样，在第（2）项中，给定例外 $\mathbb{CP}^1$ 的体积及渐近模型后，有唯一的渐近锥形度量。这两个度量均在\BookSectionRef{sec:7.1.1}{7.1.1}提到过。

\SourceParagraph{ch07-7.2.3-009}{541}
给定 $T^{n-1}$ 不变的三元组 $(Y_\Sigma,\omega,\Omega)$，就可以应用\BookRef{prop:7.6}{命题7.6}。有矩映射 $\mu\colon Y_\Sigma\to\mathbb R^{n-1}$。给定正则值 $p\in\mathbb R^{n-1}$，相应辛商
\[
Y_p':=\mu^{-1}(p)/T^{n-1}=Y_\Sigma/\!/T^{n-1}
\]
是一条 Riemann 曲线。由 $n_0$ 定义的单项式 $w=\prod_{i=1}^nz_i$ 在 $(\mathbb C^*)^{n-1}$ 作用下不变，因此在 $Y_p'$ 上定义复坐标。将 $\Omega$ 的常数相位与诱导形式的实常数因子作相应归一化后，容易验证 $\Omega'=\mathrm dw$。由于 $Y_p'$ 为实二维，任意一维子流形自动为 Lagrange 的；特殊 Lagrange 子流形因而由 $\operatorname{Im}w=\text{常数}$ 给出。于是，映射
\[
f\colon Y_\Sigma\longrightarrow\mathbb R^n,
\qquad f(y)=\bigl(\mu(y),\operatorname{Im}w(y)\bigr)
\]
的光滑纤维为特殊 Lagrange 子流形。这是\BookRef{ex:7.7}{例题7.7}的推广。

\SourceParagraph{ch07-7.2.3-010}{542}
详细分析这些纤维化很有意思。不过，若对 $\mu$ 不进一步了解，很难做到。若矩映射 $\mu$ 为适当映射，就可以描述 $f\colon Y_\Sigma\to\mathbb R^n$ 的纤维；例如，当度量为 ALE 时，情形就是如此。这里略去推导细节，参见\BookCite{ref199}定理 2.2。关键在于，这一条件保证映射
\[
q\colon Y_\Sigma\longrightarrow\mathbb R^{n-1}\times\mathbb C,
\qquad q=(\mu,w)
\]
实际上为 $T^{n-1}$ 作用的商映射。因此，对于 $x\in\mathbb R^n$，$f^{-1}(x)$ 是在一条直线上纤维化的 $T^{n-1}$ 轨道之并。如果这些轨道都为 $(n-1)$ 维，则 $f^{-1}(x)=T^{n-1}\times\mathbb R$；若其中一个轨道降维，则 $f^{-1}(x)$ 为奇异纤维。因此，$f$ 的判别集恰好是退化 $T^{n-1}$ 轨道在 $f$ 下的像，这很容易理解。

\SourceParagraph{ch07-7.2.3-011}{542}
$Y_\Sigma$ 是 $M\otimes\mathbb C^*$ 轨道之并，这些轨道与 $\Sigma$ 的锥一一对应。设 $\tau$ 为这样的一个锥。环面几何的一个标准事实是，对应轨道被子环面 $(\mathbb R\tau\cap M)\otimes\mathbb C^*$ 固定，其中 $\mathbb R\tau$ 为 $\tau$ 张成的实向量空间。若稳定子与 $T^{n-1}$ 的交具有正维数，则这一环面分层上的 $T^{n-1}$ 轨道退化。而先前对 $T^{n-1}$ 的描述表明，它是代数环面 $n_0^\perp\otimes\mathbb C^*$ 的紧部分。因此，若 $n_0^\perp\cap\mathbb R\tau\neq\{0\}$，对应于 $\tau$ 的环面分层就是退化 $T^{n-1}$ 轨道之并。由于 $\Sigma$ 的所有一维锥都由 $P$ 中的元素生成，且 $\langle n_0,P\rangle=1$，有
\[
n_0^\perp\cap\mathbb R\tau\neq\{0\}
\quad\Longleftrightarrow\quad\dim\tau\geq2.
\]
因此，$f$ 的判别集是 $Y_\Sigma$ 中所有余维至少为二的环面分层的像。

\SourceParagraph{ch07-7.2.3-012}{542}
函数 $w$ 为 $n_0$ 定义的单项式；由于 $n_0$ 在 $\Sigma$ 的每条射线上都非零，$w$ 在每个余维至少为一的分层上都必须为零。因此，判别集包含在 $\mathbb R^{n-1}\times\{0\}$ 中，实际上就是余维至少为二的分层在矩映射 $\mu$ 下的像；在拓扑上，也就是这些分层在 $T^{n-1}$ 作用下的商。这很容易描述。例如，\BookRef{ex:7.8}{例题7.8}第（1）项的判别集如\BookRef{fig:7.4}{图4}所示；第（2）项则随小消解的选择，得到\BookRef{fig:7.5}{图5}中的一种图景。注意，内部线段的长度由 $\omega$ 的 Kähler 类决定。
\begin{figure}[H]
\centering
\BookImage{01_Numbered_Figures/Chapter_07/Ch07_Fig04_PDFp543_Bookp533.pdf}
\BookNumberOnlyCaption\label{fig:7.4}
\end{figure}
\begin{figure}[H]
\centering
\BookImage{01_Numbered_Figures/Chapter_07/Ch07_Fig05_PDFp543_Bookp533.pdf}
\BookNumberOnlyCaption\label{fig:7.5}
\end{figure}

\SourceParagraph{ch07-7.2.3-013}{542}
只要复 $n$ 维概 Calabi--Yau 流形具有 $T^{n-1}$ 作用，就会得到类似图景。不过，环面几何还给出另一个很好的情形，即环面奇点的光滑化。这里描述最简单的例子：普通双点的光滑化。

\SourceParagraph{ch07-7.2.3-014}{543}
考虑 $\mathbb C^4$ 中由 $xy-zw=t$ 给出的仿射二次曲面 $Q_t$，其中 $t$ 为固定数。则 $Q_t$ 具有 $T^2$ 作用
\[
(x,y,z,w)\longmapsto
\bigl(e^{\mathrm i\theta_1}x,e^{-\mathrm i\theta_1}y,
e^{\mathrm i\theta_2}z,e^{-\mathrm i\theta_2}w\bigr).
\]
此外，$Q_t$ 具有全纯三形式
\[
\begin{aligned}
\Omega
&=-\frac{\mathrm dx\wedge\mathrm dy\wedge\mathrm dz}{z}
=\frac{\mathrm dx\wedge\mathrm dy\wedge\mathrm dw}{w}\\
&=-\frac{\mathrm dy\wedge\mathrm dz\wedge\mathrm dw}{y}
=\frac{\mathrm dx\wedge\mathrm dz\wedge\mathrm dw}{x}.
\end{aligned}
\]
这些等式由 $Q_t$ 上的关系 $x\,\mathrm dy+y\,\mathrm dx=z\,\mathrm dw+w\,\mathrm dz$ 得到。显然，$\Omega$ 在给定的 $T^2$ 作用下不变。

\SourceParagraph{ch07-7.2.3-015}{543}
此仿射二次曲面具有 Ricci 平坦度量，即\BookRef{ex:7.3}{例题7.3}中的 Stenzel 度量，并且它也在 $T^2$ 作用下不变。因此，可以再次应用\BookRef{prop:7.6}{命题7.6}。函数 $xy$ 在 $Q_t$ 上不变，因而下降为 $\mathbb C=\mu^{-1}(p)/T^2$ 上的复坐标，其中 $p\in\mathbb R^2$ 为正则值。此外，$\Omega'=\mathrm d(xy)$；计算 $\iota(X_1,X_2)\Omega$ 就可以验证，这里 $X_1,X_2$ 为生成 $T^2$ 作用的向量场，诱导形式的常数实因子按前述方式归一化。因此，$f=(\mu,\operatorname{Im}(xy))$ 给出特殊 Lagrange 纤维化。与先前一样，一般纤维为 $\mathbb R\times T^2$，判别集是 $Q_t$ 上所有降维 $T^2$ 轨道之并的像。

\SourceParagraph{ch07-7.2.3-016}{544}
在 $t\neq0$ 时，$Q_t$ 上具有非平凡稳定子的点集，是两条双曲线：$x=y=0$、$zw=-t$，以及 $z=w=0$、$xy=t$。函数 $xy$ 在这两条双曲线上分别取值 $0$ 和 $t$；它们在矩映射下的像，是 $\mathbb R^2$ 中两条不平行的直线。因此，判别集由两条直线组成，高度分别为 $0$ 和 $\operatorname{Im}t$。注意，若 $\operatorname{Im}t=0$，这两条线相交。还可以令 $t\to0$，得到奇异的 $Q_0$，随后作小消解。从这一角度看，上述过程给出了消除普通双点的三种不同方式，见\BookRef{fig:7.6}{图6}。
\begin{figure}[H]
\centering
\BookImage{01_Numbered_Figures/Chapter_07/Ch07_Fig06_PDFp544_Bookp534.pdf}
\BookNumberOnlyCaption\label{fig:7.6}
\end{figure}

\subsection{\texorpdfstring{$\mathbb C^3$}{ℂ³} 中的 \texorpdfstring{$S^1$}{S¹} 不变例子}\label{sec:7.2.4}

\SourceParagraph{ch07-7.2.4-001}{544}
在一般 Calabi--Yau 流形上考虑 $S^1$ 作用，实在过于困难。因此，这里使用 $\mathbb C^3$ 上的作用
\[
(z_1,z_2,z_3)\longmapsto
\bigl(e^{\mathrm i\theta}z_1,e^{-\mathrm i\theta}z_2,z_3\bigr),
\]
并考察 $S^1$ 不变的特殊 Lagrange 子流形。这已经很有启发性，因为开始出现与 $T^{n-1}$ 情形很不相同的行为。Dominic Joyce 在一系列论文\BookCite{ref266,ref267,ref268}中广泛研究了这一情形。

\SourceParagraph{ch07-7.2.4-002}{544}
相对于标准辛形式 $\omega=\tfrac{\mathrm i}{2}\sum_{i=1}^3\mathrm dz_i\wedge\mathrm d\bar z_i$，为方便计算，将 Lie 代数生成元重标为标准生成元的 $-2$ 倍，取相应矩映射
\[
\mu(z_1,z_2,z_3)=|z_1|^2-|z_2|^2.
\]
下面计算 $\mu^{-1}(2a)/S^1$ 上约化的辛形式与全纯形式；因子二是一个方便的归一化。对应的生成向量场为
\[
X=2\mathrm i\bigl(\bar z_1\partial_{\bar z_1}-z_1\partial_{z_1}
-\bar z_2\partial_{\bar z_2}+z_2\partial_{z_2}\bigr).
\]
因此，
\[
\begin{aligned}
\iota(X)\Omega
&=-2\mathrm i\bigl(z_1\,\mathrm dz_2\wedge\mathrm dz_3
+z_2\,\mathrm dz_1\wedge\mathrm dz_3\bigr)\\
&=-2\mathrm i\,\mathrm d(z_1z_2)\wedge\mathrm dz_3.
\end{aligned}
\]
显然，$w=z_1z_2$ 与 $z_3$ 是不变函数，确定 $\mu^{-1}(2a)/S^1$ 上的复结构，因此，除一个常数因子外，$\Omega'$ 为 $\mathrm dw\wedge\mathrm dz_3$。

\SourceParagraph{ch07-7.2.4-003}{545}
另一方面，$\mu^{-1}(2a)/S^1$ 上的 $\omega'$ 依赖于 $a$。计算方法如下。考虑映射
\[
\pi\colon\mathbb C^3\longrightarrow\mathbb C^2,
\qquad\pi(z_1,z_2,z_3)=(z_1z_2,z_3)=(w,z_3).
\]
在点 $(z_1,z_2,z_3)$，切向量的推前为
\[
\pi_*\left(\frac\partial{\partial z_1}\right)
=z_2\frac\partial{\partial w},\qquad
\pi_*\left(\frac\partial{\partial z_2}\right)
=z_1\frac\partial{\partial w},\qquad
\pi_*\left(\frac\partial{\partial z_3}\right)
=\frac\partial{\partial z_3}.
\]
要计算点 $(w,z_3)$ 的两个切向量 $v_1,v_2$ 上的 $\omega'(v_1,v_2)$，先将此点提升到 $p\in\mu^{-1}(2a)$，再将 $v_1,v_2$ 提升为 $p$ 处切于该水平集的向量 $\widetilde v_1,\widetilde v_2$。于是，$\omega'(v_1,v_2)=\omega(\widetilde v_1,\widetilde v_2)$。具体地，在 $w\neq0$ 的坐标片中，利用 $S^1$ 作用将 $z_1$ 取为正实数，可以使用提升
\[
p=\left(
\sqrt{\frac{2a+\sqrt{4a^2+4|w|^2}}2},
\sqrt{\frac{-2a+\sqrt{4a^2+4|w|^2}}2}\frac w{|w|},
z_3\right)=(z_1,z_2,z_3).
\]
在其他非固定点，可以改用另一坐标片或取相应极限。

\SourceParagraph{ch07-7.2.4-004}{545}
向量 $\partial/\partial z_3$ 提升为 $p$ 处同名向量，它切于 $\mu^{-1}(2a)$；$\partial/\partial w$ 则提升为某个线性组合 $\alpha\partial/\partial z_1+\beta\partial/\partial z_2$。它切于水平集的条件为
\[
\left(\alpha\frac\partial{\partial z_1}
+\beta\frac\partial{\partial z_2}\right)
(|z_1|^2-|z_2|^2)=0,
\quad\text{即}\quad\alpha\bar z_1-\beta\bar z_2=0.
\]
而其推前为 $\partial/\partial w$ 的条件，是 $\alpha z_2+\beta z_1=1$。解这两个方程，得到
\[
\alpha=\frac{\bar z_2}{|z_1|^2+|z_2|^2},\qquad
\beta=\frac{\bar z_1}{|z_1|^2+|z_2|^2}.
\]
因此，
\[
\begin{aligned}
\omega'\left(\frac\partial{\partial w},\frac\partial{\partial\bar w}\right)
&=\omega\left(
\alpha\frac\partial{\partial z_1}+\beta\frac\partial{\partial z_2},
\bar\alpha\frac\partial{\partial\bar z_1}+\bar\beta\frac\partial{\partial\bar z_2}\right)\\
&=\frac{\mathrm i}2(|\alpha|^2+|\beta|^2)
=\frac{\mathrm i}2\frac1{|z_1|^2+|z_2|^2}.
\end{aligned}
\]

\SourceParagraph{ch07-7.2.4-005}{546}
由此可见，
\[
\begin{aligned}
\omega'&=\frac{\mathrm i}2\left(
\frac1{2\sqrt{a^2+|w|^2}}\,\mathrm dw\wedge\mathrm d\bar w
+\mathrm dz_3\wedge\mathrm d\bar z_3\right),\\
\Omega'&=-2\mathrm i\,\mathrm dw\wedge\mathrm dz_3.
\end{aligned}
\]
注意，当 $a=0$ 时，$\omega'$ 沿 $w=0$ 的点集奇异。这是一个重要特征：矩映射的临界点，会导致约化辛形式出现奇点。

\SourceParagraph{ch07-7.2.4-006}{546}
为寻找 $\mathbb C^2$ 中相对于这两个形式的特殊 Lagrange 子流形，选择一个 $\mathbb R^2\subset\mathbb C^2$，尝试寻找其上函数图像形式的子流形。例如，在 $\mathbb R^2$ 上取坐标 $x,y$，写成
\[
M'=\left\{(w,z_3)\in\mathbb C^2\ \middle|\
\begin{aligned}
w&=v(x,y)+\mathrm i y,\\
z_3&=x+\mathrm i u(x,y),\quad(x,y)\in S\subseteq\mathbb R^2
\end{aligned}\right\},
\]
其中 $S$ 为 $\mathbb R^2$ 中的区域。在此情形，$M'$ 相对于 $\omega',\Omega'$ 为特殊 Lagrange 的条件为
\[
\begin{aligned}
0&=\operatorname{Im}\bigl(\mathrm i\,\mathrm d(v+\mathrm i y)
\wedge\mathrm d(x+\mathrm i u)\bigr)\\
&=\mathrm dv\wedge\mathrm dx-\mathrm dy\wedge\mathrm du\\
&=\left(\frac{\partial v}{\partial y}
-\frac{\partial u}{\partial x}\right)\mathrm dy\wedge\mathrm dx,
\end{aligned}
\]
以及
\[
\begin{aligned}
0&=\frac{\mathrm d(v+\mathrm i y)\wedge\mathrm d(v-\mathrm i y)}
{2\sqrt{a^2+|w|^2}}
+\mathrm d(x+\mathrm i u)\wedge\mathrm d(x-\mathrm i u)\\
&=\frac{-2\mathrm i\,\mathrm dv\wedge\mathrm dy}{2\sqrt{a^2+|w|^2}}
-2\mathrm i\,\mathrm dx\wedge\mathrm du\\
&=2\mathrm i\left(\frac1{2\sqrt{a^2+|w|^2}}\frac{\partial v}{\partial x}
+\frac{\partial u}{\partial y}\right)\mathrm dy\wedge\mathrm dx.
\end{aligned}
\]
因此，$u,v$ 必须满足方程
\setcounter{equation}{0}
\begin{equation}\label{eq:7.1}
\frac{\partial v}{\partial y}=\frac{\partial u}{\partial x},
\end{equation}
\begin{equation}\label{eq:7.2}
\frac{\partial v}{\partial x}
+2\sqrt{a^2+y^2+v^2}\frac{\partial u}{\partial y}=0.
\end{equation}

\SourceParagraph{ch07-7.2.4-007}{546}
这两个方程可以看作 Cauchy--Riemann 方程 $\partial v/\partial y=\partial u/\partial x$、$\partial v/\partial x=-\partial u/\partial y$ 的扰动。若 $S$ 为 $\mathbb R^2$ 中的单连通区域，第一个方程使我们能够找到势函数 $f\colon S\to\mathbb R$，满足 $\partial f/\partial x=v$ 和 $\partial f/\partial y=u$。于是，$f$ 满足一个方程
\begin{equation}\label{eq:7.3}
\frac{\partial^2f}{\partial x^2}
+2\sqrt{a^2+y^2+(\partial f/\partial x)^2}
\frac{\partial^2f}{\partial y^2}=0.
\end{equation}
这是二阶拟线性椭圆方程，但当 $a=0$、$y=0$ 且 $\partial f/\partial x=0$ 时，$\partial^2f/\partial y^2$ 的系数为零，方程退化。

\SourceParagraph{ch07-7.2.4-008}{547}
下面列出一些解，见 Joyce\BookCite{ref266}第 5 节。
\BookStatementNumber{10}
\begin{example}\label{ex:7.10}
\textup{（1）}当 $a=0$ 时，\BookEqRef{eq:7.1}{7.1}与\BookEqRef{eq:7.2}{7.2}的一个显式全局解为
\[
u(x,y)=y\tanh x,\qquad
v(x,y)=\frac12y^2\operatorname{sech}^2x-\frac12\cosh^2x.
\]
注意，当 $y=0$ 时，$v(x,y)\neq0$。这意味着，由这两个函数定义的曲面 $M'\subseteq\mathbb C^2$ 上，$w$ 从不为零。因此，特殊 Lagrange 子流形 $M\subseteq\mathbb C^3$ 不与 $S^1$ 作用的不动点集 $z_1=z_2=0$ 相交。由此看到，$M$ 为 $S^1\times\mathbb R^2$，且非奇异。

\SourceParagraph{ch07-7.2.4-009}{547}
\textup{（2）}这是一个熟悉的例子。与其给出 $u,v$，不如用更有用的方式写出：对于 $a\geq0$，取
\[
M_a=\left\{(z_1,z_2,z_3)\in\mathbb C^3\ \middle|\
\begin{gathered}
|z_1|^2-|z_2|^2=2a,\quad |z_2|^2=|z_3|^2,\\
\operatorname{Im}(z_1z_2z_3)=0,\quad\operatorname{Re}(z_1z_2z_3)\geq0
\end{gathered}\right\}.
\]
$M_a$ 为 $S^1$ 不变的，事实上也为 $T^2$ 不变的。它包含在\BookRef{ex:7.7}{例题7.7}的映射的某条纤维中，但由于不等式的存在，只是该纤维的一部分。然而，$a>0$ 时，$M_a$ 又是 $S^1\times\mathbb R^2$，没有奇点。事实上，给定 $|z_3|^2$，$M_a$ 中的点满足 $|w|^2=(|z_3|^2+2a)|z_3|^2$。而 $\operatorname{Im}(wz_3)=0$、$\operatorname{Re}(wz_3)\geq0$ 将 $w,z_3$ 的选择固定为一个 $S^1$，除非 $w=z_3=0$。于是，$S^1$ 轨道集由一个 $\mathbb R^2$ 参数化，可以用极坐标描述。$a=0$ 时，得到双环面锥
\[
|z_1|^2=|z_2|^2=|z_3|^2,\qquad\operatorname{Im}(z_1z_2z_3)=0
\]
的一半。直观上，$M_a$ 中的一个 $S^1$ 轨道缩为零维轨道 $z_1=z_2=z_3=0$，从而产生奇点。

\SourceParagraph{ch07-7.2.4-010}{547}
这一例子可以用\BookEqRef{eq:7.1}{7.1}和\BookEqRef{eq:7.2}{7.2}表达。条件 $\operatorname{Im}(wz_3)=0$ 化为 $xy+uv=0$，所以在 $u\neq0$ 时，$v=-xy/u$；而 $(|z_3|^2+2a)|z_3|^2=|w|^2$ 化为
\begin{equation}\label{eq:7.4}
\begin{aligned}
(x^2+u^2+2a)(x^2+u^2)
&=y^2+v^2\\
&=y^2+\frac{x^2y^2}{u^2}.
\end{aligned}
\end{equation}
由此得到关于 $u^2$ 的三次方程
\[
u^6+(2x^2+2a)u^4
\!+\bigl(-y^2+(x^2+2a)x^2\bigr)u^2-x^2y^2=0.
\]
需要验证，对于每个 $x,y$，这一方程连同原来的约束，确定唯一允许的非负 $u^2$；在 $y\neq0$ 时，它严格为正。分为四种情形。

\SourceParagraph{ch07-7.2.4-011}{548}
\begin{enumerate}[label=\textup{(\roman*)}]
\item $x\neq0$、$y\neq0$，且有三个实根 $u^2=\gamma_1,\gamma_2,\gamma_3$。则 $\gamma_1\gamma_2\gamma_3=x^2y^2>0$，而 $\gamma_1+\gamma_2+\gamma_3=-(2x^2+2a)<0$。因此，其中两个根为负，一个为正。
\item $x\neq0$、$y\neq0$，且有一个实根 $\gamma$，另两个根为共轭复数 $\delta,\bar\delta$。与第（i）种情形一样，$\gamma|\delta|^2>0$，所以 $\gamma>0$ 为唯一的正实根。
\item $y=0$。方程为
\[
u^2(u^2+x^2+2a)(u^2+x^2)=0.
\]
唯一的非负根为零；$x=0$ 时可能为重根。
\item $x=0$，但 $y\neq0$。方程为
\[
u^2(u^4+2au^2-y^2)=0.
\]
它有根 $u^2=0$，以及一个正根，因为第二个因子在 $u^2=0$ 时为负。不过，若 $u^2=x^2=0$，则由\BookEqRef{eq:7.4}{7.4}的第一行，$y^2+v^2=0$，即 $y=v=0$。因此，在 $y\neq0$ 的假设下，只有一个允许的正根。
\end{enumerate}

\SourceParagraph{ch07-7.2.4-012}{548}
这给出函数 $u$ 的存在性。为了满足 $\operatorname{Re}(wz_3)\geq0$，在 $y\neq0$ 时，取 $u$ 的符号与 $y$ 相反，并由 $v=-xy/u$ 确定 $v$；在 $y=0$ 时，原方程给出 $u=0$、$v=x\sqrt{x^2+2a}$。于是得到\BookEqRef{eq:7.1}{7.1}、\BookEqRef{eq:7.2}{7.2}的解。这些解显然已经相当复杂，却说明存在有趣的非平凡解。可以验证，$a>0$ 时解光滑；$a=0$ 时，除 $x=y=0$ 外仍光滑，而在该点仅连续，见\BookCite{ref266}第 5 节。
\end{example}

\SourceParagraph{ch07-7.2.4-013}{548}
一般而言，给定某个 $a$ 对应的\BookEqRef{eq:7.1}{7.1}、\BookEqRef{eq:7.2}{7.2}的解 $u,v$，若 $a\neq0$，就得到非奇异特殊 Lagrange 流形，它以 $S^1$ 为纤维，在 $\mathbb R^2$ 的一个区域上纤维化。然而，$a=0$ 时，三维特殊 Lagrange 流形与零维轨道之并相交的地方可能产生奇点；零维轨道的位置为 $z_1=z_2=0$。由于在 $\mu^{-1}(0)$ 中，$|z_1|^2=|z_2|^2$，条件 $z_1=z_2=0$ 等价于 $w=0$，也就是 $y=v=0$。因此，若在 $y=0$ 的点上 $v$ 也为零，就出现奇点。在\BookRef{ex:7.10}{例题7.10}第（2）项中，这只在 $x=y=0$ 时发生，对应于锥顶。

\SourceParagraph{ch07-7.2.4-014}{548}
进一步讨论\BookEqRef{eq:7.1}{7.1}、\BookEqRef{eq:7.2}{7.2}之前，先利用上述例子，介绍 Joyce 的一个非常重要的反例。

\SourceParagraph{ch07-7.2.4-015}{548}
\BookStatementNumber{11}
\begin{example}[Joyce\BookCite{ref261}]\label{ex:7.11}
首先，将\BookRef{ex:7.10}{例题7.10}第（2）项改写为
\[
M_a=\left\{(z_1,z_2,z_3)\in\mathbb C^3\ \middle|\
\begin{gathered}
|z_1|^2-a=|z_2|^2+a=|z_3|^2+|a|,\\
\operatorname{Im}(z_1z_2z_3)=0,\quad\operatorname{Re}(z_1z_2z_3)\geq0
\end{gathered}\right\}.
\]
容易验证，对于 $a<0$，它也为特殊 Lagrange 的。$a\geq0$ 时，$M_a$ 是纤维 $f^{-1}(2a,2a,0)$ 的一半，其中 $f$ 为\BookRef{ex:7.7}{例题7.7}的映射；$a\leq0$ 时，它则是 $f^{-1}(2a,0,0)$ 的一半。现在，在 $z_3$ 方向引入平移 $c$：对每个 $c\in\mathbb C$，定义
\[
M_{a,c}=\{(z_1,z_2,z_3)\in\mathbb C^3\mid(z_1,z_2,z_3-c)\in M_a\}.
\]

\SourceParagraph{ch07-7.2.4-016}{549}
这给出 $\mathbb C^3$ 中特殊 Lagrange 子流形的三参数族。实际上，它给出 $\mathbb C^3$ 的特殊 Lagrange 纤维化，由 $F(z_1,z_2,z_3)=(a,b)$ 定义，其中 $2a=|z_1|^2-|z_2|^2$，而
\[
b=\begin{cases}
z_3,&a=z_1=z_2=0,\\
z_3-\bar z_1\bar z_2/|z_1|,&a\geq0,\ z_1\neq0,\\
z_3-\bar z_1\bar z_2/|z_2|,&a<0.
\end{cases}
\]

\SourceParagraph{ch07-7.2.4-017}{549}
\begin{proof}
验证 $a\geq0$ 时，$F^{-1}(a,c)=M_{a,c}$；$a<0$ 的论证类似。注意，$(z_1,z_2,z_3)\in M_{a,c}$ 当且仅当
\[
|z_1|^2-|z_2|^2=2a,\qquad |z_2|^2=|z_3-c|^2,
\qquad z_1z_2(z_3-c)\in[0,\infty).
\]
因此，$z_1z_2(z_3-c)=|z_1||z_2||z_3-c|=|z_1||z_2|^2$。只要 $z_1z_2\neq0$，除以 $z_1z_2$，便得
\[
c=z_3-\frac{|z_1||z_2|^2}{z_1z_2}
=z_3-\frac{\bar z_1\bar z_2}{|z_1|}.
\]
另一方面，若 $z_1z_2=0$，则或者 $a=0$，此时 $z_1=z_2=0$、$z_3=c$；或者 $a>0$，此时 $z_2=z_3-c=0$、$z_1\neq0$，再次有 $z_3-\bar z_1\bar z_2/|z_1|=c$。故 $M_{a,c}\subseteq F^{-1}(a,c)$，反向论证给出等号。
\end{proof}

\SourceParagraph{ch07-7.2.4-018}{549}
于是，$a\neq0$ 时，$F^{-1}(a,c)$ 同胚于 $\mathbb R^2\times S^1$；而 $a=0$ 时，它奇异，是 $T^2$ 上的锥。因此，$F$ 的判别集为 $\{0\}\times\mathbb C$，其余维为一！
\end{example}

\SourceParagraph{ch07-7.2.4-019}{549}
这个例子对理解 SYZ 猜想至关重要。余维一的判别集，与\BookChapterRef{6}以及上一小节的 $T^2$ 不变纤维化中看到的良好的余维二判别集，给出非常不同的图景。一旦知道余维一的行为也可能发生，就必须重新思考完整 SYZ 猜想的某些方面。下一节将回到这些观念性的问题。

\SourceParagraph{ch07-7.2.4-020}{549}
首先，回到方程\BookEqRef{eq:7.1}{7.1}与\BookEqRef{eq:7.2}{7.2}的行为。下面介绍 Joyce 对这些方程的分析结果，它们说明存在许多有趣的解。研究\BookEqRef{eq:7.3}{7.3}更方便，此时 $u=\partial f/\partial y$、$v=\partial f/\partial x$。Joyce 解决了该方程的 Dirichlet 问题。
\BookStatementNumber{12}
\begin{theorem}\label{thm:7.12}
设 $S$ 是 $\mathbb R^2$ 中具有光滑边界的有界严格凸区域，在 $(x,y)\mapsto(x,-y)$ 下不变；设整数 $k\geq0$、$\alpha\in(0,1)$。取 $a\in\mathbb R$、$\phi\in C^{k+3,\alpha}(\partial S)$。若 $a\neq0$，则存在唯一的 $f\in C^{k+3,\alpha}(\overline S)$，满足\BookEqRef{eq:7.3}{7.3}，且在 $\partial S$ 上 $f=\phi$。若 $a=0$，则存在唯一的 $f\in C^1(\overline S)$，在 $\partial S$ 上 $f=\phi$，并且二阶弱可微，以弱导数满足\BookEqRef{eq:7.3}{7.3}。
\end{theorem}

\SourceParagraph{ch07-7.2.4-021}{549}
\begin{proof}
见\BookCite{ref266}定理 7.6，以及\BookCite{ref267}定理 9.20、9.21。
\end{proof}

\SourceParagraph{ch07-7.2.4-022}{549}
当 $a\neq0$ 时，这一结果的证明，只需要拟线性椭圆偏微分方程理论中一些相当标准的技术；但 $a=0$ 时，方程可能退化，技术上非常困难，这里不提供细节。

\SourceParagraph{ch07-7.2.4-023}{550}
此定理给出 $\mathbb C^3$ 中大量带边界的特殊 Lagrange 子流形例子。$a\neq0$ 时，它们同胚于 $\overline S\times S^1$；$a=0$ 时，则可能出现奇点，位置为 $v=y=0$，也就是 $\partial f/\partial x=y=0$。对于“一般”的解，奇点为 $T^2$ 上的锥，如\BookRef{ex:7.10}{例题7.10}第（2）项所示。

\SourceParagraph{ch07-7.2.4-024}{550}
乍看起来，只知道势 $f$ 在 $\partial S$ 上的值，而不知道 $u,v$ 的值，似乎很难理解这些例子的几何。不过，Joyce 还找到一种方法，可以获得解族更明确的信息。思路如下。设有区域 $S\subseteq\mathbb R^2$，对于给定的 $a$，有\BookEqRef{eq:7.1}{7.1}、\BookEqRef{eq:7.2}{7.2}的两个解 $(u_1,v_1)$ 和 $(u_2,v_2)$。希望理解这两个解定义的 $\mathbb C^3$ 中特殊 Lagrange 子流形 $M_1,M_2$ 的交。当 $(u_1,v_1)=(u_2,v_2)$ 时，$\mathbb C^2$ 中作为两个解图像的 $M_1',M_2'$ 相交；于是 $M_1,M_2$ 沿该交点上方的一个圆相交，除非 $a=v=y=0$，此时它们交于一点。因此，尝试通过计数 $(u_1,v_1)-(u_2,v_2)$ 的零点，来计算 $M_1',M_2'$ 的交点数。

\SourceParagraph{ch07-7.2.4-025}{550}
这可以类比复分析的一个标准结果。回顾，闭曲线 $\gamma\colon S^1\to\mathbb C\setminus\{0\}$ 的缠绕数定义为
\[
\frac1{2\pi\mathrm i}\int_\gamma\frac{\mathrm dz}{z}.
\]
若 $S\subseteq\mathbb C$ 为具有适当光滑边界的有界区域，$f$ 在 $\overline S$ 的邻域全纯，且在 $\partial S$ 上非零，则 $f|_{\partial S}\colon\partial S\to\mathbb C\setminus\{0\}$ 为闭曲线，它的缠绕数等于 $f$ 在 $S$ 内部的零点总数，按重数计。由于\BookEqRef{eq:7.1}{7.1}与\BookEqRef{eq:7.2}{7.2}是 Cauchy--Riemann 方程的扰动，可以期待有类似结果，事实也确实如此。首先，需要定义零点的重数。

\SourceParagraph{ch07-7.2.4-026}{550}
\BookStatementNumber{13}
\begin{definition}\label{def:7.13}
设 $a\in\mathbb R$，$(u_1,v_1),(u_2,v_2)\colon S\to\mathbb R^2$ 为\BookEqRef{eq:7.1}{7.1}、\BookEqRef{eq:7.2}{7.2}的两个解。假定 $a\neq0$ 时，它们属于 $C^1(S)$；$a=0$ 时，属于 $C^0(S)$，并以相应弱解意义满足方程。$(u_1,v_1)-(u_2,v_2)$ 的零点 $(b,c)$ 称为奇异零点，是指 $a=c=0$，且 $v_1(b,0)=v_2(b,0)=0$。因此，奇异零点对应于 $M_1,M_2$ 在一个奇点处相交。

\SourceParagraph{ch07-7.2.4-027}{550}
设 $(b,c)\in\operatorname{Int}(S)$ 为 $(u_1,v_1)-(u_2,v_2)$ 的孤立零点。其重数定义为该差沿圆
\[
\gamma\colon S^1\longrightarrow S,\qquad
\gamma(\theta)=(b+\epsilon\cos\theta,c+\epsilon\sin\theta)
\]
的缠绕数，其中 $\epsilon$ 取得充分小，使 $(b,c)$ 为圆内唯一的零点。
\end{definition}

\SourceParagraph{ch07-7.2.4-028}{550}
Joyce 证明了下面的定理，见\BookCite{ref266}第 6 节及\BookCite{ref268}第 7 节。
\BookStatementNumber{14}
\begin{theorem}\label{thm:7.14}
在\BookRef{def:7.13}{定义7.13}的情形，
\begin{enumerate}
\item $\operatorname{Int}(S)$ 中任意孤立零点的重数，都是正整数。
\item 或者 $(u_1,v_1)\equiv(u_2,v_2)$，或者 $(u_1,v_1)-(u_2,v_2)$ 在 $\operatorname{Int}(S)$ 中至多有可数多个零点，且全部孤立。
\end{enumerate}
\end{theorem}

\SourceParagraph{ch07-7.2.4-029}{551}
这一分析使 Joyce 能够证明下面的结果，它类似于复分析中的一个标准结果。
\BookStatementNumber{15}
\begin{theorem}\label{thm:7.15}
在\BookRef{def:7.13}{定义7.13}的情形，设 $S$ 为有界区域，解在其闭包上连续，且在 $\partial S$ 的每点都有 $(u_1,v_1)\neq(u_2,v_2)$。则 $(u_1,v_1)-(u_2,v_2)$ 在 $S$ 中只有有限多个零点，全部孤立。设其重数为 $k_1,\ldots,k_n$，则沿 $\partial S$，该差关于零的缠绕数为 $\sum_{i=1}^nk_i$。
\end{theorem}

\SourceParagraph{ch07-7.2.4-030}{551}
\begin{proof}
$a\neq0$ 时，见\BookCite{ref266}定理 6.7；$a=0$ 时，见\BookCite{ref268}定理 7.7。
\end{proof}

\SourceParagraph{ch07-7.2.4-031}{551}
还有一个类似结果，在 Dirichlet 问题的解，也就是\BookRef{thm:7.12}{定理7.12}的背景中更有用。
\BookStatementNumber{16}
\begin{theorem}\label{thm:7.16}
采用\BookRef{thm:7.12}{定理7.12}的假设。给定两个不同的边界条件 $\phi_1,\phi_2\in C^{k+3,\alpha}(\partial S)$，相应的解 $f_1,f_2$ 给出 $(u_1,v_1)$、$(u_2,v_2)$。设 $\phi_1-\phi_2$ 在 $\partial S$ 上恰有 $l$ 个局部极大值和 $l$ 个局部极小值。则 $(u_1,v_1)-(u_2,v_2)$ 在 $S\setminus\partial S$ 中只有有限多个零点，全部孤立。设重数为 $k_1,\ldots,k_n$，则
\[
\sum_{i=1}^nk_i\leq l-1.
\]
\end{theorem}

\SourceParagraph{ch07-7.2.4-032}{551}
\begin{proof}
$a\neq0$ 时，见\BookCite{ref266}定理 7.11；$a=0$ 时，见\BookCite{ref268}定理 7.10。
\end{proof}

\SourceParagraph{ch07-7.2.4-033}{551}
特别地，若 $l=1$，而 $M_1,M_2$ 是分别由 $(u_1,v_1)$、$(u_2,v_2)$ 在 $S$ 内部上方定义的 $\mathbb C^3$ 中特殊 Lagrange 子流形，就可以保证 $M_1\cap M_2=\varnothing$。

\SourceParagraph{ch07-7.2.4-034}{551}
可以据此构造特殊 Lagrange 纤维化的例子。取 $S$ 为 $\mathbb R^2$ 中在 $(x,y)\mapsto(x,-y)$ 下不变的有界严格凸光滑区域，$U$ 为 $\mathbb R^3$ 的开子集。假设给定连续映射 $\Phi\colon U\to C^{3,\alpha}(\partial S)$，使得若 $(a,b,c)\neq(a,b',c')$，则 $\Phi(a,b,c)-\Phi(a,b',c')$ 只有一个局部极大值和一个局部极小值。对于 $\beta\in U$，令 $f_\beta$ 为\BookEqRef{eq:7.3}{7.3}的 Dirichlet 问题的解，满足 $f_\beta|_{\partial S}=\Phi(\beta)$，且 $a$ 为 $\beta$ 的第一个分量。当 $a\neq0$ 时，$f_\beta\in C^{3,\alpha}(\overline S)$；当 $a=0$ 时，采用\BookRef{thm:7.12}{定理7.12}给出的 $C^1$ 弱解。$f_\beta$ 定义特殊 Lagrange 子流形 $M_\beta\subseteq\mathbb C^3$，不同参数的成员互不相交：同一 $a$ 时由上述定理得到，不同 $a$ 时则位于不同矩映射水平集。可以证明，$M_\beta$ 随 $\beta$ 连续变化，因此，在 $\mathbb C^3$ 的某个子集上得到特殊 Lagrange 纤维化。

\SourceParagraph{ch07-7.2.4-035}{551}
例如，可以固定某个 $\phi\in C^{3,\alpha}(\partial S)$，取 $\Phi(a,b,c)=\phi+bx+cy$。

\subsection{具有锥形奇点的特殊 Lagrange 子流形}\label{sec:7.2.5}

\SourceParagraph{ch07-7.2.5-001}{551}
已经看到，在复三维情形，$\operatorname U(1)$ 不变的特殊 Lagrange 子流形的“一般”奇点，看起来像 $T^2$ 上的锥。因此，为更好地理解特殊 Lagrange 子流形的几何，研究具有较温和奇点的子流形的变形理论很有用。回顾\BookSectionRef{sec:6.1.1}{6.1.1}：若 $X$ 为 Calabi--Yau 流形，$M\subseteq X$ 为紧非奇异特殊 Lagrange 子流形，则 $M$ 在 $X$ 中的特殊 Lagrange 变形模空间，在对应于 $M$ 的点光滑，切空间为 $H^1(M,\mathbb R)$；更自然地说，切空间为 $M$ 上调和一形式的空间。

\SourceParagraph{ch07-7.2.5-002}{552}
先更一般地研究锥，在 $\mathbb C^m$ 中工作，$m\geq3$。
\begin{definition*}
奇异特殊 Lagrange 子流形 $C\subseteq\mathbb C^m$ 称为锥，是指对所有 $t>0$，均有 $C=tC:=\{tx\mid x\in C\}$。锥的链环为 $\Sigma=S^{2m-1}\cap C$，其中 $S^{2m-1}$ 为 $\mathbb C^m$ 的单位球面。链环的维数为 $m-1$。以下讨论原点为唯一奇点、链环光滑的锥。
\end{definition*}

\SourceParagraph{ch07-7.2.5-003}{552}
写作 $C'=C\setminus\{0\}$，则它是映射 $(0,\infty)\times\Sigma\to\mathbb C^m$、$(r,s)\mapsto rs$ 的像。在此识别下，$C'$ 的度量为 $\mathrm dr^2+r^2g_\Sigma$，其中 $g_\Sigma$ 是 $\Sigma$ 的度量。

\SourceParagraph{ch07-7.2.5-004}{552}
为稍作解释即将给出的公式，注意，$C'$ 上的 Laplace 算子 $\Delta$ 可以用 $\Sigma$ 上的 $\Delta_\Sigma$ 计算。对于函数 $f\colon\Sigma\to\mathbb R$，有
\begin{equation}\label{eq:7.5}
\Delta(r^\alpha f)
=r^{\alpha-2}\bigl(\Delta_\Sigma f-\alpha(\alpha+m-2)f\bigr).
\end{equation}
利用相对于局部坐标 $x_i$、度量分量为 $g_{ij}$ 时 Laplace 算子的显式公式，就能得到这一结果：
\[
\Delta f=-\frac1{\sqrt{\det(g_{ij})}}
\sum_{k,l}\frac\partial{\partial x_k}
\left(\sqrt{\det(g_{ij})}\,g^{kl}\frac{\partial f}{\partial x_l}\right).
\]

\SourceParagraph{ch07-7.2.5-005}{552}
因此，只要 $\alpha(\alpha+m-2)$ 为 $\Delta_\Sigma$ 的特征值、$f$ 为相应特征函数，$r^\alpha f$ 就是 $C'$ 上的调和函数。令
\[
\mathcal D_\Sigma
=\{\alpha\in\mathbb R\mid\alpha(\alpha+m-2)
\text{ 为 }\Delta_\Sigma\text{ 的特征值}\}.
\]
由于 $\Sigma$ 紧，$\mathcal D_\Sigma$ 离散。每个特征值还有重数，即相应特征函数空间的维数；用 $m_\Sigma\colon\mathcal D_\Sigma\to\mathbb N$ 表示这一重数函数。$\mathcal D_\Sigma$ 的意义，可以粗略解释如下。Lockhart 和 McOwen\BookCite{ref337}以及 Melrose\BookCite{ref357}发展了带端流形上的椭圆偏微分方程理论。若考虑 $C'$ 上调和函数的空间，要求在 $r=0$ 附近的渐近行为类似于 $r^\alpha f$，其中 $f$ 为 $\Sigma$ 上的函数，则相应带权空间中的解维数依赖于 $\alpha$；当 $\alpha$ 穿过 $\mathcal D_\Sigma$ 的元素时，维数会发生变化。几页之后将提供更多细节。

\SourceParagraph{ch07-7.2.5-006}{552}
最后，设 $G$ 为 $\operatorname{SU}(m)$ 中保持 $C$ 的子群。
\begin{definition*}
Joyce 将锥 $C$ 的稳定性指标定义为
\[
\operatorname{s\!\text{-}\!ind}(C)
=-b_0(\Sigma)-m^2-2m+1+\dim G
+\sum_{\alpha\in\mathcal D_\Sigma\cap[0,2]}m_\Sigma(\alpha).
\]
这里，$b_0(\Sigma)$ 是零阶 Betti 数，即 $\Sigma$ 的连通分支个数。若 $\operatorname{s\!\text{-}\!ind}(C)=0$，则称 $C$ 稳定。
\end{definition*}

\SourceParagraph{ch07-7.2.5-007}{553}
已知许多特殊 Lagrange 锥的例子，第一个就是环面上的锥
\[
C=\{0\}\cup\left\{(z_1,\ldots,z_m)\in\mathbb C^m\ \middle|\
\begin{gathered}
\operatorname{Im}(z_1\cdots z_m)=0,\quad
\operatorname{Re}(z_1\cdots z_m)>0,\\
|z_1|=\cdots=|z_m|
\end{gathered}\right\}.
\]
链环 $\Sigma$ 是环面
\[
T^{m-1}=\left\{
\frac1{\sqrt m}\bigl(e^{\mathrm i\theta_1},\ldots,
e^{\mathrm i\theta_{m-1}},e^{-\mathrm i(\theta_1+\cdots+\theta_{m-1})}\bigr)
\ \middle|\ \theta_1,\ldots,\theta_{m-1}\in\mathbb R
\right\}\subseteq\mathbb C^m.
\]

\SourceParagraph{ch07-7.2.5-008}{553}
环面 $T^{m-1}=\mathbb R^{m-1}/2\pi\mathbb Z^{m-1}$ 上的诱导度量计算如下。它的切向量 $\partial/\partial\theta_i$ 映为 $\mathbb C^m$ 中的向量
\[
\frac1{\sqrt m}\left(
\begin{aligned}
&-\sin\theta_i\frac\partial{\partial x_i}
+\cos\theta_i\frac\partial{\partial y_i}\\
&+\sin(\theta_1+\cdots+\theta_{m-1})\frac\partial{\partial x_m}
-\cos(\theta_1+\cdots+\theta_{m-1})\frac\partial{\partial y_m}
\end{aligned}\right),
\]
这里 $z_i=x_i+\mathrm i y_i$ 为 $\mathbb C^m$ 的坐标。因此，
\[
g_{ij}=g\left(\frac\partial{\partial\theta_i},
\frac\partial{\partial\theta_j}\right)=\frac1m(\delta_{ij}+1),
\qquad g^{ij}=m\delta_{ij}-1.
\]
此平坦环面上 $\Delta$ 的特征函数当然具有形式
\[
f=e^{\mathrm i(n_1\theta_1+\cdots+n_{m-1}\theta_{m-1})},
\qquad(n_1,\ldots,n_{m-1})\in\mathbb Z^{m-1},
\]
且 $\Delta f=\sum_{k,l}g^{kl}n_kn_lf$。所以，$f$ 的特征值为
\[
m\sum_{i=1}^{m-1}n_i^2-\sum_{i,j=1}^{m-1}n_in_j.
\]

\SourceParagraph{ch07-7.2.5-009}{553}
由此，经过繁琐计算可得，对于 $m=3,4,5,6,\ldots$，分别有
\[
\sum_{\alpha\in\mathcal D_\Sigma\cap[0,2]}m_\Sigma(\alpha)
=13,27,51,93,\ldots.
\]
另一方面，保持 $C$ 的 Lie 群的单位分支为 $\operatorname U(1)^{m-1}$。因此，
\[
\operatorname{s\!\text{-}\!ind}(C)=0,6,20,50,\ldots.
\]
事实上，$C$ 仅在 $m=3$ 时稳定。这是当时唯一已知的非平凡稳定奇异锥例子。

\SourceParagraph{ch07-7.2.5-010}{554}
接下来定义具有锥形奇点的特殊 Lagrange 子流形。
\begin{definition*}\label{def:7.conical-singularity}
设 $X$ 为 Calabi--Yau 流形，$M\subseteq X$ 为具有孤立奇点的特殊 Lagrange 子流形。设 $x\in M$ 为奇点，$B_R$ 为 $\mathbb C^m$ 中半径 $R$ 的球，带标准辛形式；$D\colon B_R\to X$ 为辛嵌入，满足 $D(0)=x$，其在原点的线性识别与给定 Calabi--Yau 结构相容。也就是说，$D$ 给出 $x$ 的 Darboux 邻域。设 $C$ 为 $\mathbb C^m$ 中的特殊 Lagrange 锥，链环为 $\Sigma$。令 $\varphi\colon\Sigma\times(0,R)\to B_R$ 为嵌入 $\varphi(y,r)=ry$。假定存在光滑映射 $\varphi'\colon\Sigma\times(0,R')\to B_R$，使 $D\circ\varphi'$ 微分同胚地映到 $M$ 在 $x$ 附近去掉 $x$ 后的部分。最后，取 $2<\mu<3$，且 $(2,\mu]\cap\mathcal D_\Sigma=\varnothing$。若在 $r\to0$ 时，
\[
|\varphi-\varphi'|=O(r^{\mu-1}),\qquad
|\nabla(\varphi-\varphi')|=O(r^{\mu-2}),
\]
则称 $M$ 在 $x$ 处有收敛率 $\mu$、模型锥 $C$ 的锥形奇点。这里，范数及 Levi--Civita 联络相对于 $\Sigma\times(0,R)$ 上的锥度量计算。
\end{definition*}

\SourceParagraph{ch07-7.2.5-011}{554}
实质上，这就是说，在适当度量意义下，$M$ 在 $x$ 附近接近锥 $C$。注意，选择 $\mu<3$，是因为保持原点线性识别的 Darboux 坐标变化，会使 $|\varphi-\varphi'|$ 改变 $O(r^2)$。涉及 $\mathcal D_\Sigma$ 的约束，则是为使下述变形理论成立。

\SourceParagraph{ch07-7.2.5-012}{554}
Joyce 对具有锥形奇点的特殊 Lagrange 子流形，证明了若干结果。下面介绍最重要的几个。

\SourceParagraph{ch07-7.2.5-013}{554}
第一个结果推广了 McLean 的变形理论。给定紧特殊 Lagrange 子流形 $M\subseteq X$，它只有锥形奇点 $x_1,\ldots,x_n$，模型锥为 $C_1,\ldots,C_n$，可以尝试在保持这些奇点类型的同时变形 $M$。称特殊 Lagrange 子流形 $M'\subseteq X$ 为 $M$ 的这种变形，是指它有奇点 $x_1',\ldots,x_n'$，对应的锥 $C_1',\ldots,C_n'$ 分别与原锥 $\operatorname{SU}(m)$ 等价，且保持相应收敛率；存在同胚 $\iota\colon M\to M'$，满足 $\iota(x_i)=x_i'$，在奇点之外为微分同胚，并且嵌入 $M\hookrightarrow X$ 与复合映射 $M\xrightarrow{\iota}M'\hookrightarrow X$ 作为连续映射同伦，即有连接两者的连续映射族 $M\to X$。

\SourceParagraph{ch07-7.2.5-014}{554}
令 $\mathcal M_{\mathrm{SL}}(M)$ 表示这类变形的集合；可以赋予此空间一个拓扑。Joyce 在\BookCite{ref263}中如下描述它在 $M$ 附近的结构。令 $M^\circ=M\setminus\{x_1,\ldots,x_n\}$，有从紧支撑 de Rham 上同调到通常 de Rham 上同调的映射
\[
H_c^1(M^\circ,\mathbb R)\longrightarrow H^1(M^\circ,\mathbb R).
\]
将其像记为 $I_M$。存在一个维数为 $\sum_{i=1}^n\operatorname{s\!\text{-}\!ind}(C_i)$ 的空间 $O_M$，称为阻碍空间，以及 $C^\infty$ 映射 $\Phi\colon U\to O_M$，其中 $U$ 为 $I_M$ 中零的某个开邻域，使 $\Phi^{-1}(0)$ 同胚于 $\mathcal M_{\mathrm{SL}}(M)$ 中 $M$ 的一个开邻域。

\SourceParagraph{ch07-7.2.5-015}{554}
非奇异情形没有这些阻碍。为粗略理解它们如何出现，下面给出证明概要。
\begin{proof}[证明概要]
\SourceParagraph{ch07-7.2.5-016}{555}
首先证明一个 Lagrange 邻域定理：$M^\circ$ 在 $X$ 中的某个管状邻域，辛同胚于余切丛 $T^*M^\circ$ 中零截面的某个开邻域。因此，避开变形后的奇点，$M$ 的任意小变形都可以识别为 $T^*M^\circ$ 的一个截面的图像。由于希望截面为 Lagrange 的，它实际上是某个闭一形式 $\alpha$ 的图像，定义上同调类 $[\alpha]\in H^1(M^\circ,\mathbb R)$。当然，$\alpha$ 来自 $M$ 的小变形，还必须满足额外条件，即在奇点附近满足相应衰减条件。这迫使 $\alpha$ 在各奇点邻域实际上恰当，因而 $[\alpha]$ 位于 $H_c^1(M^\circ,\mathbb R)$ 的像中。于是，选择紧支撑闭一形式的向量空间 $I_M$，代表上述像，再选择与 $\alpha$ 同类的 $\beta\in I_M$。对于 $M^\circ$ 上某个 $C^\infty$ 函数 $f$，有 $\alpha=\beta+\mathrm df$。变形后的子流形为特殊 Lagrange 的条件，可以写成二阶非线性偏微分方程。具体形式依赖于复结构，但在 $\beta=f=0$ 处，其线性化为 $\mathrm d^*(\beta+\mathrm df)=0$，即 $\Delta f=-\mathrm d^*\beta$。

\SourceParagraph{ch07-7.2.5-017}{555}
因此，必须理解 $M^\circ$ 上 Laplace 算子的余核。利用 Lockhart 和 McOwen\BookCite{ref337}发展出的柱形端流形的偏微分方程理论，考虑带权 Sobolev 函数空间 $L^p_{k,\alpha}(M^\circ)$。其中，函数及前 $k$ 阶导数以相应权重属于 $L^p$；在模型锥 $C_i$、链环 $\Sigma_i$ 的奇点 $x_i$ 附近，权重描述形如 $r^\alpha g$ 的渐近行为，$g$ 为 $\Sigma_i$ 上的函数。由\BookEqRef{eq:7.5}{7.5}可见，$\Delta$ 将 $L^p_{k,\mu}(M^\circ)$ 映到 $L^p_{k-2,\mu-2}(M^\circ)$。取 $\mu\in(2,3)$，且对每个 $i$，$\mathcal D_{\Sigma_i}\cap(2,\mu]=\varnothing$，则\BookCite{ref337}给出
\[
\dim\operatorname{coker}\bigl(
\Delta\colon L^p_{k,\mu}(M^\circ)
\longrightarrow L^p_{k-2,\mu-2}(M^\circ)\bigr)
=\sum_{i=1}^n\sum_{\alpha\in\mathcal D_{\Sigma_i}\cap[0,2]}
m_{\Sigma_i}(\alpha).
\]
这还不等于 $\sum_i\operatorname{s\!\text{-}\!ind}(C_i)$。不过，允许奇点移动、模型锥作相应旋转，并允许 $f$ 在各端具有“常数加 $r^\mu g$”的渐近行为后，就得到正确的余核，它本质上就是阻碍空间 $O_M$。

\SourceParagraph{ch07-7.2.5-018}{555}
可见，最终导致阻碍的是 $\Delta_{\Sigma_i}$ 的特征函数。
\end{proof}

\SourceParagraph{ch07-7.2.5-019}{555}
注意，若 $M$ 的所有奇点都稳定，则 $\mathcal M_{\mathrm{SL}}(M)$ 实际上光滑。在复三维、锥形奇点恰为\BookRef{ex:7.10}{例题7.10}所述 $T^2$ 锥的情形，正是如此。

\SourceParagraph{ch07-7.2.5-020}{555}
实际上，还希望更多地了解特殊 Lagrange 子流形模空间。特别地，给定奇异特殊 Lagrange 子流形 $M\subseteq X$，希望知道它能否变形为非奇异子流形。若 $M$ 只有锥形奇点，可以再次利用 Joyce 的结果\BookCite{ref264,ref265}。首先，需要研究 $\mathbb C^m$ 中渐近锥形的特殊 Lagrange 子流形。

\SourceParagraph{ch07-7.2.5-021}{556}
\begin{definition*}\label{def:7.ac-special-lagrangian}
设 $C\subseteq\mathbb C^m$ 为只在原点有孤立奇点的特殊 Lagrange 锥，链环 $\Sigma=C\cap S^{2m-1}$；像通常一样，设 $\iota\colon\Sigma\times(0,\infty)\to\mathbb C^m$ 为嵌入 $\iota(\sigma,r)=r\sigma$。特殊 Lagrange 子流形 $M\subseteq\mathbb C^m$ 称为具有收敛率 $\lambda<2$、模型锥 $C$ 的渐近锥形子流形，是指存在紧子集 $K\subseteq M$，以及对于某个 $T>0$ 的微分同胚 $\varphi\colon\Sigma\times(T,\infty)\to M\setminus K$，使得在 $r\to\infty$ 时，
\[
|\varphi-\iota|=O(r^{\lambda-1}),\qquad
|\nabla(\varphi-\iota)|=O(r^{\lambda-2}).
\]
这里的梯度与范数相对于 $\Sigma\times(T,\infty)$ 上的锥度量计算。
\end{definition*}

\SourceParagraph{ch07-7.2.5-022}{556}
\BookStatementNumber{17}
\begin{example}\label{ex:7.17}
这类例子有不少，但这里只使用渐近于环面锥的例子，它们是\BookRef{ex:7.10}{例题7.10}第（2）项的自然推广。令 $a=(a_1,\ldots,a_m)$，所有 $a_i\geq0$，恰有两个为零，其余为正。定义
\[
M^a=\left\{(z_1,\ldots,z_m)\in\mathbb C^m\ \middle|\
\begin{gathered}
\operatorname{Re}(z_1\cdots z_m)\geq0,
\quad\operatorname{Im}(z_1\cdots z_m)=0,\\
|z_1|^2-a_1=\cdots=|z_m|^2-a_m
\end{gathered}\right\}.
\]
与前述例子一样，$M^a$ 实际上是无边界流形，渐近于锥 $M^0$，收敛率为零。

\SourceParagraph{ch07-7.2.5-023}{556}
注意，$M^a$ 同胚于 $T^{m-2}\times\mathbb R^2$。事实上，不妨设 $a_{m-1}=a_m=0$。则 $1\leq i\leq m-2$ 时，$|z_i|^2>0$，因而 $M^a$ 上的 $T^{m-2}$ 作用
\[
\begin{split}
(z_1,\ldots,z_m)\longmapsto
\bigl(&e^{\mathrm i\theta_1}z_1,\ldots,
e^{\mathrm i\theta_{m-2}}z_{m-2},\\
&e^{-\mathrm i(\theta_1+\cdots+\theta_{m-2})}z_{m-1},z_m\bigr)
\end{split}
\]
是自由的。此作用的商为 $\mathbb C=\mathbb R^2$，复坐标为 $z_m$。

\SourceParagraph{ch07-7.2.5-024}{556}
因此，两个零 $a_i$ 的每一种选择，都给出一个模型锥为 $M^0$ 的 $(m-2)$ 维渐近锥形特殊 Lagrange 子流形族；这样的选择共有 $\binom m2$ 种。特别地，这给出 $M^0$ 的 $\binom m2$ 种不同光滑化！
\end{example}

\SourceParagraph{ch07-7.2.5-025}{556}
给定收敛率 $\lambda<2$、模型锥 $C$ 的渐近锥形特殊 Lagrange 子流形 $M$，用 $\mathcal M_{\mathrm{SL}}^\lambda(M)$ 表示所有收敛率 $\lambda$、模型锥 $C$ 的渐近锥形特殊 Lagrange 子流形 $M'$ 的模空间，其中 $M'$ 作为 $\mathbb C^m$ 中的渐近锥形子流形，与 $M$ 同痕。

\SourceParagraph{ch07-7.2.5-026}{556}
注意，若 $M$ 的收敛率为 $\lambda$，则对任意 $\lambda'\in[\lambda,2)$，其收敛率也可以取为 $\lambda'$；$\mathcal M_{\mathrm{SL}}^\lambda(M)$ 可能依赖于 $\lambda$。

\SourceParagraph{ch07-7.2.5-027}{556}
Marshall\BookCite{ref349}和 Pacini\BookCite{ref385}证明了下述结果，另见\BookCite{ref262}定理 6.8。
\BookStatementNumber{18}
\begin{theorem}\label{thm:7.18}
设 $M$ 为 $\mathbb C^m$ 中渐近锥形的特殊 Lagrange 子流形，模型锥 $C$、收敛率 $\lambda<2$。令 $\Sigma=C\cap S^{2m-1}$，$\mathcal D_\Sigma,m_\Sigma$ 为 $\Delta_\Sigma$ 确定的数据。则
\begin{enumerate}
\item 若 $\lambda\in(0,2)\setminus\mathcal D_\Sigma$，则 $\mathcal M_{\mathrm{SL}}^\lambda(M)$ 为流形，且
\[
\dim\mathcal M_{\mathrm{SL}}^\lambda(M)
=b_1(M)-b_0(M)
+\sum_{\alpha\in\mathcal D_\Sigma\cap[0,\lambda]}m_\Sigma(\alpha).
\]
\item 若 $\lambda\in(2-m,0)$，则 $\mathcal M_{\mathrm{SL}}^\lambda(M)$ 为维数 $b_{m-1}(M)$ 的流形。
\end{enumerate}
\end{theorem}

\SourceParagraph{ch07-7.2.5-028}{557}
这里，$b_i(M)=\dim_{\mathbb R}H^i(M,\mathbb R)$ 为第 $i$ 阶 Betti 数。注意，若 $0<\lambda<\inf(\mathcal D_\Sigma\cap(0,\infty))$，第一项求和只有 $m_\Sigma(0)$ 作出贡献，它是 $\Sigma$ 上调和函数空间的维数。由于 $\Sigma$ 紧，这就是连通分支数 $b_0(\Sigma)$。

\SourceParagraph{ch07-7.2.5-029}{557}
\BookStatementNumber{19}
\begin{example}\label{ex:7.19}
在\BookRef{ex:7.17}{例题7.17}中，取 $0<\lambda<\inf(\mathcal D_\Sigma\cap(0,\infty))$。由于 $b_0(M^a)=b_0(\Sigma)=1$，有
\[
\dim\mathcal M_{\mathrm{SL}}^\lambda(M^a)
=b_1(T^{m-2}\times\mathbb R^2)=m-2.
\]
它与参数 $a$ 的选择空间的维数相同。
\end{example}

\SourceParagraph{ch07-7.2.5-030}{557}
为讨论后面要用的光滑化结果，需要定义一个上同调不变量。给定渐近锥形特殊 Lagrange 子流形 $M\subseteq\mathbb C^m$，辛形式 $\omega$ 在 $M$ 上的限制当然为零。因此，它代表相对 de Rham 上同调群 $H^2(\mathbb C^m,M;\mathbb R)$ 的一个元素。此群用 $\mathbb C^m$ 中在 $M$ 上限制为零的闭二形式，模去在 $M$ 上限制为零的一形式的微分来计算。而相对上同调正合列为
\[
\begin{split}
0=H^1(\mathbb C^m,\mathbb R)&\longrightarrow H^1(M,\mathbb R)
\longrightarrow H^2(\mathbb C^m,M;\mathbb R)\\
&\longrightarrow H^2(\mathbb C^m,\mathbb R)=0.
\end{split}
\]
所以，$H^1(M,\mathbb R)\cong H^2(\mathbb C^m,M;\mathbb R)$。另一方面，存在紧集 $K\subseteq M$，使 $M\setminus K\cong\Sigma\times(T,\infty)$，从而有自然限制映射 $H^1(M,\mathbb R)\to H^1(\Sigma,\mathbb R)$。于是，$\omega$ 确定 $H^1(M,\mathbb R)$ 的一个类，将其在 $H^1(\Sigma,\mathbb R)$ 中的像记为 $Y(M)$。

\SourceParagraph{ch07-7.2.5-031}{557}
\BookStatementNumber{20}
\begin{example}\label{ex:7.20}
在\BookRef{ex:7.17}{例题7.17}中，$\Sigma\cong T^{m-1}$，所以 $H^1(\Sigma,\mathbb R)=\mathbb R^{m-1}$。计算 $Y(M)$，可以利用相对同调的对偶序列
\[
\begin{split}
0=H_2(\mathbb C^m,\mathbb R)&\longrightarrow H_2(\mathbb C^m,M;\mathbb R)
\longrightarrow H_1(M,\mathbb R)\\
&\longrightarrow H_1(\mathbb C^m,\mathbb R)=0.
\end{split}
\]
因此，来自 $[\omega]\in H^2(\mathbb C^m,M;\mathbb R)$ 的 $H^1(M,\mathbb R)$ 类，可以通过在边界位于 $M$ 上的曲面上积分 $\omega$ 计算。

\SourceParagraph{ch07-7.2.5-032}{557}
不妨对例题 7.17 中 $a_1,\ldots,a_{m-2}>0$、$a_{m-1}=a_m=0$ 的光滑化 $M^a$ 计算。对于某个固定 $r>0$，用渐近端的截面代表 $\Sigma$：
\[
\left\{\bigl(e^{\mathrm i\theta_1}\sqrt{r+a_1},\ldots,
e^{\mathrm i\theta_{m-1}}\sqrt{r+a_{m-1}},
e^{-\mathrm i(\theta_1+\cdots+\theta_{m-1})}\sqrt r\bigr)\right\}.
\]
$H_1(\Sigma,\mathbb R)$ 的生成元为环路，对于 $1\leq i\leq m-1$，
\[
\gamma_i=\left\{\bigl(\sqrt{r+a_1},\ldots,
e^{\mathrm i\theta}\sqrt{r+a_i},\ldots,
\sqrt{r+a_{m-1}},e^{-\mathrm i\theta}\sqrt r\bigr)
\ \middle|\ \theta\in\mathbb R\right\}.
\]

\SourceParagraph{ch07-7.2.5-033}{558}
它是映射 $\varphi_i\colon D=\{z\in\mathbb C\mid|z|\leq1\}\to\mathbb C^m$ 的像的边界，其中
\[
\varphi_i(z)=\bigl(\sqrt{r+a_1},\ldots,
z\sqrt{r+a_i},\ldots,\sqrt{r+a_{m-1}},\bar z\sqrt r\bigr).
\]
于是，
\[
\int_{\varphi_i(D)}\omega
=\int_D\varphi_i^*\omega
=\frac{\mathrm i}2\int_Da_i\,\mathrm dz\wedge\mathrm d\bar z
=\pi a_i.
\]
所以，$Y(M^a)=(\pi a_1,\ldots,\pi a_{m-2},0)$。

\SourceParagraph{ch07-7.2.5-034}{558}
由例题 7.17 中 $M^a$ 的显式描述可见，$Y(M^a)$ 位于 $H^1(M^a,\mathbb R)\to H^1(\Sigma,\mathbb R)$ 的像中；此像为 $(m-2)$ 维。
\end{example}

\SourceParagraph{ch07-7.2.5-035}{558}
现在，可以陈述 Joyce 光滑化结果中与本书讨论最相关的版本。
\BookStatementNumber{21}
\begin{theorem}\label{thm:7.21}
设 $X$ 为 Calabi--Yau $m$ 维流形，$2<m<6$；$M\subseteq X$ 为紧特殊 Lagrange 子流形，具有锥形奇点 $x_1,\ldots,x_n$，模型锥为 $C_1,\ldots,C_n$。设 $M_1,\ldots,M_n$ 为 $\mathbb C^m$ 中渐近锥形的特殊 Lagrange 子流形，模型锥分别为 $C_1,\ldots,C_n$，收敛率 $\lambda_1,\ldots,\lambda_n$ 均满足 $\lambda_i\leq0$。假定 $M^\circ=M\setminus\{x_1,\ldots,x_n\}$ 连通，且存在类 $\rho\in H^1(M^\circ,\mathbb R)$，在自然映射
\[
H^1(M^\circ,\mathbb R)\longrightarrow
\bigoplus_{i=1}^nH^1(\Sigma_i,\mathbb R)
\]
下的像为 $(Y(M_1),\ldots,Y(M_n))$。则存在 $\epsilon>0$，以及紧非奇异特殊 Lagrange 子流形的光滑族 $\{\widetilde M^t\mid t\in(0,\epsilon]\}$，使 $t\to0$ 时，$\widetilde M^t\to M$ 在流的意义下成立。
\end{theorem}

\SourceParagraph{ch07-7.2.5-036}{558}
基本思想如下。对任意 $t>0$，$tM_i$ 也渐近锥形，因此总有至少这一单参数变形族。$t\to0$ 时，$tM_i\to C_i$ 在流的意义下成立。对于小 $t$，用 $tM_i$ 替换 $M$ 的每个奇点 $x_i$ 的一个邻域；该邻域看起来像 $C_i$。这一粘合过程给出 $M$ 的一个光滑化，但它还不是特殊 Lagrange 的。不过，它已接近特殊 Lagrange；需要精细的分析，才能证明 $t$ 充分小时，可以将其变形为真正的特殊 Lagrange 子流形。然而，还存在一个由各 $Y(M_i)$ 给出的上同调阻碍。注意，自然映射 $H^1(M^\circ,\mathbb R)\to H^1(\Sigma_i,\mathbb R)$ 来自下述事实：若 $U$ 为 $x_i$ 在 $M$ 中的小邻域，则 $U\cap M^\circ=(0,r)\times\Sigma_i$，此映射就是限制映射。

\SourceParagraph{ch07-7.2.5-037}{559}
至此，对 Joyce 关于锥形特殊 Lagrange 子流形工作的介绍完成。下一小节，将应用这些结果，探讨复三维特殊 Lagrange 纤维化的真实形态。

\subsection{特殊 Lagrange \texorpdfstring{$T^3$}{T³} 纤维化}\label{sec:7.2.6}

\SourceParagraph{ch07-7.2.6-001}{559}
如前所述，Joyce 的例子说明，必须放弃特殊 Lagrange 纤维化一般都有余维二判别集的想法。因此，不能期待\BookChapterRef{6}中的拓扑对偶图景原样成立。这里介绍 Joyce 的推测：自然出现的特殊 Lagrange 纤维化可能有怎样的行为。当然，对于完整群恰为 $\operatorname{SU}(3)$ 的紧 Calabi--Yau 三维流形，仍然不知道任何一个特殊 Lagrange 纤维化的例子，因此也应记住，这些纤维化甚至可能不存在。不过，Joyce 的工作给出了许多局部例子，支持这里将描述的基本图景。以下思想均出自 Joyce。

\SourceParagraph{ch07-7.2.6-002}{559}
基本思想是，期待判别集为“变形虫”。Gelfand、Kapranov 和 Zelevinsky\BookCite{ref173}最初提出的变形虫概念，是 $(\mathbb C^*)^2$ 中全纯曲线在映射 $(z_1,z_2)\mapsto(\log|z_1|,\log|z_2|)$ 下的像。这些集合往往看起来像变形虫。

\SourceParagraph{ch07-7.2.6-003}{559}
考虑一个拓扑例子。设曲面 $S\subseteq(\mathbb C^*)^2$ 由 $z_1+z_2+1=0$ 定义，取映射
\[
\log\colon(\mathbb C^*)^2\longrightarrow\mathbb R^2,
\qquad\log(z_1,z_2)=(\log|z_1|,\log|z_2|).
\]
$\log(S)$ 的形状可以如下看出。先考虑 $S$ 在绝对值映射下的像，也就是先不取对数；此像如\BookRef{fig:7.7}{图7}所示。
\begin{figure}[H]
\centering
\BookImage{01_Numbered_Figures/Chapter_07/Ch07_Fig07_PDFp559_Bookp549.pdf}
\BookNumberOnlyCaption\label{fig:7.7}
\end{figure}

\SourceParagraph{ch07-7.2.6-004}{560}
线段 $r_1+r_2=1$、$r_1,r_2>0$，是 $\{(-a,a-1)\mid0<a<1\}\subseteq S$ 的像；射线 $r_2=r_1+1$、$r_1>0$，是 $\{(-a,a-1)\mid a<0\}$ 的像；射线 $r_1=r_2+1$ 则是 $\{(-a,a-1)\mid a>1\}$ 的像。映射 $S\to|S|$ 在 $|S|$ 的边界一一对应，在内部则二对一，$(z_1,z_2)$ 与 $(\bar z_1,\bar z_2)$ 映到同一点。对此图景取对数，得到 $S$ 的变形虫 $\log(S)$，如\BookRef{fig:7.8}{图8}所示。
\begin{figure}[H]
\centering
\BookImage{01_Numbered_Figures/Chapter_07/Ch07_Fig08_PDFp560_Bookp550.pdf}
\BookNumberOnlyCaption\label{fig:7.8}
\end{figure}

\SourceParagraph{ch07-7.2.6-005}{560}
现在，考虑 $\overline S=S\times\{0\}\subseteq\overline Y=(\mathbb C^*)^2\times\mathbb R=T^2\times\mathbb R^3$。利用\BookRef{constr:6.89}{构造6.89}，可以得到一个六维空间 $\overline X$ 及映射 $\pi\colon\overline X\to\overline Y$：它在 $\overline Y\setminus\overline S$ 上为 $S^1$ 丛，在 $\overline S$ 上退化，使 $\pi^{-1}(\overline S)\xrightarrow{\sim}\overline S$。将 $\pi$ 与
\[
(\log,\operatorname{id})\colon(\mathbb C^*)^2\times\mathbb R
\longrightarrow\mathbb R^3=B
\]
复合，就得到 $\overline X$ 上的 $T^3$ 纤维化 $f\colon\overline X\to\mathbb R^3$。显然，判别集为 $\log(S)\times\{0\}$。若 $b$ 位于其内部，则 $f^{-1}(b)$ 在拓扑上由 $T^3=T^2\times S^1$ 上的两个圆 $\{p_1\}\times S^1$ 和 $\{p_2\}\times S^1$ 收缩为点而得到。这正是特殊 Lagrange 情形中熟悉的锥形奇点。不过，目前还不知道如何将此纤维化变形为特殊 Lagrange 纤维化。

\SourceParagraph{ch07-7.2.6-006}{560}
若 $b\in\partial(\log(S)\times\{0\})$，则 $f^{-1}(b)$ 有一个稍复杂的奇点，但只有一个。这里不讨论这些纤维，而考察“一般”奇异纤维如何与上一小节的结果相符。特别地，对于 $\log(S)\times\{0\}$ 内部的 $b$，判别集在局部将 $B$ 分成两个区域，分别代表 $f^{-1}(b)$ 的两种可能光滑化。下面利用\BookRef{thm:7.21}{定理7.21}，说明这一行为如何可能发生。

\SourceParagraph{ch07-7.2.6-007}{561}
现在，假设 $f\colon X\to B$ 是拓扑如上述例子的特殊 Lagrange 纤维化，其判别集 $\Delta$ 为变形虫。取 $b\in\operatorname{Int}(\Delta)$，令 $M=f^{-1}(b)$，并设 $M^\circ=M\setminus\{x_1,x_2\}$，其中 $x_1,x_2$ 为两个锥形奇点。假定它们的切锥 $C_1,C_2$ 均为\BookRef{ex:7.17}{例题7.17}所考虑的 $M^0$ 型锥。则链环 $\Sigma_1,\Sigma_2$ 均为 $T^2$，期待它们在拓扑上可如下描述。注意，$M^\circ\cong(T^2\setminus\{y_1,y_2\})\times S^1$，$y_1,y_2$ 为 $T^2$ 中的两个点。假设 $\Sigma_i$ 具有形式 $\gamma_i\times S^1$，其中 $\gamma_i$ 为绕 $y_i$ 的简单环路。若这些假设成立，为用定理 7.21 判断 $M$ 的光滑化，考虑上同调限制映射
\[
H^1(M^\circ,\mathbb R)\longrightarrow
H^1(\Sigma_1,\mathbb R)\oplus H^1(\Sigma_2,\mathbb R).
\]

\SourceParagraph{ch07-7.2.6-008}{561}
不难看出，此映射的像为二维。事实上，取 $H^1(\Sigma_i,\mathbb R)$ 的基 $e_1^i,e_2^i$，其中 $e_1^i$ 为 $[\gamma_i]\times\mathrm{pt}$ 的 Poincaré 对偶，$e_2^i$ 为 $\mathrm{pt}\times S^1$ 的 Poincaré 对偶。限制映射的像由 $(e_1^1,e_1^2)$ 和 $(e_2^1,-e_2^2)$ 张成。上述具体拓扑纤维化模型并非特殊 Lagrange，因此，尤其不知道 $M$ 在 $x_1,x_2$ 的切锥在 $\mathbb C^3$ 中究竟如何嵌入，无法直接与渐近锥形光滑化比较。为给出合理性论证，选择 $H^1(\Sigma_i,\mathbb R)$ 的新基 $f_1^i,f_2^i$，使例题 7.17 给出的两处奇点切锥的三种光滑化 $M_i^{(a,0,0)}$、$M_i^{(0,a,0)}$、$M_i^{(0,0,a)}$ 满足
\[
\begin{aligned}
Y(M_i^{(a,0,0)})&=\pi af_1^i,\\
Y(M_i^{(0,a,0)})&=\pi af_2^i,\\
Y(M_i^{(0,0,a)})&=-\pi a(f_1^i+f_2^i).
\end{aligned}
\]
如\BookRef{ex:7.20}{例题7.20}那样计算，就可以验证这些确为正确的取值。

\SourceParagraph{ch07-7.2.6-009}{561}
假设在新基下，限制映射的像由 $(f_1^1,rf_2^2)$ 与 $(rf_2^1,f_1^2)$ 张成，其中 $r>0$、$r\neq1$。由\BookRef{thm:7.21}{定理7.21}，$M$ 有两种可能的光滑化方式：或者分别在 $x_1,x_2$ 粘入 $M_1^{(a,0,0)}$ 与 $M_2^{(0,ra,0)}$；或者分别粘入 $M_1^{(0,ra,0)}$ 与 $M_2^{(a,0,0)}$。这可能对应于将 $M$ 变形到 $f$ 判别集两侧的纤维。更多细节和其他可能性，见\BookCite{ref262}第 10.3 节。例如，加入 $r\neq1$，是为排除存在三种光滑化的情况；对于限制映射的像在基 $\{f_j^i\}$ 下的其他取值，也可能没有光滑化，或只有一种。

\SourceParagraph{ch07-7.2.6-010}{561}
这至少为拓扑类型如 $f$ 的特殊 Lagrange 纤维化的存在，提供了一种合理性论证。不过，当时还没有构造出这种纤维化。

\SourceParagraph{ch07-7.2.6-011}{561}
Joyce 最初的\BookRef{ex:7.11}{例题7.11}，给出余维一判别集、奇异纤维具有 $T^2$ 锥奇点的特殊 Lagrange 纤维化；加上上述对 $\operatorname U(1)$ 不变特殊 Lagrange 子流形及锥形奇点的更深入分析，迫使人们显著调整 Strominger--Yau--Zaslow 图景。不能只期待良好的余维二判别集，而必须面对特殊 Lagrange 纤维化中的余维一判别集。在这一图景下，不可避免地得出：某种“强形式”的 Strominger--Yau--Zaslow 猜想不能成立。特别地，若 $f\colon X\to B$ 与 $\check f\colon\check X\to B$ 为对偶特殊 Lagrange 纤维化，便不能期待它们的判别集相同。因此，不能指望逐纤维定义对偶化过程，必须细化纤维化对偶化的概念。\BookSectionRef{sec:7.3.6}{7.3.6}将再次讨论。这里先解释，为什么对偶化时判别集必须改变。

\SourceParagraph{ch07-7.2.6-012}{562}
关键在于\BookSectionRef{sec:6.4.2}{6.4.2}中正型与负型顶点的行为。从显式构造已经看到，正型情形的纤维化局部模型，其临界集为三条全纯曲线之并；负型情形的临界集则为裤形曲面。在“一般”的特殊 Lagrange 纤维化中，期待临界集大致保持原样，但它在底空间 $B$ 中的像变厚。负型情形的像，是\BookRef{fig:7.8}{图8}那样的变形虫。正型顶点的临界集，至少局部为三条全纯曲线之并，因此，期待判别集为三个不同变形虫之并。\BookRef{fig:7.9}{图9}显示了 Joyce 所预测的这两种新判别集。对偶化会交换正型和负型顶点，因此判别集必须改变。
\begin{figure}[H]
\centering
\BookImage{01_Numbered_Figures/Chapter_07/Ch07_Fig09_PDFp562_Bookp552.pdf}
\caption{负型与正型顶点的加厚。左图为负型顶点，判别集加厚为二维平面变形虫。右图中，判别集的三条支线各自在一个平面中加厚为变形虫状区域，这些平面沿一条直线横截相交。}\label{fig:7.9}
\end{figure}

\SourceParagraph{ch07-7.2.6-013}{563}
当然，这些讨论都带有很强的推测性，背后假设判别集只是\BookChapterRef{6}中图的加厚。不过，为涵盖这种可能性，显然需要新的对偶化概念。\BookSectionRef{sec:7.3.6}{7.3.6}将解释一个修订版本的 SYZ 猜想，处理这些问题。

\section{大复结构极限}\label{sec:7.3}
\subsection{紧流形上 Ricci 平坦度量的近似}\label{sec:7.3.1}

\SourceParagraph{ch07-7.3.1-001}{563}
\BookSectionRef{sec:7.1}{7.1}提到过，对于紧流形，还没有已知的非平坦 Ricci 平坦度量的显式例子。不过，在某些极限情形，可以写出非常精确的近似。为说明这一点，先考察 R. Kobayashi\BookCite{ref307}研究的 Kummer 曲面。

\SourceParagraph{ch07-7.3.1-002}{563}
\BookStatementNumber{22}
\begin{example}\label{ex:7.22}
设 $X$ 为复二维环面 $Y$ 的 Kummer 曲面，即 $\overline X=Y/\{\pm1\}$ 的极小消解。这里，$\{\pm1\}\cong\mathbb Z_2$ 在 $Y$ 上作用，$-1$ 为取负映射。$\overline X$ 有 16 个节点，$X$ 有 $(-2)$ 曲线 $E_1,\ldots,E_{16}$。考虑 $X$ 上形如
\[
\kappa-\sum_{i=1}^{16}a_ie_i
\]
的 Kähler 类，其中 $e_i$ 为 $E_i$ 的上同调类，$\kappa$ 为 $\overline X$ 上某个轨形 Kähler 类的拉回，各 $a_i$ 是小正数。Kobayashi 如下构造近似 Ricci 平坦度量。

\SourceParagraph{ch07-7.3.1-003}{563}
至少在局部解析意义下，每个节点都具有形式 $\mathbb C^2/\{\pm1\}$，即 $\mathbb C^3$ 中的二次锥；其消解为 $T^*\mathbb{CP}^1$，微分同胚于 $S^2$ 的实切丛。\BookRef{ex:7.1}{例题7.1}显式描述了这一消解上的 Eguchi--Hanson 度量：它是 $\operatorname{SO}(3)$ 不变的 Ricci 平坦 ALE Kähler 度量，按照那里所用的归一化，Kähler 类为 $-2a[E]$。此度量中，二球面 $E$ 的曲率为 $a^{-1}$，体积为 $4\pi a$。

\SourceParagraph{ch07-7.3.1-004}{563}
在 $E$ 补集中的任意紧集上，当 $a\to0$ 时，Eguchi--Hanson 度量收敛到轨形平坦度量。在原先的 K3 轨形 $\overline X$ 中，每个节点有一个解析同构于 $B/\{\pm1\}$ 的邻域，其中 $B$ 为 $\mathbb C^2$ 原点附近的球。取 $B$ 中一个固定圆环区域，将环外平坦度量与环内消解上的 Eguchi--Hanson 度量的 Kähler 势粘合；对第 $i$ 个节点，使用例题 7.1 中参数为 $a_i$ 的度量。这样，就在 Kummer 曲面 $X$ 上，得到属于所需 Kähler 类的度量 $\omega_a$。$a=(a_1,\ldots,a_{16})$ 很小时，它几乎 Ricci 平坦，曲率只在例外曲线附近变大。

\SourceParagraph{ch07-7.3.1-005}{563}
现在实行 Yau 的求解程序，得到适当归一化的函数 $u_a$，以及 Ricci 平坦度量
\[
\widetilde\omega_a=\omega_a+\mathrm i\partial\bar\partial u_a.
\]
也就是说，若 $\Omega$ 为全纯二形式，则 $(\tfrac{\mathrm i}2)^2\Omega\wedge\overline\Omega$ 为体积形式 $\widetilde\omega_a^2/2$ 的一个常数倍。

\SourceParagraph{ch07-7.3.1-006}{564}
此外，\BookCite{ref307}对 Yau 证明的分析，给出 $u_a$ 各种范数的显式估计；参见下文\BookSectionRef{sec:7.3.5}{7.3.5}。
\begin{enumerate}[label=\textup{(\Roman*)}]
\item $C^0$ 估计：$\|u_a\|_{C^0}\leq C|a|^2$。
\item $C^2$ 估计。假设 $a\to0$ 时，所有 $a_i$ 之间的比值保持有界，最方便的写法为
\[
(1-C|a|^{1/2})\omega_a\leq\widetilde\omega_a
\leq(1+C|a|^{1/2})\omega_a.
\]
适当缩放后，在每条 $(-2)$ 曲线的邻域中，还有 $i=2,3,4$ 时的 $C^i$ 估计。
\item 在 $(-2)$ 曲线补集中的任意固定相对紧子区域，有 $C^{k,\alpha}$ 估计
\[
\|u_a\|_{C^{k,\alpha}}\leq C|a|^2.
\]
$C^{k,\alpha}$ 的定义见\BookSectionRef{sec:7.1.1}{7.1.1}。这里，先应用一般结果\BookCite{ref179}定理 17.14，从 $C^2$ 估计得到 $C^{2,\alpha}$ 估计，再用正则性提升与 Schauder 估计的标准论证，得到更强的结论。
\end{enumerate}
\end{example}

\SourceParagraph{ch07-7.3.1-007}{564}
\BookStatementNumber{23}
\begin{note}[numbered]\label{note:7.23}
对于任意有理双点奇点，可以粘入 Kronheimer 的 ALE 引力瞬子\BookCite{ref316}，代替 Eguchi--Hanson 度量，同样的论证仍成立\BookCite{ref307}。
\end{note}

\SourceParagraph{ch07-7.3.1-008}{564}
对于任意 K3 曲面向具有有理双点曲面的退化，Kobayashi 重复上述论证，以轨形 Ricci 平坦 Kähler 度量为背景度量；其存在性来自 Yau 定理。由于此度量局部仅在三阶近似意义下与 Euclid 度量一致，所以粘合不再使用固定圆环区域，而使用距奇点约 $a^{3/8}$ 的圆环区域。

\SourceParagraph{ch07-7.3.1-009}{564}
这一更一般情形的估计稍差。例如，第（I）与第（III）项中，$|a|^2$ 要换为 $|a|^{13/8}$。注意，与 Kummer 情形不同，这里仅在 $(-2)$ 曲线附近，才有 Calabi--Yau 度量的显式近似。

\SourceParagraph{ch07-7.3.1-010}{564}
当 K3 曲面的 Kähler 结构退化到具有有理双点的结构时，上述近似 Ricci 平坦度量的描述，使我们能够推出：Ricci 平坦度量的 Gromov--Hausdorff 极限，见下文\BookSectionRef{sec:7.3.6}{7.3.6}，恰为轨形 Ricci 平坦度量。Anderson\BookCite{ref013}一般地研究了 K3 曲面上体积不塌缩、直径有界的 Ricci 平坦度量的极限，证明可以取子列收敛到轨形度量。对于体积统一有界而直径无界的情形，Anderson 证明必发生塌缩：若将度量缩放到直径固定，则 Gromov--Hausdorff 极限的维数严格降低。对于这里讨论的 K3 曲面 Kähler 退化，主要有两种情形：上述轨形退化，以及将在\BookSectionRef{sec:7.3.3}{7.3.3}至\BookSectionRef{sec:7.3.6}{7.3.6}研究的椭圆纤维化情形。后一情形将通过超 Kähler 旋转，与 K3 曲面的大复结构极限联系起来。

\subsection{K3 曲面的大复结构极限}\label{sec:7.3.2}

\SourceParagraph{ch07-7.3.2-001}{565}
Calabi--Yau 流形的大复结构极限点的定义，包括如下条件：它们是复模空间边界上，具有极大幺幂单值变换的点。文献\BookCite{ref365}对 Calabi--Yau 三维流形的定义，还包括一个整性条件。本小节只研究 K3 曲面，以周期来表述，就能得到这类极限点的具体描述，以及镜像对称的显式公式。这里的描述最初来自 Pinkham、Nikulin 和 Dolgachev 的工作，细节见\BookCite{ref126}。

\SourceParagraph{ch07-7.3.2-002}{565}
用 $L$ 表示 K3 格，考虑带有同构 $H^2(X,\mathbb Z)\cong L$ 的 K3 曲面 $X$ 的模空间。这样的带标记 K3 曲面的大复结构极限点，是具有极大幺幂单值变换的边界点；它给出本原各向同性向量 $E\in L$，即 $E^2=0$，生成 $H^2(X,\mathbb Q)$ 的单值变换权滤过的 $W_0$ 部分，也可比较\BookRef{thm:6.104}{定理6.104}。从类 $E$ 出发，将构造三元组 $(X,B+\mathrm i\omega,\Omega)$ 模空间上的对合，其中 $X$ 为带标记 K3 曲面，$\Omega$ 为全纯二形式的类，$\omega\in E^\perp\otimes\mathbb R$ 为 $X$ 的 Kähler 类，$B$ 场属于 $(E^\perp/\langle E\rangle)\otimes\mathbb R$。此外，将 $\Omega$ 归一化为
\[
\operatorname{Im}\Omega\in E^\perp\otimes\mathbb R,
\qquad\omega^2=(\operatorname{Re}\Omega)^2=(\operatorname{Im}\Omega)^2,
\]
并选择 $E$ 的符号，使 $E\cdot\operatorname{Re}\Omega>0$。

\SourceParagraph{ch07-7.3.2-003}{565}
考虑复维数二十的周期域
\[
\begin{aligned}
\mathcal D&=\{[\Omega]\in\mathbb P(L_{\mathbb C})
\mid\Omega^2=0,\ \Omega\cdot\overline\Omega>0\}\\
&\cong\operatorname O(3,19)/
\bigl(\operatorname{SO}(2)\times\operatorname O(1,19)\bigr).
\end{aligned}
\]
投影映射
\[
q\colon\mathbb P(L_{\mathbb C})\setminus\{[E]\}
\longrightarrow\mathbb P(L_{\mathbb C}/\langle E\rangle)
\]
限制到 $\mathcal D$ 后，将子集 $\{[\Omega]\in\mathcal D\mid\Omega\cdot E\neq0\}$，识别为仿射空间
\[
\{v\in L_{\mathbb C}/\langle E\rangle\mid v\cdot E=1\}
\]
的一个开子集。

\SourceParagraph{ch07-7.3.2-004}{565}
将 $\Omega$ 归一化为 $\operatorname{Im}\Omega\cdot E=0$，并选择 $\sigma_0\in L$，满足 $\sigma_0^2=-2$、$\sigma_0\cdot E=1$，使 $\langle E,\sigma_0\rangle$ 为双曲格。从 $q([\Omega])$ 的上述归一化代表减去 $\sigma_0$，就将其识别为 $E^\perp/\langle E\rangle$ 的一个点。模去 $E$ 的倍数，得到 $L_{\mathbb C}$ 中的点 $\check B+\mathrm i\check\omega$。显式地，
\[
\check B\equiv(E\cdot\operatorname{Re}\Omega)^{-1}
\operatorname{Re}\Omega-\sigma_0\pmod E,
\]
且
\begin{equation}\label{eq:7.6}
\check\omega\equiv(E\cdot\operatorname{Re}\Omega)^{-1}
\operatorname{Im}\Omega\pmod E.
\end{equation}

\SourceParagraph{ch07-7.3.2-005}{566}
所需对合交换 $(X,B+\mathrm i\omega,\Omega)$ 与 $(\check X,\check B+\mathrm i\check\omega,\check\Omega)$，其中
\begin{equation}\label{eq:7.7}
\check\Omega\equiv(E\cdot\operatorname{Re}\Omega)^{-1}
(\sigma_0+B+\mathrm i\omega)\pmod E,
\end{equation}
而 $\check X$ 表示周期为 $\check\Omega$ 的带标记 K3 曲面。注意，$\check\Omega$ 与 $\check\omega$ 的实际类，完全由关系
\[
(\operatorname{Re}\check\Omega)^2
=(\operatorname{Im}\check\Omega)^2=\check\omega^2,
\]
以及
\[
\check\omega\cdot\operatorname{Re}\check\Omega
=\check\omega\cdot\operatorname{Im}\check\Omega
=\operatorname{Re}\check\Omega\cdot\operatorname{Im}\check\Omega=0
\]
确定。

\SourceParagraph{ch07-7.3.2-006}{566}
将论证，这一公式与镜像对称的 SYZ 图景一致。关键在于，利用所谓超 Kähler 技巧，很容易在 K3 曲面上构造特殊 Lagrange 环面纤维化。具有 Ricci 平坦，从而为超 Kähler 度量的 K3 曲面，有一个 $S^2$ 参数族的相容复结构。超 Kähler 流形有三个与度量相容的复结构 $I,J,K$，满足 $I^2=J^2=K^2=-1$、$IJ=K$、$JK=I$、$KI=J$。一般的相容复结构为 $aI+bJ+cK$，其中 $a^2+b^2+c^2=1$。以形式表述：若复结构 $I$ 的 Kähler 形式为 $\omega$，全纯二形式为 $\Omega$，且归一化为 $\omega^2=(\operatorname{Re}\Omega)^2=(\operatorname{Im}\Omega)^2$，则适当标记旋转后的 $K$ 复结构，其全纯二形式与 Kähler 形式为
\[
\Omega_K=\operatorname{Im}\Omega+\mathrm i\omega,
\qquad\omega_K=\operatorname{Re}\Omega.
\]
总可以先将 $\Omega$ 乘以任意相位因子 $e^{\mathrm i\theta}$，从而使 $\operatorname{Im}\Omega\cdot E=0$。于是，$\Omega_K\cdot E=0$，结合 $E\cdot\omega_K>0$，蕴含 $E$ 为一个有效除子的类，且 $E^2=0$。对于一般的数据选择，$E$ 为椭圆纤维化 $f\colon X_K\to\mathbb{CP}^1$ 的纤维类。由于 $\Omega_K$ 在该映射的纤维上限制为零，$f\colon X_I\to\mathbb{CP}^1$ 的一般纤维便为特殊 Lagrange 环面。

\SourceParagraph{ch07-7.3.2-007}{566}
再取 $f$ 的 $C^\infty$ 截面 $\sigma_0$，对应于 $X_K\to\mathbb{CP}^1$ 的 Jacobian 纤维化的零截面。具体地，\BookCite{ref210}第 482 页说明，由
\[
\Omega'=\Omega_K-f^*\sigma_0^*\Omega_K
\]
确定的复结构，恰为 Jacobian 纤维化的复结构。回顾，椭圆纤维化的 Jacobian 纤维化，具有与原纤维同构的纤维，并且有一个截面。此截面的类也记为 $\sigma_0$。

\SourceParagraph{ch07-7.3.2-008}{566}
假定 $f$ 的所有纤维不可约，则将镜像对称解释为对偶化特殊 Lagrange 纤维化 $f\colon X\to\mathbb{CP}^1$，与上述由 $E,\sigma_0$ 确定的显式公式一致。文献\BookCite{ref209}第 4 节，在 $\sigma_0$ 为 $f$ 的特殊 Lagrange 截面的假设下验证这一结论；换言之，$X_K\to\mathbb{CP}^1$ 为 Jacobian 纤维化，具有全纯截面。一般情形则由\BookCite{ref210}第 1 节的讨论得到。若 $f\colon X\to B$ 为上述特殊 Lagrange 纤维化，$B=S^2$，则 Leray 谱序列
\[
H^i(S^2,R^jf_*\mathbb Q)\Longrightarrow H^{i+j}(X,\mathbb Q)
\]
在 $E_2$ 页退化，非零项如下图所示，见\BookCite{ref209}命题 4.1。
\begin{figure}[H]
\centering
\BookImage{02_Unnumbered_Diagrams/Chapter_07/Ch07_Diagram_PDFp567_01_Bookp557.pdf}
\end{figure}

\SourceParagraph{ch07-7.3.2-009}{567}
由此给出滤过
\[
0\subseteq F_0\subseteq F_1\subseteq F_2=H^2(X,\mathbb Q),
\]
其中 $F_1=E^\perp$，且
\[
\begin{aligned}
F_0&\cong H^2(S^2,f_*\mathbb Q)\cong\mathbb QE,\\
F_1/F_0&\cong H^1(S^2,R^1f_*\mathbb Q),\\
F_2/F_1&\cong H^0(S^2,R^2f_*\mathbb Q)\cong\mathbb Q\sigma_0.
\end{aligned}
\]
这里最后的 $\sigma_0$ 表示其在商空间中的类。

\SourceParagraph{ch07-7.3.2-010}{567}
适当重新编号后，这一滤过实际上与大复结构极限点决定的权滤过一致。\BookSectionRef{sec:6.2.3}{6.2.3}的镜像对称描述提示，在上同调层面，镜像对合应交换 $\Omega$ 与 $\exp(B+\mathrm i\omega)$。在上述滤过上，可以如下看出这种交换。先通过 Poincaré 对偶，将 $X\to S^2$ 的对偶纤维化与原纤维化识别。若进一步归一化 $\Omega$，使 $\operatorname{Re}\Omega\cdot E=1$，即 $\Omega$ 在纤维上的限制为归一化体积形式，则
\[
\operatorname{Im}\Omega\in F_1\otimes_{\mathbb Q}\mathbb R,
\qquad\operatorname{Re}\Omega\in\sigma_0+F_1\otimes_{\mathbb Q}\mathbb R.
\]
因此，$\Omega-\sigma_0\in F_1\otimes_{\mathbb Q}\mathbb C$，其在 $(F_1/F_0)\otimes_{\mathbb Q}\mathbb C$ 中的像，应与 $\check B+\mathrm i\check\omega$ 一致。反过来，可以将 $B+\mathrm i\omega$ 看作 $(F_1/F_0)\otimes_{\mathbb Q}\mathbb C$ 的元素，并提升为 $F_1\otimes_{\mathbb Q}\mathbb C$ 中唯一的类 $\alpha$，使 $\alpha+\sigma_0=\check\Omega$ 满足 $\check\Omega^2=0$。这正是\BookEqRef{eq:7.6}{7.6}、\BookEqRef{eq:7.7}{7.7}的镜像对称描述。在形式层面对这一对偶化过程的详细解释，见\BookCite{ref201}。

\SourceParagraph{ch07-7.3.2-011}{567}
有了镜像映射的上述显式公式，现在直接将 $\check X$ 的大复结构极限，定义为 $X$ 的大 Kähler 结构极限的镜像。在后者中，保持 $X$ 的复结构固定，让 Kähler 形式趋于无穷。更准确地说，若 $\{B_l+\mathrm i t_l\omega\}$ 为 $X$ 上复化 Kähler 形式的序列，$t_l>0$、$t_l\to\infty$，则称此序列趋近于 $X$ 的复化 Kähler 模空间中的大 Kähler 极限。

\SourceParagraph{ch07-7.3.2-012}{568}
\BookStatementNumber{24}
\begin{definition}\label{def:7.24}
对于每个 $l$，设 $\check X_l$ 为数据 $(\check X_l,\check B_l+\mathrm i\check\omega_l,\check\Omega_l)$ 给出的 K3 曲面，这组数据与 $(X,B_l+\mathrm i t_l\omega,t_l\Omega)$ 互为镜像。称曲面序列 $\{\check X_l\}$ 趋近于一个大复结构极限点。
\end{definition}

\SourceParagraph{ch07-7.3.2-013}{568}
实际上，这里稍作了简化，仅沿一条射线趋近大 Kähler 极限。更一般的方法，可以容许更一般的 Kähler 形式序列；不过，那样更难处理，因为下文出现的 Jacobian 纤维化会随之变化，掩盖主要目标。采用上述定义，就有用 $t_l,\Omega,\omega,B_l$ 表示 $\check\Omega_l$、$\check\omega_l$ 的公式。

\SourceParagraph{ch07-7.3.2-014}{568}
Kähler 类 $\check\omega_l$ 由 Ricci 平坦度量 $\check g_l$ 代表，希望理解 $t_l\to\infty$ 时这一度量的行为。方便的做法，是实施超 Kähler 旋转：$\check g_l$ 也为 K3 曲面 $\check X_{l,K}$ 上的 Kähler 度量，其形式为
\[
\check\Omega_{l,K}=\operatorname{Im}\check\Omega_l+\mathrm i\check\omega_l,
\qquad\check\omega_{l,K}=\operatorname{Re}\check\Omega_l.
\]
这些等式在微分形式层面成立。具体地，在上同调中，\BookCite{ref210}第 481 页给出
\[
\begin{aligned}
\check\Omega_{l,K}
&=(E\cdot\operatorname{Re}\Omega)^{-1}
\Bigl(\omega+\mathrm i\operatorname{Im}\Omega\\
&\hspace{25mm}-\bigl((\omega+\mathrm i\operatorname{Im}\Omega)
\cdot(\sigma_0+B_l)\bigr)E\Bigr),\\
\check\omega_{l,K}
&\equiv(t_lE\cdot\operatorname{Re}\Omega)^{-1}(\sigma_0+B_l)\pmod E.
\end{aligned}
\]
假定数据选择一般，$E$ 代表椭圆纤维化 $f_l\colon\check X_{l,K}\to\mathbb{CP}^1$ 的纤维类；这一椭圆纤维化，与 $\check X_l$ 上的特殊 Lagrange $T^2$ 纤维化一致。注意，相对于 $\check g_l$，$f_l$ 的纤维面积为 $(t_lE\cdot\operatorname{Re}\Omega)^{-1}$，随 $t_l\to\infty$ 而趋于零。

\SourceParagraph{ch07-7.3.2-015}{568}
$\check\Omega_{l,K}$ 虽依赖于 $l$，但这些类之间只相差从 $\mathbb{CP}^1$ 拉回的类。实际上，这说明椭圆 K3 曲面 $\check X_{l,K}$ 密切相关；特别地，它们具有同一个 Jacobian $\check J_K$。这是具有全纯截面的唯一 Jacobian 椭圆 K3 曲面，其复结构由某个 $\mathbb{CP}^1$ 上二形式 $\alpha$ 对应的 $\check\Omega_{l,K}+f_l^*\alpha$ 诱导。

\SourceParagraph{ch07-7.3.2-016}{568}
由此得到下面的问题。
\BookStatementNumber{25}
\begin{problem}\label{prob:7.25}
设 $j\colon J\to\mathbb{CP}^1$ 是具有截面的椭圆 K3 曲面，$f_l\colon X_l\to\mathbb{CP}^1$ 是 Jacobian 均为 $j\colon J\to\mathbb{CP}^1$ 的椭圆 K3 曲面序列。设 $\omega_l$ 为 $X_l$ 上的 Ricci 平坦 Kähler 度量，$\operatorname{Vol}(X_l)$ 与 $l$ 无关。令
\[
\epsilon_l=\operatorname{Area}_{\omega_l}(f_l^{-1}(y)),
\qquad y\in\mathbb{CP}^1,
\]
并假设 $l\to\infty$ 时，$\epsilon_l\to0$。描述 $\omega_l$ 在 $l\to\infty$ 时的行为。
\end{problem}

\SourceParagraph{ch07-7.3.2-017}{568}
文献\BookCite{ref210}在 $j$ 有 24 条 $I_1$ 型 Kodaira 纤维的情形，解决了这一问题；一般的 K3 椭圆纤维化满足此条件。从现在起，对奇异纤维作这一假设。下面将描述如何构造度量的非常精确的近似。为简单起见，只处理所有椭圆 K3 曲面 $f_l\colon X_l\to\mathbb{CP}^1$ 都与 Jacobian 纤维化一致的情形；在上述表述中，这对应于 $B_l$ 取特殊值。一般情形也类似，只需对论证略作修改\BookCite{ref210}。

\subsection{远离奇异纤维的度量}\label{sec:7.3.3}

\SourceParagraph{ch07-7.3.3-001}{569}
从现在起，设 $f\colon X\to\mathbb{CP}^1$ 为固定的椭圆 K3 曲面，有全纯截面 $\sigma_0$，以及 24 条 $I_1$ 型奇异纤维。令 $\Omega$ 为 $X$ 上的全纯二形式，$\omega$ 为 Kähler 形式，满足 $[\omega]^2=[\operatorname{Re}\Omega]^2=[\operatorname{Im}\Omega]^2$。设 $\epsilon=\int_{X_b}\omega$ 为纤维面积。若 $g$ 为 $X$ 上相应的 Ricci 平坦度量，希望理解 $\epsilon$ 变小时 $g$ 的行为。结果表明，奇异纤维附近的行为，与其他地方非常不同。本小节描述远离奇异纤维时 $g$ 的良好近似，即所谓半平坦度量。为显式描述它们，需要更自然地描述纤维化，并选择一套正则坐标。

\SourceParagraph{ch07-7.3.3-002}{569}
用 $T^*_{\mathbb{CP}^1}$ 表示底空间的全纯余切丛。给定 $\mathbb{CP}^1$ 的开集 $U$ 与光滑截面 $\alpha\in\Gamma(U,T^*_{\mathbb{CP}^1})$，有自然作用 $F_\alpha\colon f^{-1}(U)\to f^{-1}(U)$，定义如下。先设 $\alpha$ 在 $U$ 中紧支撑，则一形式 $f^*\alpha$ 在 $f^{-1}(U)$ 中紧支撑。利用 $X$ 的全纯辛形式 $\Omega$，得到唯一的实向量场 $V$，切于 $f$ 的纤维，满足 $\iota(V)\Omega=f^*\alpha$；积分后，得到对所有实 $t$ 定义的流 $\varphi_t\colon f^{-1}(U)\to f^{-1}(U)$。取 $F_\alpha=\varphi_1$。若进一步假设 $\alpha$ 闭合，容易验证，$F_\alpha$ 保持 $\Omega$，并且逐纤维作用。现在观察，$V|_{f^{-1}(b)}$ 仅依赖于 $\alpha$ 在 $b\in U$ 处的值，因此，对于任意光滑截面 $\alpha$，仍可以定义逐纤维作用 $F_\alpha$。

\SourceParagraph{ch07-7.3.3-003}{569}
由此定义映射 $\pi\colon T^*_{\mathbb{CP}^1}\to X$，将 $(P,\alpha)$ 映为 $F_\alpha(\sigma_0(P))$；余切丛的零截面对应于给定全纯截面 $\sigma_0$。此映射的分析见\BookCite{ref201}第 2 节及第 7 节。在当前情形，$\pi$ 的像恰为 $f$ 临界集在 $X$ 中的补集 $X^0$；其核为余切丛纤维中全纯变化、在奇异纤维处退化的格族。可以将此核识别为 $R^1f_*\mathbb Z\subset T^*_{\mathbb{CP}^1}$，这里采用对椭圆 K3 曲面成立的自然识别 $R^1f_*\mathcal O_X\cong\omega_{\mathbb{CP}^1}$，其中 $\omega_{\mathbb{CP}^1}$ 为底空间的典范层，也就是 $T^*_{\mathbb{CP}^1}$ 的全纯截面层。因此，$X^0$ 识别为全纯余切丛对 $\Lambda\cong R^1f_*\mathbb Z$ 的商。

\SourceParagraph{ch07-7.3.3-004}{569}
在 $\mathbb{CP}^1$ 上局部取余切丛的全纯正则坐标 $(x,y)$，其中 $y$ 为底空间的局部全纯坐标，$(x,y)$ 对应于余切向量 $x\,\mathrm dy$；零截面局部由 $x=0$ 给出。全纯余切丛有典范全纯辛形式，局部为 $\mathrm d(x\,\mathrm dy)=\mathrm dx\wedge\mathrm dy$。$\pi^*\Omega$ 也是余切丛上处处非零的全纯二形式；按照上述流的归一化，它恰与典范全纯辛形式相等，改变整体归一化时至多相差一个常数。若 $f\colon X\to\mathbb{CP}^1$ 不是 Jacobian 纤维化，仍有类似理论，但需要对上述描述略作修改，涉及从底空间拉回的二形式，见\BookCite{ref201}第 7 节。

\SourceParagraph{ch07-7.3.3-005}{570}
回到具有给定全纯截面 $\sigma_0$ 的椭圆 K3 曲面。在远离奇异纤维处，可以将纤维化描述为 $f_0\colon X_0\to B_0$，其中 $B_0$ 在坐标片中为 $\mathbb C$ 的开子集，$X_0=T^*_{B_0}/\Lambda$，$\Lambda$ 为全纯变化的格族。取 $X_0$ 的局部全纯坐标 $(x,y)$，对应于 $T^*_{B_0}$ 的全纯正则坐标，$y=y_1+\mathrm i y_2$、$x=x_1+\mathrm i x_2$。$X_0$ 的全纯二形式由 $T^*_{B_0}$ 上的 $\Omega=\mathrm dx\wedge\mathrm dy$ 诱导。对于原纤维化 $f\colon X\to\mathbb{CP}^1$，在避开 24 个点构成的判别集 $\Delta$ 的开集 $B_0$ 上，就使用这些坐标。现在，将 Kähler 形式局部写为
\[
\begin{aligned}
\omega=\frac{\mathrm i}2W\Bigl(&\mathrm dx\wedge\mathrm d\bar x
+\bar b\,\mathrm dx\wedge\mathrm d\bar y
+b\,\mathrm dy\wedge\mathrm d\bar x\\
&+(W^{-2}+|b|^2)\,\mathrm dy\wedge\mathrm d\bar y\Bigr)\\
=\frac{\mathrm i}2\Bigl(&W(\mathrm dx+b\,\mathrm dy)
\wedge(\mathrm d\bar x+\bar b\,\mathrm d\bar y)
+W^{-1}\mathrm dy\wedge\mathrm d\bar y\Bigr).
\end{aligned}
\]

\SourceParagraph{ch07-7.3.3-006}{570}
这里，函数 $W,b$ 由这一表达式定义，$\mathrm dy\wedge\mathrm d\bar y$ 的系数保证归一化 $\omega^2=(\operatorname{Im}\Omega)^2$。$W$ 为正实值函数，$b$ 为复值函数。\BookRef{prop:6.25}{命题6.25}用形式刻画 Ricci 平坦度量，从而保证，只要上述形式为 Kähler 的，相应度量就 Ricci 平坦。Kähler 条件为 $\mathrm d\omega=0$，可以写成
\[
\partial_yW=\partial_x(Wb),\qquad
\partial_y(W\bar b)=\partial_x\bigl(W(W^{-2}+|b|^2)\bigr).
\]
展开第二式，得到
\[
\begin{aligned}
W\partial_y\bar b+\bar b\partial_yW
={}&-W^{-2}\partial_xW+(\partial_xW)|b|^2\\
&+W\bigl(b\partial_x\bar b+\bar b\partial_xb\bigr).
\end{aligned}
\]
用第一式替换 $\partial_yW$ 并化简，可知上述两个方程等价于
\begin{equation*}\tag{K1}\label{eq:7.K1}
(\partial_y-b\partial_x)\bar b=-W^{-3}\partial_xW,
\end{equation*}
\begin{equation*}\tag{K2}\label{eq:7.K2}
(\partial_y-b\partial_x)W=W\partial_xb.
\end{equation*}

\SourceParagraph{ch07-7.3.3-007}{570}
在这一框架中，显然应研究\BookSectionRef{sec:6.2.1}{6.2.1}所述的半平坦度量。按照上述描述，半平坦意味着 $W$ 在每条纤维上为常数。设 $\tau_1(y),\tau_2(y)$ 为 $B_0$ 上两个局部全纯函数，使 $\tau_1(y)\,\mathrm dy$、$\tau_2(y)\,\mathrm dy$ 在每个 $y\in B_0$ 生成格 $\Lambda(y)\subseteq T^*_{B_0,y}$，从而给出全纯变化的格族 $\Lambda\subseteq T^*_{B_0}=B_0\times\mathbb C$。通常可以容许 $\tau_1,\tau_2$ 多值。不妨假设 $\operatorname{Im}(\bar\tau_1\tau_2)>0$，则 $X_0=(B_0\times\mathbb C)/\Lambda$ 上的 Ricci 平坦度量由数据
\[
\begin{aligned}
W&=\frac\epsilon{\operatorname{Im}(\bar\tau_1\tau_2)},\\
b&=-\frac W\epsilon\bigl[
\operatorname{Im}(\tau_2\bar x)\partial_y\tau_1
+\operatorname{Im}(\bar\tau_1x)\partial_y\tau_2\bigr]
\end{aligned}
\]
给出。容易验证，它们满足\BookEqRef{eq:7.K1}{K1}与\BookEqRef{eq:7.K2}{K2}。此度量起初定义在 $T^*_{B_0}$ 上，随后下降到 $X_0$；$f$ 的每条纤维面积为 $\epsilon$。称它为标准半平坦度量，Kähler 形式记为 $\omega_{\mathrm{SF}}$。文献\BookCite{ref196}首先考虑了这一度量。

\SourceParagraph{ch07-7.3.3-008}{571}
可以显式验证，此度量与 $\Lambda$ 的生成元选择无关，所以 $\tau_1,\tau_2$ 的多值性不会造成问题。此外，它与坐标 $y$ 的选择无关；注意，改变 $y$ 时，正则坐标 $x$ 及函数 $\tau_1,\tau_2$ 也必须相应改变。这还可以如下理解。由包含关系 $R^1f_{0*}\mathbb Z\cong\Lambda\subseteq T^*_{B_0}$，可以将 $(R^1f_{0*}\mathbb R)\otimes\mathcal C^\infty_{B_0}$ 识别为 $T^*_{B_0}$ 的底层光滑实向量丛，并赋予其 Gauss--Manin 联络 $\nabla_{\mathrm{GM}}$；这是平坦联络，平坦截面为 $R^1f_{0*}\mathbb R$ 的截面。标准半平坦度量，是满足下列条件的唯一半平坦 Ricci 平坦 Kähler 度量：
\begin{enumerate}
\item 每条纤维的面积为 $\epsilon$。
\item $\omega_{\mathrm{SF}}^2=(\operatorname{Re}\Omega)^2=(\operatorname{Im}\Omega)^2$。
\item 每点垂直切空间的正交补，恰为 $\nabla_{\mathrm{GM}}$ 在该点给出的水平切空间。
\end{enumerate}

\SourceParagraph{ch07-7.3.3-009}{571}
不过，应注意，若 $T_\sigma\colon X_0\to X_0$ 表示沿全纯截面 $\sigma$ 的平移，则 $T_\sigma^*\omega_{\mathrm{SF}}$ 可能给出另一半平坦度量，满足第（1）、第（2）项，却不满足第（3）项。但若 $\sigma$ 除全纯外，还是 Gauss--Manin 联络的平坦截面，即 $\sigma(y)=a_1\tau_1(y)+a_2\tau_2(y)$，其中 $a_1,a_2$ 为实常数，则 $T_\sigma$ 为等距变换，且 $T_\sigma^*\omega_{\mathrm{SF}}=\omega_{\mathrm{SF}}$、$T_\sigma^*\Omega=\Omega$。

\SourceParagraph{ch07-7.3.3-010}{571}
度量的 Kähler 势也很有用。它是满足 $\omega=\tfrac{\mathrm i}{2}\partial\bar\partial\varphi$ 的函数 $\varphi$。设 $\phi_1,\phi_2$ 分别为 $\tau_1,\tau_2$ 的原函数，则可以取
\[
\begin{aligned}
\varphi={}&\frac\epsilon{\operatorname{Im}(\bar\tau_1\tau_2)}
\left(-\frac{\bar x^2\tau_1}{2\bar\tau_1}
+|x|^2-\frac{x^2\bar\tau_1}{2\tau_1}\right)\\
&+\frac{\mathrm i}{2\epsilon}
(\phi_1\bar\phi_2-\bar\phi_1\phi_2).
\end{aligned}
\]
它在单连通 $U$ 上的子集 $T^*_U\subseteq T^*_{B_0}$ 中定义良好，但不下降到 $T^*_{B_0}/\Lambda$。

\SourceParagraph{ch07-7.3.3-011}{571}
综上，上述半平坦度量 Ricci 平坦，纤维面积为 $\epsilon$。在底空间远离奇异纤维的任意固定紧区域中，纤维直径为 $O(\epsilon^{1/2})$，曲率张量 $R$ 的 $C^0$ 范数可以计算为 $\|R\|=O(\epsilon)$；下文会说明这一计算。更多细节见\BookCite{ref210}。

\subsection{奇异纤维邻域中的度量}\label{sec:7.3.4}

\SourceParagraph{ch07-7.3.4-001}{572}
上一小节看到，远离奇异纤维时，适于计算的是复作用角坐标。奇异纤维附近，这些坐标不再合适；回顾，已假设奇异纤维为 $I_1$ 型。此时，合适的坐标来自\BookSectionRef{sec:7.1.2}{7.1.2}介绍的 Gibbons--Hawking 拟设。

\SourceParagraph{ch07-7.3.4-002}{572}
本小节的目标，是在椭圆纤维化 K3 曲面的每条奇异纤维的邻域，描述一个超 Kähler 度量，并推导相关估计。假定纤维体积为 $\epsilon$，则 $\epsilon$ 很小时，远离奇异纤维，这一度量迅速接近半平坦度量。由于奇异纤维为 $I_1$ 型 Kodaira 纤维，其附近有一个周期在单值变换下不变；事实上，适当选择底空间的全纯坐标 $y$，可令此周期为常数一。另一个周期则多值，并趋向无穷。将定义纤维化光滑部分上的 $S^1$ 不变度量，如\BookSectionRef{sec:7.1.3}{7.1.3}所述；最方便的描述方式是 Gibbons--Hawking 坐标。Ooguri 和 Vafa\BookCite{ref383}首次以稍有不同的形式写出这一度量，另见\BookCite{ref416}，因此称为 Ooguri--Vafa 度量。

\SourceParagraph{ch07-7.3.4-003}{572}
设 $D_r\subset\mathbb C$ 为以原点为中心、半径 $r<1$ 的圆盘，$f\colon\overline X\to D_r$ 为椭圆纤维化，原点上方有 $I_1$ 型奇异纤维。令
\[
\overline Y=D_r\times(\mathbb R/\epsilon\mathbb Z),\qquad
Y=\bigl((D_r\times\mathbb R)\setminus(\{0\}\times\epsilon\mathbb Z)\bigr)/\epsilon\mathbb Z.
\]
容易验证，有诱导的光滑流形映射 $\bar\pi\colon\overline X\to\overline Y$，其在 $Y$ 上限制为 $S^1$ 丛 $\pi\colon X\to Y$，局部 Chern 数为 $\pm1$，符号依赖于纤维定向。$\bar\pi$ 恰为\BookRef{constr:6.86}{构造6.86}中 $\pi$ 的拓扑紧化，这并不令人意外：椭圆纤维化 $I_1$ 纤维上的奇点，对应于在点 $(\{0\}\times\epsilon\mathbb Z)/\epsilon\mathbb Z$ 上方插入 NUT 模型。现在的计划，是对 $\pi\colon X\to Y$ 应用 Gibbons--Hawking 拟设，定义 $X$ 的超 Kähler 度量，再证明它延伸到 $\overline X$。不过，必须保证所定义的度量与 $\overline X$ 上给定的复结构相容。

\SourceParagraph{ch07-7.3.4-004}{572}
按照\BookSectionRef{sec:7.1.2}{7.1.2}与\BookSectionRef{sec:7.1.3}{7.1.3}，在
\[
U=(D\times\mathbb R)\setminus(\{0\}\times\epsilon\mathbb Z)
\]
上应用 Gibbons--Hawking 拟设，其中 $D\subset\mathbb C$ 为单位圆盘。用 $y_1,y_2$ 表示 $u_1,u_2$，$u=u_3$ 为 $\mathbb R$ 坐标。希望找到 $U$ 上的调和函数 $V$，关于 $u$ 以 $\epsilon$ 为周期，在 $\{0\}\times\epsilon\mathbb Z$ 的各点具有正确类型的奇点。例如，在原点附近，它应为一个调和函数加上 $1/(4\pi|x|)$，其中这里 $x=(y_1,y_2,u)$ 为实三维位置向量。这样，$X$ 上得到的度量就能延伸到 $\overline X$；这本质上是\BookRef{ex:7.5}{例题7.5}中产生 Taub--NUT 度量的计算。因此，取 $V=V_0+f(y_1,y_2)/\epsilon$，其中 $f$ 为 $D$ 上的调和函数，而
\[
V_0=\frac1{4\pi}\sum_{n=-\infty}^{\infty}
\left(\frac1{\sqrt{(u+n\epsilon)^2+y_1^2+y_2^2}}-a_{|n|}\right).
\]
对于 $n>0$，取 $a_n=1/(n\epsilon)$，以保证适当收敛；仔细选择 $a_0$，使改变 $\epsilon$ 时周期不变，从而度量定义在固定的椭圆纤维化上。

\SourceParagraph{ch07-7.3.4-005}{573}
直接的复分析说明，上述级数确实收敛到 $U$ 上的调和函数。取
\[
a_0=\frac{2(-\gamma+\log(2\epsilon))}{\epsilon},
\]
其中 $\gamma$ 为 Euler 常数。对于这一选择，可以证明下述结果，参见\BookCite{ref210}引理 3.1。
\BookStatementNumber{26}
\begin{lemma}\label{lem:7.26}
采用上述记号，当 $|y|\neq0$ 时，$V_0$ 有展开式
\[
V_0=-\frac1{4\pi\epsilon}\log|y|^2
+\frac1{2\pi\epsilon}\sum_{\substack{m\in\mathbb Z\\m\neq0}}
e^{2\pi\mathrm i mu/\epsilon}K_0(2\pi|my|/\epsilon),
\]
其中 $y=y_1+\mathrm i y_2$，$K_0$ 为修正 Bessel 函数。由此，存在常数 $C$，使对任意 $0<r_0<1$，有 $\epsilon_0>0$，当 $0<\epsilon<\epsilon_0$、$|y|>r_0$ 时，
\[
\left|V_0+\frac1{4\pi\epsilon}\log|y|^2\right|
\leq\frac C\epsilon e^{-2\pi|y|/\epsilon}.
\]
因此，在 $D$ 的固定圆环区域中，$\epsilon\to0$ 时，有 $\epsilon V_0\sim-\tfrac1{2\pi}\log r$，其中 $r^2=y_1^2+y_2^2$。此外，若 $r\leq1$，$f$ 为半径 $r$ 的闭圆盘邻域上的调和函数，且对于 $|y|\leq r$，有 $f(y)-\tfrac1{4\pi}\log|y|^2>0$，则存在 $\epsilon_0$，使对于所有 $0<\epsilon<\epsilon_0$，在整个 $D_r\times\mathbb R$ 去掉极点后，有
\[
V_0+f(y)/\epsilon>0.
\]
在极点处，理解为函数趋于正无穷。
\end{lemma}

\SourceParagraph{ch07-7.3.4-006}{573}
现设 $D_r\subseteq\mathbb C$ 为原点中心、半径 $r<1$ 的圆盘。令 $h(y)=f(y_1,y_2)+\mathrm i g(y_1,y_2)$ 为 $D_r$ 上的全纯函数，使
\[
-\frac1{4\pi}\log|y|^2+f(y_1,y_2)>0.
\]
设 $V_0$ 为\BookRef{lem:7.26}{引理7.26}定义的调和函数，取 $V=V_0+f/\epsilon$；如有需要，将圆盘略缩小，再选择足够小的 $\epsilon$，使 $V$ 在 $Y$ 上正定。

\SourceParagraph{ch07-7.3.4-007}{574}
\BookStatementNumber{27}
\begin{proposition}\label{prop:7.27}
$X$ 上存在联络一形式 $\theta$，满足 $\mathrm d\theta/(2\pi\mathrm i)=*\mathrm dV$，从而通过 Gibbons--Hawking 拟设定义一个超 Kähler 度量。此度量延伸到 $\overline X$，给出全纯椭圆纤维化 $\overline X\to D_r$，周期为一，以及
\[
\frac1{2\pi\mathrm i}\log y+\mathrm i h(y)+C,
\]
其中 $C$ 为某个实常数。适当选择 $\theta$，可以令 $C=0$。
\end{proposition}

\SourceParagraph{ch07-7.3.4-008}{574}
下面简要说明命题中关于周期的陈述。
\begin{proof}[周期部分的证明]
由构造，一个周期自动为常数一。另一个周期 $\tau(y)$ 在 $y\neq0$ 时，局部由 $\int_\gamma\mathrm dx$ 给出，其中 $x,y$ 为全纯正则坐标，$\gamma$ 为纤维 $X_y$ 中同构映到 $\{y\}\times S^1\subseteq Y$ 的圆。利用\BookSectionRef{sec:7.1.3}{7.1.3}中正则全纯坐标与 Gibbons--Hawking 坐标的关系，写成
\[
\tau(y)=\int_\gamma\theta_0-\mathrm i\int_\gamma V\,\mathrm du,
\]
其中，像先前一样，$\theta_0=\theta/(2\pi\mathrm i)$ 为实一形式。计算虚部，并利用引理 7.26 的 Fourier 展开，得到
\[
\int_\gamma\mathrm dx_2=-\int_\gamma V\,\mathrm du
=\pm\left(\frac1{4\pi}\log|y|^2-f(y_1,y_2)\right).
\]
选择 $\gamma$ 的定向，使这里取负号，即 $\int_\gamma\mathrm du=-\epsilon$。于是，$\int_\gamma\mathrm dx_1$ 必为 $-\tfrac1{4\pi}\log|y|^2+f(y_1,y_2)$ 的相应局部调和共轭，周期因而为
\[
\frac1{2\pi\mathrm i}\log y+\mathrm i h(y)+C,
\qquad C\in\mathbb R.
\]
可以给 $\theta_0$ 加上 $a\,\mathrm du$，$a\in\mathbb R$，而不改变 $\mathrm d\theta_0=*\mathrm dV$。作此改变后，周期实部变为
\[
\int_\gamma(\theta_0+a\,\mathrm du)
=-a\epsilon+\int_\gamma\theta_0.
\]
因此，适当选择 $a$，即可令 $C=0$，得到所述周期。
\end{proof}

\SourceParagraph{ch07-7.3.4-009}{574}
这样，就保证了所定义的全部超 Kähler 度量，与 $\overline X$ 的固定复结构相容。此外，在 $D$ 的固定圆环区域上，引理 7.26 的 $V_0$ 渐近公式，可以重新解释为：Ooguri--Vafa 度量以指数速度趋近半平坦度量。注意，周期为上述一与 $\tfrac1{2\pi\mathrm i}\log y+\mathrm i h(y)$ 时，半平坦度量对应于
\[
V=\epsilon^{-1}\left(-\frac1{4\pi}\log|y|^2+f(y)\right)=W^{-1},
\]
其中 $W=\epsilon/\operatorname{Im}\tau$，与上一小节相同。

\SourceParagraph{ch07-7.3.4-010}{575}
Ooguri--Vafa 度量公式的显式性，特别是引理 7.26 的 $V_0$ Fourier 展开，使我们能精确估计直径与曲率。当纤维面积为 $\epsilon$ 时，\BookCite{ref210}命题 3.5证明，奇异纤维的直径为 $O((\epsilon\log\epsilon^{-1})^{1/2})$，而固定光滑纤维的直径仅为 $O(\epsilon^{1/2})$，与半平坦情形的计算一致。几何上，这意味着，若将度量缩放，使纤维面积为一，奇异纤维的直径为 $O((\log\epsilon^{-1})^{1/2})$，并在 $\epsilon\to0$ 时变得任意大。

\SourceParagraph{ch07-7.3.4-011}{575}
要估计 Ooguri--Vafa 度量的曲率 $R$，使用\BookSectionRef{sec:7.1.2}{7.1.2}给出的 $\|R\|^2$ 公式；它适用于任意由 Gibbons--Hawking 拟设得到的度量。再次利用 $V_0$ 的 Fourier 展开，通过稍繁复的论证，\BookCite{ref210}命题 3.8给出下述估计。
\BookStatementNumber{28}
\begin{proposition}\label{prop:7.28}
令 $R(\epsilon)$ 为上述 Ooguri--Vafa 度量的曲率张量，纤维面积为 $\epsilon$。在固定的局部模型区域中，存在与 $\epsilon$ 无关的正常数 $C,C'$，使对所有充分小的 $\epsilon$，
\[
C'\epsilon^{-1}(\log\epsilon^{-1})^{-2}
<\|R(\epsilon)\|_{C^0}
<C\epsilon^{-1}\log\epsilon^{-1},
\]
其中 $\|\cdot\|_{C^0}$ 为通常的 $C^0$ 范数。
\end{proposition}

\SourceParagraph{ch07-7.3.4-012}{575}
因此，Ooguri--Vafa 度量在椭圆纤维化奇异纤维的节点附近，出现曲率爆破。这与远离奇异纤维的固定光滑纤维附近，或者半平坦度量的行为很不相同：在那里，$\|R(\epsilon)\|_{C^0}=O(\epsilon)$。后一事实可以直接证明，也可以将半平坦度量写成 Gibbons--Hawking 拟设，再次使用\BookSectionRef{sec:7.1.2}{7.1.2}的公式。

\subsection{Ricci 平坦度量的精确近似}\label{sec:7.3.5}

\SourceParagraph{ch07-7.3.5-001}{575}
现在，为\BookRef{prob:7.25}{问题7.25}中的 Ricci 平坦度量，构造在 $\epsilon=\epsilon_l\to0$ 时非常精确的近似。像先前一样，设 $f\colon X\to B=\mathbb{CP}^1$ 为具有全纯截面 $\sigma_0$ 的椭圆纤维化，$\Delta=\{p_1,\ldots,p_{24}\}$ 为判别集，其上方纤维具有 $I_1$ 奇点。令 $B_0=B\setminus\Delta$、$X_0=f^{-1}(B_0)$。\BookSectionRef{sec:7.3.3}{7.3.3}给出了 $X_0$ 上的标准半平坦度量 $\omega_{\mathrm{SF}}$，纤维面积为 $\epsilon$；以下要始终记住，此度量依赖于 $\epsilon$。所用方法，是在每条奇异纤维的邻域，将经适当平移调整的 Ooguri--Vafa 度量，粘入这一背景度量，构造 $X$ 上的近似 Ricci 平坦度量。

\SourceParagraph{ch07-7.3.5-002}{576}
先集中于固定点 $p\in\Delta$。设 $U$ 为 $p$ 的可缩邻域，$U\cap\Delta=\{p\}$；令 $U^*=U\setminus\{p\}$，并将 $f^{-1}(U)$、$f^{-1}(U^*)$ 分别记为 $X_U$、$X_{U^*}$。为进行粘合，需要\BookCite{ref210}引理 4.3证明的下述事实。

\SourceParagraph{ch07-7.3.5-003}{576}
\BookStatementNumber{29}
\begin{lemma}\label{lem:7.29}
设 $\omega_{\mathrm{SF}}$ 为 $X_0$ 上纤维面积 $\epsilon$ 的标准半平坦度量的 Kähler 形式，$\omega$ 为 $X_U$ 上同样具有纤维面积 $\epsilon$ 的 Kähler 形式。则
\[
[\omega_{\mathrm{SF}}-\omega]=0\quad\text{于 }H^2(X_{U^*},\mathbb R).
\]
此外，存在 $f\colon X_U\to U$ 的全纯截面 $\sigma$，以及 $X_{U^*}$ 上的实函数 $\varphi$，使
\[
\omega_{\mathrm{SF}}-T_\sigma^*\omega
=\mathrm i\partial\bar\partial\varphi,
\]
其中 $T_\sigma$ 为沿截面 $\sigma$ 的平移。
\end{lemma}

\SourceParagraph{ch07-7.3.5-004}{576}
在 $p$ 的邻域选择全纯坐标 $y$，将 $f$ 的全纯周期写为 $\tau_1(y),\tau_2(y)$，其中 $\tau_1$ 单值。在 $T^*_U$ 中，这对应于全纯微分 $\tau_1(y)\,\mathrm dy$。局部存在函数 $g(y)$，满足 $\mathrm dg=\tau_1(y)\,\mathrm dy$；由于 $\tau_1(p)\neq0$，可以使用 $g$ 作为 $p$ 附近的局部全纯坐标。以 $g$ 代替 $y$，就可以假设 $\tau_1(y)=1$，并令 $p$ 对应于 $y=0$。上一小节的结果说明，对所有 $0<\epsilon<\epsilon_0$，可以在 $f^{-1}(U)$ 上构造纤维面积 $\epsilon$ 的 Ooguri--Vafa 度量 $\omega_{\mathrm{OV}}$。这里，$U=\{y\mid|y|<r\}$，$r$ 仅依赖于周期 $\tau_2$ 及 $\epsilon_0$，与 $\epsilon$ 无关。固定 $r_1<r_2<r$，令 $U_i=\{y\mid|y|<r_i\}$，并记圆环区域为 $U_{r_1,r_2}=\{y\mid r_1<|y|<r_2\}$。

\SourceParagraph{ch07-7.3.5-005}{576}
利用适当的固定 $C^\infty$ 截断函数 $\psi$，满足 $\psi(s)=1$ 当 $s\leq r_1^2$、$\psi(s)=0$ 当 $s\geq r_2^2$，就能将局部度量 $\omega_{\mathrm{OV}}$ 粘入背景度量 $\omega_{\mathrm{SF}}$。关键在于，由\BookRef{lem:7.29}{引理7.29}，存在全纯截面 $\sigma$ 与 $X_{U^*}$ 上的函数 $\varphi$，在 $U^*$ 上方满足
\[
\omega_{\mathrm{SF}}-T_\sigma^*\omega_{\mathrm{OV}}
=\mathrm i\partial\bar\partial\varphi.
\]
考虑
\[
\omega_{\mathrm{SF}}
-\mathrm i\partial\bar\partial\bigl(\psi(|y|^2)\varphi\bigr),
\]
就得到闭实 $(1,1)$ 形式：$|y|\leq r_1$ 时，它与 $T_\sigma^*\omega_{\mathrm{OV}}$ 一致；$|y|\geq r_2$ 时，与 $\omega_{\mathrm{SF}}$ 一致。在每条奇异纤维处实施这一粘合，得到全局闭实 $(1,1)$ 形式 $\omega_{\mathrm{new}}$。可以验证，当 $\epsilon$ 足够小时，它正定，因而定义 $X$ 上的 Kähler 度量。

\SourceParagraph{ch07-7.3.5-006}{576}
另一方面，还不能保证 $\int_X\omega_{\mathrm{new}}^2=\int_X(\operatorname{Re}\Omega)^2$。不过，回顾 $\epsilon$ 很小时，$\omega_{\mathrm{SF}}$ 与 $\omega_{\mathrm{OV}}$ 在 $U_{r_1,r_2}$ 上方只相差 $O(e^{-C/\epsilon})$。由构造，在各粘合圆环上方之外，有 $\omega_{\mathrm{new}}^2=(\operatorname{Re}\Omega)^2$；在粘合区域内，二者仅相差 $O(e^{-C/\epsilon})$。所以，
\[
\int_X\omega_{\mathrm{new}}^2-\int_X(\operatorname{Re}\Omega)^2
=O(e^{-C/\epsilon}).
\]
注意，$E$ 为纤维类时，
\[
([\omega_{\mathrm{new}}]+aE)^2
=[\omega_{\mathrm{new}}]^2+2a\epsilon.
\]
因此，可以在 $B$ 上找到支撑于粘合圆环区域的二形式 $\alpha$，对每个 $k$，其 $C^k$ 范数为 $O(e^{-C/\epsilon})$，使
\[
\int_X(\omega_{\mathrm{new}}+f^*\alpha)^2
=\int_X(\operatorname{Re}\Omega)^2.
\]
令 $\omega_\epsilon=\omega_{\mathrm{new}}+f^*\alpha$。由于 $\alpha$ 仍很小，$\omega_\epsilon$ 保持正定，定义所需的 Kähler 度量。

\SourceParagraph{ch07-7.3.5-007}{577}
这一粘合度量在粘合区域之外显然 Ricci 平坦；粘合区域就是 $\Delta$ 各点附近所选圆环的上方。它也非常接近全局 Ricci 平坦，Ricci 形式至多为 $O(e^{-C/\epsilon})$。此外，设
\[
F_\epsilon=\log\left(\frac{\Omega\wedge\overline\Omega/2}{\omega_\epsilon^2}\right),
\]
则 $F_\epsilon$ 及其相对于 $\omega_\epsilon$ 的 Laplace 算子的 $C^0$ 范数，均至多为 $O(e^{-C/\epsilon})$。前两小节的计算说明，$\omega_\epsilon$ 的直径至多为 $O(\epsilon^{-1/2})$；曲率张量 $R$ 满足 $\|R\|_{C^0}\to\infty$，但至多以 $O(\epsilon^{-1}\log\epsilon^{-1})$ 的速度增长。在任意固定光滑纤维附近，有 $\|R\|\leq C\epsilon$，其中 $C$ 依赖于该纤维。

\SourceParagraph{ch07-7.3.5-008}{577}
现在断言，以此方式构造的度量，是 Ricci 平坦度量的非常精确的近似。为证明这一点，分析 Yau 的 Ricci 平坦度量存在性证明，其方式类似于\BookSectionRef{sec:7.3.1}{7.3.1}中 Kobayashi 对 K3 曲面向有理双点曲面的 Kähler 退化的处理。实行该程序的起始度量，当然为 $\omega_\epsilon$。Yau 的证明给出 $X$ 上的函数 $u_\epsilon$，满足
\[
(\omega_\epsilon+\mathrm i\partial\bar\partial u_\epsilon)^2
=e^{F_\epsilon}\omega_\epsilon^2,
\qquad\int_Xu_\epsilon\omega_\epsilon^2=0.
\]
进一步分析证明，得到 $u_\epsilon$ 的各种估计，证实 $\omega_\epsilon$ 确为 Ricci 平坦度量的精确近似；将后者记为 $\widetilde\omega_\epsilon$。这一分析在\BookCite{ref210}第 5 节进行。

\SourceParagraph{ch07-7.3.5-009}{577}
\begin{enumerate}[label=\textup{(\Roman*)}]
\item 所得 $C^0$ 估计为：存在与 $\epsilon$ 无关的正常数 $C_1,C_2$，使
\[
\|u_\epsilon\|_\infty\leq C_1\epsilon^{-5}e^{-C_2/\epsilon}.
\]
证明利用 Sobolev 不等式。上述直径估计，保证相应归一化的 Sobolev 常数有 $c\epsilon^5$ 的下界。再利用 $F_\epsilon$ 的 $C^0$ 范数至多为 $O(e^{-C/\epsilon})$，就得到此估计。当然，$\epsilon$ 很小时，指数项才是主导项，因此 $C^0$ 范数随 $\epsilon\to0$ 指数趋于零。
\item 所得 $C^2$ 估计为：存在与 $\epsilon$ 无关的常数 $C>1$，使对所有充分小的 $\epsilon$，
\[
C^{-1}\omega_\epsilon\leq\widetilde\omega_\epsilon\leq C\omega_\epsilon.
\]
这一估计来自 Yau 使用极大值原理的论证，需要上述关于 $\omega_\epsilon$ 的各种界。它是 Yau 证明以及 $u_\epsilon$ 估计的核心。这已经蕴含：在 Ricci 平坦度量 $\widetilde\omega_\epsilon$ 下，$f\colon X\to\mathbb{CP}^1$ 的纤维随 $\epsilon\to0$ 塌缩，参见下文\BookSectionRef{sec:7.3.6}{7.3.6}。
\item 所得 $C^{k,\alpha}$ 估计，$k\geq2$、$0<\alpha<1$，如下：对于任意单连通开集 $U\subset B$，其闭包 $\overline U\subset B_0=B\setminus\Delta$，存在依赖于 $U,k,\alpha$、但与 $\epsilon$ 无关的正常数 $C_1,C_2$，使在 $U$ 上方，对充分小的 $\epsilon$，
\[
\|u_\epsilon\|_{C^{k,\alpha}}\leq C_1e^{-C_2/\epsilon}.
\]
先应用一般结果\BookCite{ref179}定理 17.14，从 $C^2$ 估计得到 $C^{2,\alpha}$ 估计，再通过正则性提升与 Schauder 估计的标准论证，得到更强的结论。
\end{enumerate}

\SourceParagraph{ch07-7.3.5-010}{578}
这些估计的证明之所以成立，根本原因是：当 $\epsilon\to0$ 时，直径、曲率等发散量，至多以 $\epsilon^{-1}$ 的多项式速度增长；而 $F_\epsilon$、Ricci 曲率等趋于零的量，以指数速度衰减，由 $O(e^{-C/\epsilon})$ 控制。分析 Yau 证明给出的先验估计，可以看出，后一行为始终支配前一行为。

\SourceParagraph{ch07-7.3.5-011}{578}
因此，$\omega_\epsilon$ 在全局上与 Ricci 平坦度量 $\widetilde\omega_\epsilon$ 相差至多一个固定倍数；在 $\mathbb{CP}^1\setminus\Delta$ 的相对紧子集上，两者在任意 $C^{k,\alpha}$ 范数下都任意接近，$\epsilon$ 足够小时，误差为 $O(e^{-C/\epsilon})$。下一小节描述度量空间的 Gromov--Hausdorff 收敛时，将看到这些结论的几何意义。

\subsection{Gromov--Hausdorff 收敛}\label{sec:7.3.6}

\SourceParagraph{ch07-7.3.6-001}{578}
前文多次提到度量收敛，现在给出精确定义。上一小节的结果，将蕴含一个关于 K3 曲面复结构向大复结构极限退化的重要结论；等价地，也就是上述椭圆 K3 曲面的 Kähler 退化。适当归一化后，在极限中，K3 曲面实际上收敛到带有标准度量的二球面。由此将引出一个一般猜想，描述 Calabi--Yau 流形的复结构向大复结构极限退化时的度量极限；它推动了近来多位作者的工作。

\SourceParagraph{ch07-7.3.6-002}{578}
\BookStatementNumber{30}
\begin{definition}\label{def:7.30}
设 $(X,d_X)$、$(Y,d_Y)$ 为两个紧度量空间。它们的 Gromov--Hausdorff 距离定义为
\[
d_{\mathrm{GH}}(X,Y)
=\inf_{Z,\iota_X,\iota_Y}
d_H^Z\bigl(\iota_X(X),\iota_Y(Y)\bigr),
\]
其中遍历所有共同度量空间 $Z$ 及等距嵌入 $\iota_X\colon X\hookrightarrow Z$、$\iota_Y\colon Y\hookrightarrow Z$，$d_H^Z$ 为 Hausdorff 距离，即使两个集合均包含在对方的 $\epsilon$ 邻域中所需 $\epsilon$ 的下确界。

\SourceParagraph{ch07-7.3.6-003}{579}
也可以用近似映射刻画接近程度。设存在映射 $f\colon X\to Y$、$g\colon Y\to X$，不必连续，使对所有 $x_1,x_2\in X$，
\[
|d_X(x_1,x_2)-d_Y(f(x_1),f(x_2))|<\epsilon,
\]
且对所有 $x\in X$，$d_X(x,g\circ f(x))<\epsilon$；同时满足交换 $X,Y$ 及 $f,g$ 所得的两个对称条件。则 $d_{\mathrm{GH}}(X,Y)\leq\epsilon$。所有这类 $\epsilon$ 的下确界记为 $\delta(X,Y)$，有
\[
d_{\mathrm{GH}}(X,Y)\leq\delta(X,Y)\leq2d_{\mathrm{GH}}(X,Y),
\]
因此也可以通过这些近似映射的误差趋于零来定义相同的收敛概念。
\end{definition}

\SourceParagraph{ch07-7.3.6-004}{579}
Gromov--Hausdorff 距离在紧度量空间的等距同构类上定义一个拓扑，从而给出收敛概念。由 Gromov 的结果，例如\BookCite{ref391}第 281 页推论 1.11，固定维数、直径至多为 $D$ 的紧 Ricci 平坦流形族预紧。任意序列都可以取一个子列，在 Gromov--Hausdorff 意义下收敛到紧长度空间。长度空间是指两点距离等于连接它们的路径长度的下确界的度量空间；在紧情形，这个下确界由某条连接两点的曲线达到。特别地，若一列 Calabi--Yau $n$ 维流形的复结构趋近于大复结构极限点，或者其他边界点，并且度量直径有统一上界，则存在收敛子列。现在的基本问题是，极限究竟是什么？若极限的实维数小于 $2n$，称发生塌缩。

\SourceParagraph{ch07-7.3.6-005}{579}
对于上一小节研究的 K3 曲面退化，那里的结果回答了这个问题，不过要将度量归一化为直径有界，而非体积有界；将所有度量乘以 $\epsilon$ 即可达到这一归一化。下面以更一般的形式陈述结果，不假设椭圆纤维化 K3 曲面一定具有全纯截面。

\SourceParagraph{ch07-7.3.6-006}{579}
\BookStatementNumber{31}
\begin{theorem}\label{thm:7.31}
设 $j\colon J\to B$ 为具有截面的椭圆纤维化 K3 曲面，奇异纤维均为 $I_1$ 型；设 $f_i\colon X_i\to B$ 为 Jacobian 均为 $j$ 的椭圆纤维化 K3 曲面序列。设 $\omega_i$ 对应于 $X_i$ 上的 Ricci 平坦 Kähler 度量，$\int_{X_i}\omega_i^2$ 与 $i$ 无关，且
\[
\int_{f_i^{-1}(b)}\omega_i=\epsilon_i\longrightarrow0
\quad\text{当 }i\to\infty.
\]
则 Riemann 流形序列 $(X_i,\epsilon_i\omega_i)$，在 Gromov--Hausdorff 意义下收敛到 $B$。$B$ 的度量为判别集补集上的标准 Riemann 度量的度量完备化。局部坐标中，将周期写为 $\tau_1(y),\tau_2(y)$，且 $\operatorname{Im}(\bar\tau_1\tau_2)>0$，则此 Riemann 度量为
\[
\operatorname{Im}(\bar\tau_1\tau_2)
(\mathrm dy_1^2+\mathrm dy_2^2),\qquad y=y_1+\mathrm i y_2,
\]
也就是将 $\operatorname{Im}(\bar\tau_1\tau_2)\,\mathrm dy\otimes\mathrm d\bar y$ 理解为实对称度量。
\end{theorem}

\SourceParagraph{ch07-7.3.6-007}{579}
必须指出，此定理一点也不显然。仅仅纤维面积趋于零，并不意味着纤维不会有异常行为，例如越来越长、越来越细，以致面积趋于零，却仍然保持很“大”。

\SourceParagraph{ch07-7.3.6-008}{580}
\begin{proof}[证明思路]
按照\BookSectionRef{sec:7.3.5}{7.3.5}近似 Ricci 平坦度量 $\omega_i$。远离奇异纤维，度量非常迅速地收敛到半平坦度量，从而在度量意义下，收敛到带上述 Riemann 度量的 $B$ 的相应开子集。在奇异纤维的固定邻域内，只能使用第（II）项的 $C^2$ 估计，尚不足以得到所说的度量收敛。不过，若容许奇异纤维的邻域收缩，使 $B$ 中所选圆环的半径随 $i\to\infty$ 适当地趋于零，再在奇异纤维附近使用 $C^2$ 估计，就可以直接论证，Gromov--Hausdorff 极限确如所述。可比较\BookRef{note:7.23}{笔记7.23}中任意 K3 曲面向具有有理双点曲面的退化。

\SourceParagraph{ch07-7.3.6-009}{580}
判别集补集 $B_0$ 上的极限 Riemann 度量，恰为\BookSectionRef{sec:6.1.1}{6.1.1}中 McLean 度量的一个倍数；它对应于将复结构通过超 Kähler 旋转回到 $I$ 后得到的特殊 Lagrange 环面纤维化 $X_0\to B_0$，并使用半平坦度量。事实上，回顾，对于 $b\in B_0$ 处切向量 $u,v$，McLean 度量定义为
\[
\begin{aligned}
g(u,v)
&=-\int_{f^{-1}(b)}\iota(u)\omega_I\wedge\iota(v)\operatorname{Im}\Omega_I\\
&=-\int_{f^{-1}(b)}\iota(u)\operatorname{Im}\Omega_K
\wedge\iota(v)\operatorname{Re}\Omega_K,
\end{aligned}
\]
局部坐标中，$\Omega_K=\mathrm dx\wedge\mathrm dy$。容易计算得
\[
2g\left(\frac\partial{\partial y},\frac\partial{\partial\bar y}\right)
=\int_{f^{-1}(b)}\mathrm dx_1\wedge\mathrm dx_2
=\operatorname{Im}(\bar\tau_1\tau_2).
\]
\end{proof}

\SourceParagraph{ch07-7.3.6-010}{580}
上述定理引出一个关于 Calabi--Yau $n$ 维流形的猜想，提出于\BookCite{ref210}第 6 节；Kontsevich、Soibelman 和 Todorov 也独立提出了它。
\BookStatementNumber{32}
\begin{conjecture}\label{conj:7.32}
设 $\overline{\mathcal M}$ 是单连通 Calabi--Yau $n$ 维流形 $X$ 的复变形模空间的一个紧化，$X$ 的完整群为 $\operatorname{SU}(n)$；设 $p\in\overline{\mathcal M}$ 为大复结构极限点，其精确 Hodge 理论定义见\BookCite{ref365}。设 $(X_i,g_i)$ 为带 Ricci 平坦 Kähler 度量的 Calabi--Yau 流形序列，均为 $X$ 的复变形，$[X_i]\in\mathcal M$ 以适当方式趋近 $p$，且对所有 $i$，
\[
C_1\geq\operatorname{Diam}(X_i)\geq C_2>0.
\]
则 $(X_i,g_i)$ 的某个子列收敛到度量空间 $(X_\infty,d_\infty)$，其中 $X_\infty$ 同胚于 $S^n$。此外，存在余维二的子集 $\Delta\subseteq X_\infty$，使 $X_\infty\setminus\Delta$ 带有仿射结构，且 $d_\infty$ 由其上的 Monge--Ampère Riemann 度量诱导。
\end{conjecture}

\SourceParagraph{ch07-7.3.6-011}{580}
“以适当方式趋近 $p$”这一说法，在\BookCite{ref210}第 6 节讨论。特别需要排除极限实维数严格小于 $n$ 的情形。例如，在镜像中，若趋近大 Kähler 结构极限的同时，也趋近射影化 Kähler 锥的边界，就可能导致这种情况。

\subsection{SYZ 猜想的未来}\label{sec:7.3.7}

\SourceParagraph{ch07-7.3.7-001}{581}
Strominger--Yau--Zaslow 猜想提出已有十多年。如本章解释的那样，它的原始形式看来极不可能得到证明。不过，\BookRef{conj:7.32}{猜想7.32}提示了 SYZ 猜想的极限版本，似乎更有希望证明。

\SourceParagraph{ch07-7.3.7-002}{581}
从 Calabi--Yau $n$ 维流形的极大幺幂退化 $\mathcal X\to D$ 出发，其中 $D$ 为圆盘，并取 $t_i\to0\in D$。假设 $t_i$ 足够接近零时，存在特殊 Lagrange $T^n$，其同调类在单值变换下不变；更具体地，它生成相应单值变换权滤过的 $W_0$。这正是我们期待与极大幺幂退化有关的特殊 Lagrange 纤维化的纤维出现的位置。令 $B_{0,i}$ 为此环面的变形模空间；其每点对应于 $\mathcal X_{t_i}$ 中一个光滑特殊 Lagrange 环面。此流形带有\BookSectionRef{sec:6.1.1}{6.1.1}定义的 McLean 度量及仿射结构。再作紧化 $B_{0,i}\subseteq B_i$，或许可以取 $B_{0,i}$ 在特殊 Lagrange 流空间中的闭包；这里不讨论细节。希望这样得到紧度量空间序列 $(B_i,d_i)$，$d_i$ 由 McLean 度量诱导。将 McLean 度量归一化，使 $B_i$ 的直径为与 $i$ 无关的常数，就可以期待它收敛到紧度量空间 $(B_\infty,d_\infty)$。SYZ 的极限形式如下。

\SourceParagraph{ch07-7.3.7-003}{581}
\BookStatementNumber{33}
\begin{conjecture}\label{conj:7.33}
若 $(\mathcal X_{t_i},g_i)$ 收敛到 $(X_\infty,g_\infty)$，且充分大的 $i$ 对应的 $(B_i,d_i)$ 非空，并收敛到 $(B_\infty,d_\infty)$，则 $B_\infty$ 与 $X_\infty$ 除尺度因子外等距。此外，存在子空间 $B_{\infty,0}\subseteq B_\infty$，使 $\Delta:=B_\infty\setminus B_{\infty,0}$ 在 $B_\infty$ 中具有 Hausdorff 余维二，$B_{\infty,0}$ 为 Monge--Ampère 流形，其 Monge--Ampère 度量诱导 $d_\infty$ 在 $B_{\infty,0}$ 上的限制。
\end{conjecture}

\SourceParagraph{ch07-7.3.7-004}{581}
本质上，这就是说，趋近大复结构极限时，特殊 Lagrange 环面填充 Calabi--Yau 流形中越来越大的部分。不过，在模空间中，没有哪个点能够保证在整个流形上得到特殊 Lagrange 纤维化。

\SourceParagraph{ch07-7.3.7-005}{581}
猜想采用这一表述，是因为它比更强的形式更有希望证明。下面概述一种证明途径。
\begin{enumerate}
\item 给定极大幺幂退化 $\mathcal X\to D$，猜测极限对象 $X_\infty=B_\infty$ 是什么，但先仅将它看作带奇点的仿射流形。
\item 解决仿射 Calabi 猜想，即在 $B_\infty$ 上找到 Monge--Ampère 度量。由此得到 $X(B_{\infty,0})$ 上的 Ricci 平坦度量。
\item 在 $X(B_{\infty,0})$ 的紧化的奇异纤维附近，寻找 Ricci 平坦度量的局部模型。
\item 粘入这些局部模型，像\BookSectionRef{sec:7.3.5}{7.3.5}那样，证明所得度量接近真正的 Ricci 平坦度量。这将很好地描述大复结构极限附近的 Calabi--Yau 流形的 Ricci 平坦度量。
\item 由于远离奇异纤维时，度量接近半平坦度量，可以期待，半平坦度量的特殊 Lagrange 纤维，会变形为真正 Ricci 平坦度量下的特殊 Lagrange 环面。当 $t_i$ 接近零时，由此得到模空间 $B_{i,0}$ 的开子集，进而比较 $X_i$ 与 $B_i$，证明极限结论。
\end{enumerate}

\SourceParagraph{ch07-7.3.7-006}{582}
第一步已在很大程度上得到理解。事实上，\BookSectionRef{sec:6.4.1}{6.4.1}给出了一些带奇点的仿射流形例子，猜测它们就是正确的极限；至少在拓扑层面，确实如此。第二、三步最具挑战性。对于维数大于二的仿射 Calabi 猜想，当时还没有成功证明，但\BookCite{ref338,ref339}有这一方向的工作。第三步方面，正型或负型顶点附近，还没有类似 Ooguri--Vafa 度量的局部模型，不过，\BookCite{ref476}有一些相关工作。对于第四步，则期待只需要类似于上述\BookCite{ref210}中所概述的技术。

\SourceParagraph{ch07-7.3.7-007}{582}
仍然留下一个问题：如何通过这一修订的 SYZ 猜想实施镜像对称？基本上，可以按下列步骤进行。
\begin{enumerate}
\item 从 Calabi--Yau 流形的极大幺幂退化 $\mathcal X\to D$ 出发，并选择极化。于是，对于每个 $t\in D\setminus\{0\}$，得到 Kähler 类 $[\omega_t]\in H^2(\mathcal X_t,\mathbb R)$，由 Ricci 平坦度量 $g_t$ 的 Kähler 形式 $\omega_t$ 代表。
\item 确定序列 $(\mathcal X_{t_i},r_ig_{t_i})$ 的 Gromov--Hausdorff 极限，其中 $t_i\to0$，$r_i$ 为使 $\mathcal X_{t_i}$ 直径固定的尺度因子。若上述猜想成立，极限就是带奇点的仿射流形 $B$，配有 Monge--Ampère 度量。
\item 作 Legendre 变换，得到新的带奇点仿射流形 $\check B$，其度量相同。
\item 对小 $\epsilon>0$，尝试紧化
\[
X_\epsilon(\check B_0):=T\check B_0/\epsilon\Lambda,
\]
得到复流形 $X_\epsilon(\check B)$，作为镜像流形。
\end{enumerate}

\SourceParagraph{ch07-7.3.7-008}{582}
实际上，最后一步还需进一步解释。问题在于，一般期待能够对辛流形 $\check X(B_0)$ 作辛紧化，从而得到作为辛流形的镜像；但不能期待直接将 $X_\epsilon(\check B_0)$ 紧化为复流形。在 Ooguri--Vafa 度量的描述中，已明确看到这一点：必须扰动半平坦超 Kähler 结构，才能越过奇异纤维延伸。虽然构造中保持椭圆纤维化的复结构固定，但旋转回 SYZ 所用的复结构时，这一扰动也改变了相应复结构。

\SourceParagraph{ch07-7.3.7-009}{583}
因此，一般期待，在能够紧化之前，必须先对 $X_\epsilon(\check B_0)$ 作一个“小”变形。以适当方式衡量，预计变形大小为 $O(e^{-C/\epsilon})$。此外，这个小变形对镜像对称至关重要：正是它给出 B 模型的瞬子修正。

\SourceParagraph{ch07-7.3.7-010}{583}
由于最后这个问题很重要，已有多位作者研究。Fukaya\BookCite{ref166}直接利用启发式思想研究它；Kontsevich 和 Soibelmann\BookCite{ref312}则调整了从仿射流形到复流形的问题，转而构造非 Archimedes 空间。将这一问题命名如下。
\BookStatementNumber{34}
\begin{problem}[重构问题]\label{prob:7.34}
给定整仿射流形 $B$，允许具有奇点，见\BookSectionRef{sec:6.4.1}{6.4.1}，构造复流形 $X_\epsilon(B)$，使其为 $X_\epsilon(B_0)$ 的一个小变形的紧化。
\end{problem}

\SourceParagraph{ch07-7.3.7-011}{583}
Gross 和 Siebert 的工作，近来在这一问题上取得许多进展。讨论这些工作超出本书范围，但其基本思想是：给定带奇点的整仿射流形 $B$ 及一些额外组合数据，可以构造 Calabi--Yau 流形的一个退化。实际上，可以完全在代数几何层面工作，由此建立退化的几何与仿射流形几何之间的对应。这给出 SYZ 猜想的代数几何版本，有望比这里讨论的微分几何版本强大得多。特别地，它直接联系有理曲线与周期。有兴趣的读者可参阅\BookCite{ref205,ref206,ref203,ref207,ref208}。
